% La Conscience et l’Inconscient — Philosophie Tle Tech — Cours signé OnBuch+
% Programme officiel MINESEC. Compiler avec : tectonic N18-conscience-inconscient.tex
\newcommand{\DOCMATIERE}{Philosophie}
\newcommand{\DOCNIVEAU}{Tle Tech}
\newcommand{\DOCMODULE}{Philosophie – classe de Terminale technique}
\newcommand{\DOCLECON}{N18}
\newcommand{\DOCTITRE}{La Conscience et l’Inconscient}
% OnBuch+ — préambule « cours signés » ENRICHI (compilé avec tectonic / XeLaTeX)
% Système visuel « Pop » d'OnBuch+ : fond crème (jamais blanc), bordures encre
% épaisses, ombres « dures » décalées, titres Archivo Black en capitales.
% Version MATHÉMATIQUES : graphes (pgfplots/tikz), boîtes pédagogiques
% (définition, propriété, méthode, attention, le savais-tu, exercice résolu,
% exercice, TP, bilan), page de garde et signature OnBuch+ ancrée en bas.
%
% Macros attendues dans main.tex AVANT \input{preamble} :
%   \DOCMATIERE  \DOCNIVEAU  \DOCMODULE  \DOCLECON (numéro)  \DOCTITRE
\documentclass[11pt]{article}

\usepackage{fontspec}
\usepackage[french]{babel}
\usepackage[a4paper,margin=2cm,top=3.4cm,bottom=2.8cm,headheight=1.4cm,footskip=1.3cm]{geometry}
\usepackage{amsmath,amssymb,amsfonts}
\usepackage{mathtools} % \xmapsto, \coloneqq... souvent utilisés par les modèles
\usepackage{cancel} % simplifications barrées (bilans de forces)
% \abs{z} (module, valeur absolue) et \norm{u} (norme), utilisés par les modèles
\providecommand{\abs}{}\let\abs\relax\DeclarePairedDelimiter{\abs}{\lvert}{\rvert}
\providecommand{\norm}{}\let\norm\relax\DeclarePairedDelimiter{\norm}{\lVert}{\rVert}
\usepackage[table]{xcolor} % [table] : fournit aussi \rowcolors (alternance de lignes)
\usepackage{pagecolor}
\usepackage{tikz}
\usetikzlibrary{arrows.meta,positioning,calc,shapes.geometric,shapes.misc,shapes.arrows,decorations.pathmorphing,decorations.pathreplacing,decorations.markings,patterns,shadows,mindmap,trees,fit,backgrounds,matrix,intersections,angles,quotes}
\usepackage{pgfplots}
\usepackage[version=4]{mhchem} % équations nucléaires \ce{^{235}_{92}U}
\usepackage[european,siunitx]{circuitikz} % schémas électriques
\pgfplotsset{compat=1.18}
\usepgfplotslibrary{fillbetween,groupplots} % aires entre courbes (calcul intégral)
\usepackage{enumitem}
\usepackage{fancyhdr}
% [most] charge toutes les bibliothèques tcolorbox (listings, minted, raster,
% documentation...) et gonfle inutilement la mémoire TeX sur un document long
% (~300 boîtes) : on ne charge que ce qui est réellement utilisé.
\usepackage[skins,breakable]{tcolorbox}
\usepackage{graphicx}
\usepackage{colortbl}
\usepackage{booktabs}
\usepackage{tabularx}
\usepackage{multirow}
\usepackage{diagbox} % case d'en-tête diagonale des tableaux à double entrée
% Colonnes extensibles centrée / gauche / droite (C, L, R), souvent utilisées par les modèles
\newcolumntype{C}{>{\centering\arraybackslash}X}
\newcolumntype{L}{>{\raggedright\arraybackslash}X}
\newcolumntype{R}{>{\raggedleft\arraybackslash}X}
\usepackage{array}
\usepackage{multicol}
\usepackage{siunitx}
% parse-numbers=false : accepte aussi \SI{2,5 \times 10^{-3}}{...}
\sisetup{locale=FR,output-decimal-marker={,},group-minimum-digits=4,parse-numbers=false}
\DeclareSIUnit\ohm{\ensuremath{\symup{\Omega}}} % U+2126 absent de la police
\DeclareSIUnit\division{div} % réglages d'oscilloscope (ms/div, V/div)
\DeclareSIUnit\tr{tr} % tours (vitesse de rotation, ex. \SI{1500}{\tr\per\minute})
\DeclareSIUnit\molar{M} % concentration molaire (ex. \SI{3}{\molar}, \SI{1}{\micro\molar})
\DeclareSIUnit\an{an} % année (le modèle confond parfois avec \year, un compteur TeX) — ex. \SI{5}{\kilo\meter\per\an}
\DeclareSIUnit\FCFA{FCFA} % franc CFA (ex. \SI{500}{\FCFA\per\kilo\gram})
\DeclareSIUnit\degreeBrix{\textdegree Brix} % teneur en sucre d'un sirop/jus (réfractométrie)
\DeclareSIUnit\month{mois} % le modèle utilise parfois \month, un compteur TeX, comme unité de durée
\DeclareSIUnit\microliter{\micro\liter} % le modèle écrit parfois \microliter en un seul token
\DeclareSIUnit\million{millions} % dénombrement (ex. \SI{15}{\million\per\mL} spermatozoïdes)
\usepackage{qrcode}
% \degree : commande fréquemment utilisée par les modèles pour un angle ou une
% température en dehors de \SI{}{} (non fournie par les paquets chargés).
\providecommand{\degree}{^{\circ}}
% \qed (fin de démonstration), utilisé par les modèles sans amsthm
\providecommand{\qed}{\hfill$\square$}
% \barre : double barre des valeurs interdites dans un tableau de variations (tkz-tab)
\providecommand{\barre}{\ensuremath{\|}}
% environnement proof (amsthm non chargé) utilisé par les modèles
\ifdefined\proof\else\newenvironment{proof}[1][Démonstration]{\par\noindent\textit{#1.}\ }{\qed\par}\fi
\usepackage{lastpage}
\usepackage{hyperref}
\hypersetup{hidelinks}

% ============================================================
% Palette « Pop » OnBuch+ — mêmes hex que la classe P (plus_ui.dart)
% ============================================================
\definecolor{poporange}{HTML}{F59321}
\definecolor{popdark}{HTML}{E07A0C}
\definecolor{popcream}{HTML}{FFF6E8}
\definecolor{popink}{HTML}{1C1714}
\definecolor{poppurple}{HTML}{7A5AE0}
\definecolor{popgreen}{HTML}{2E9E6A}
\definecolor{poppink}{HTML}{E0457B}
\definecolor{popblue}{HTML}{2F7DE1}
\definecolor{popgold}{HTML}{C98A12}
\definecolor{popmuted}{HTML}{8A7F76}
\definecolor{popline}{HTML}{EADFD2}
\definecolor{popgray}{HTML}{8A7F76}
\colorlet{popgrey}{popgray}
% Alias tolérants (le modèle nomme parfois les couleurs différemment)
\colorlet{popwhite}{white}
\colorlet{popblack}{popink}
\colorlet{popred}{poppink}
\colorlet{popyellow}{popgold}
% Teintes claires (fonds de tableaux/encadrés) — le modèle nomme parfois
% les couleurs avec un suffixe L (ex. popblueL) sans les avoir définies.
\colorlet{poporangeL}{poporange!20}
\colorlet{popdarkL}{popdark!20}
\colorlet{poppurpleL}{poppurple!20}
\colorlet{popgreenL}{popgreen!20}
\colorlet{poppinkL}{poppink!20}
\colorlet{popblueL}{popblue!20}
\colorlet{popgoldL}{popgold!20}
\colorlet{popredL}{popred!20}
\colorlet{popgrayL}{popgray!20}
% Teintes claires pour fonds de tableaux / zones
\colorlet{popblueL}{popblue!10!white}
\colorlet{poporangeL}{poporange!14!white}
\colorlet{popgreenL}{popgreen!12!white}
\colorlet{poppurpleL}{poppurple!12!white}
\colorlet{poppinkL}{poppink!12!white}
\colorlet{popgoldL}{popgold!14!white}

\pagecolor{popcream}

% ============================================================
% Typographie
% ============================================================
\defaultfontfeatures{Ligatures=TeX}
\setmainfont{PlusJakartaSans-Medium.ttf}[Path=../../../fonts/,BoldFont=PlusJakartaSans-Bold.ttf]
% Maths en sans-serif (Fira Math) avec chiffres et lettres droites de Plus
% Jakarta, pour que les formules se fondent dans le texte.
\usepackage[math-style=ISO,bold-style=ISO]{unicode-math}
\setmathfont{FiraMath-Regular.otf}[Path=../../../fonts/]
\setmathfont{PlusJakartaSans-Medium.ttf}[Path=../../../fonts/,range={up/num,up/{Latin,latin},\mathup}]
\usepackage{icomma} % virgule décimale sans espace en mode math
% Caractères mathématiques Unicode absents des polices de texte (Plus Jakarta,
% Archivo Black) : rendus par leur équivalent mathématique, en texte comme en
% formule — sinon ils s'affichent en carré vide (ex. « Arithmétique dans ℤ »).
\usepackage{newunicodechar}
\newunicodechar{ℝ}{\ensuremath{\mathbb{R}}}
\newunicodechar{ℤ}{\ensuremath{\mathbb{Z}}}
\newunicodechar{ℂ}{\ensuremath{\mathbb{C}}}
\newunicodechar{ℕ}{\ensuremath{\mathbb{N}}}
\newunicodechar{ℚ}{\ensuremath{\mathbb{Q}}}
\newunicodechar{𝔻}{\ensuremath{\mathbb{D}}}
\newunicodechar{α}{\ensuremath{\alpha}}
\newunicodechar{β}{\ensuremath{\beta}}
\newunicodechar{γ}{\ensuremath{\gamma}}
\newunicodechar{δ}{\ensuremath{\delta}}
\newunicodechar{ε}{\ensuremath{\varepsilon}}
\newunicodechar{θ}{\ensuremath{\theta}}
\newunicodechar{λ}{\ensuremath{\lambda}}
\newunicodechar{μ}{\ensuremath{\mu}}
\newunicodechar{σ}{\ensuremath{\sigma}}
\newunicodechar{τ}{\ensuremath{\tau}}
\newunicodechar{φ}{\ensuremath{\varphi}}
\newunicodechar{ψ}{\ensuremath{\psi}}
\newunicodechar{ω}{\ensuremath{\omega}}
\newunicodechar{Σ}{\ensuremath{\Sigma}}
% majuscules produites par \MakeUppercase dans les titres (α-aminés -> Α)
\newunicodechar{Α}{\ensuremath{\alpha}}
\newunicodechar{Β}{\ensuremath{\beta}}
\newunicodechar{Γ}{\ensuremath{\Gamma}}
\newunicodechar{Δ}{\ensuremath{\Delta}}
\newunicodechar{Ε}{\ensuremath{\varepsilon}}
\newunicodechar{Θ}{\ensuremath{\Theta}}
\newunicodechar{Λ}{\ensuremath{\Lambda}}
\newunicodechar{Μ}{\ensuremath{\mu}}
\newunicodechar{Τ}{\ensuremath{\tau}}
\newunicodechar{Φ}{\ensuremath{\Phi}}
\newunicodechar{Ψ}{\ensuremath{\Psi}}
\newunicodechar{Ω}{\ensuremath{\Omega}}
\newunicodechar{↦}{\ensuremath{\mapsto}}
\newunicodechar{∈}{\ensuremath{\in}}
\newunicodechar{∉}{\ensuremath{\notin}}
\newunicodechar{∧}{\ensuremath{\wedge}}
\newunicodechar{⇔}{\ensuremath{\Leftrightarrow}}
\newunicodechar{⟺}{\ensuremath{\Longleftrightarrow}}
\newunicodechar{⇒}{\ensuremath{\Rightarrow}}
\newunicodechar{⟹}{\ensuremath{\Longrightarrow}}
\newunicodechar{∘}{\ensuremath{\circ}}
\newunicodechar{∪}{\ensuremath{\cup}}
\newunicodechar{∩}{\ensuremath{\cap}}
\newunicodechar{≡}{\ensuremath{\equiv}}
\newunicodechar{⊂}{\ensuremath{\subset}}
\newunicodechar{⊥}{\ensuremath{\perp}}
\newunicodechar{∃}{\ensuremath{\exists}}
\newunicodechar{∀}{\ensuremath{\forall}}
\newunicodechar{∛}{\ensuremath{\sqrt[3]{\,}}}
\newunicodechar{∜}{\ensuremath{\sqrt[4]{\,}}}
\newunicodechar{⃗}{\ensuremath{{}^{\to}}}
\newunicodechar{ⁿ}{\ifmmode^{n}\else\textsuperscript{\ensuremath{n}}\fi}
\newunicodechar{ˣ}{\ifmmode^{x}\else\textsuperscript{\ensuremath{x}}\fi}
\newunicodechar{ᵇ}{\ifmmode^{b}\else\textsuperscript{\ensuremath{b}}\fi}
\newunicodechar{ʳ}{\ifmmode^{r}\else\textsuperscript{\ensuremath{r}}\fi}
\newunicodechar{ᵘ}{\ifmmode^{u}\else\textsuperscript{\ensuremath{u}}\fi}
\newunicodechar{ʸ}{\ifmmode^{y}\else\textsuperscript{\ensuremath{y}}\fi}
\newunicodechar{ᵃ}{\ifmmode^{a}\else\textsuperscript{\ensuremath{a}}\fi}
\newunicodechar{ᵉ}{\ifmmode^{e}\else\textsuperscript{\ensuremath{e}}\fi}
\newunicodechar{ᵏ}{\ifmmode^{k}\else\textsuperscript{\ensuremath{k}}\fi}
\newunicodechar{ᵖ}{\ifmmode^{p}\else\textsuperscript{\ensuremath{p}}\fi}
\newunicodechar{ᴺ}{\ifmmode^{N}\else\textsuperscript{\ensuremath{N}}\fi}
\newunicodechar{ᶜ}{\ifmmode^{c}\else\textsuperscript{\ensuremath{c}}\fi}
\newunicodechar{ʰ}{\ifmmode^{h}\else\textsuperscript{\ensuremath{h}}\fi}
\newunicodechar{ᵠ}{\ifmmode^{q}\else\textsuperscript{\ensuremath{q}}\fi}
\newunicodechar{ₙ}{\ifmmode_{n}\else\textsubscript{\ensuremath{n}}\fi}
\newunicodechar{ₐ}{\ifmmode_{a}\else\textsubscript{\ensuremath{a}}\fi}
\newunicodechar{ᵢ}{\ifmmode_{i}\else\textsubscript{\ensuremath{i}}\fi}
\newunicodechar{ᵣ}{\ifmmode_{r}\else\textsubscript{\ensuremath{r}}\fi}
\newunicodechar{ₖ}{\ifmmode_{k}\else\textsubscript{\ensuremath{k}}\fi}
\newunicodechar{ᵦ}{\ifmmode_{\beta}\else\textsubscript{\ensuremath{\beta}}\fi}
\newunicodechar{ₓ}{\ifmmode_{x}\else\textsubscript{\ensuremath{x}}\fi}
\newfontfamily\popdisplay{ArchivoBlack-Regular.ttf}[Path=../../../fonts/]
\newfontfamily\poplogo{SpaceGrotesk-Bold.ttf}[Path=../../../fonts/]
\newfontfamily\popsemi{PlusJakartaSans-SemiBold.ttf}[Path=../../../fonts/]
\newfontfamily\popxbold{PlusJakartaSans-ExtraBold.ttf}[Path=../../../fonts/]
\newfontfamily\popxboldls{PlusJakartaSans-ExtraBold.ttf}[Path=../../../fonts/,LetterSpace=3.0]

\setlist[enumerate]{leftmargin=1.7em,itemsep=5pt,topsep=4pt}
\setlist[itemize]{leftmargin=1.7em,itemsep=4pt,topsep=4pt,label={\color{poporange}\textbullet}}
\setlength{\parindent}{0pt}
\setlength{\parskip}{5pt}
\renewcommand{\arraystretch}{1.25}

% pgfplots : style Pop par défaut
\pgfplotsset{
  every axis/.append style={axis line style={popink,line width=1pt},tick style={popink},
    grid style={popline,line width=0.5pt},label style={font=\small},tick label style={font=\footnotesize}},
  popcurve/.style={poporange,line width=1.8pt,no markers,smooth},
}
% Même style disponible dans un \draw TikZ simple (hors pgfplots)
\tikzset{popcurve/.style={poporange,line width=1.8pt}}
\tikzset{popgrid/.style={popline,very thin}}
\colorlet{popcurve}{poporange} % le nom du style est parfois utilisé comme couleur

% ============================================================
% Pill / squiggle
% ============================================================
\newcommand{\poppill}[3]{%
  \tikz[baseline=-0.55ex]\node[draw=popink,line width=1.2pt,rounded corners=2.6pt,%
    fill=#2,inner xsep=6.5pt,inner ysep=3pt]{%
    \popxboldls\fontsize{7}{7}\selectfont\color{#3}\MakeUppercase{#1}};%
}
\newcommand{\popsquiggle}[1]{%
  \tikz[baseline=-0.4ex]{
    \draw[#1,line width=2.6pt,line cap=round]
      plot [smooth] coordinates {(0,0.09) (0.34,0.01) (0.68,0.21) (1.02,0.01) (1.36,0.10)};
  }%
}
% Mot-clé surligné (marqueur orange sous le texte)
\newcommand{\cle}[1]{{\popxbold\color{popdark}#1}}

% ============================================================
% En-tête / pied
% ============================================================
\pagestyle{fancy}
\fancyhf{}
\renewcommand{\headrulewidth}{0pt}
\renewcommand{\footrulewidth}{0pt}
\lhead{\raisebox{-0.15ex}{\poplogo\fontsize{13}{13}\selectfont\color{popink}OnBuch\color{poporange}+}}
\rhead{\poppill{\DOCMATIERE\ \textbullet\ \DOCNIVEAU\ \textbullet\ Leçon \DOCLECON}{white}{popink}}
\lfoot{\popxbold\fontsize{7.6}{7.6}\selectfont\color{popmuted}\MakeUppercase{Cours signé OnBuch+}\ \textbullet\ {\popsemi\DOCTITRE}}
\rfoot{\tikz[baseline=-0.6ex]\node[rounded corners=8.5pt,fill=popink,minimum height=17pt,minimum width=17pt,inner xsep=5pt,inner sep=0]{\popxbold\fontsize{8}{8}\selectfont\color{popcream}\thepage\,/\,\pageref*{LastPage}};}

% ============================================================
% Titres
% ============================================================
\newcommand{\chaptitre}[2]{%
  \begin{tcolorbox}[enhanced,colback=poporange,colframe=popink,boxrule=2.6pt,arc=11pt,
      drop shadow=popink,left=15pt,right=15pt,top=12pt,bottom=11pt,before skip=3pt,after skip=18pt]
    \poppill{#2}{popink}{popcream}\par\vspace{9pt}
    {\raggedright\hyphenpenalty=10000\exhyphenpenalty=10000\popdisplay\fontsize{22}{24}\selectfont\color{popcream}\MakeUppercase{#1}\par}
  \end{tcolorbox}
}
% Section numérotée : pastille orange avec le numéro + titre Archivo
\newcounter{popsec}
\newcommand{\coursec}[1]{%
  \stepcounter{popsec}\setcounter{popsub}{0}%
  \par\vspace{16pt}\noindent
  \begin{minipage}{\linewidth}
  \tikz[baseline=-0.7ex]\node[draw=popink,line width=1.6pt,fill=poporange,rounded corners=4pt,minimum size=22pt,inner sep=0]{\popdisplay\fontsize{12}{12}\selectfont\color{popcream}\thepopsec};%
  \hspace{8pt}{\popdisplay\fontsize{15}{17}\selectfont\color{popink}\MakeUppercase{#1}}\par
  \vspace{4pt}\noindent{\color{poporange}\rule{36pt}{3.6pt}}
  \end{minipage}\par\nopagebreak\vspace{8pt}
}
\newcounter{popsub}
\newcommand{\courssub}[1]{%
  \stepcounter{popsub}%
  \par\vspace{9pt}\noindent{\popxbold\fontsize{11.5}{13}\selectfont\color{popink}{\color{poporange}\thepopsec.\thepopsub}\hspace{6pt}#1}\par\nopagebreak\vspace{4pt}
}

% ============================================================
% Boîtes pédagogiques — même recette : fond blanc, bordure épaisse colorée,
% ombre dure encre, badge plein. Titre optionnel [#1].
% ============================================================
\tcbset{popbase/.style={breakable,enhanced,colback=white,boxrule=2.3pt,arc=9pt,
  shadow={3pt}{-3pt}{0pt}{popink},left=14pt,right=14pt,top=11pt,bottom=11pt,before skip=11pt,after skip=15pt,
  fontupper=\popsemi\fontsize{10.6}{14.5}\selectfont\color{popink},
  fontlower=\popsemi\fontsize{10.6}{14.5}\selectfont\color{popink}}}
% Coche dessinée (le glyphe ✓ n'existe pas dans les polices du document)
\newcommand{\popcheck}[1]{\tikz[baseline=-0.5ex]\draw[#1,line width=1.6pt,line cap=round,line join=round](-0.13,0.0)--(-0.04,-0.09)--(0.13,0.1);}
\providecommand{\coche}{\popcheck{popgreen}}
\providecommand{\resultat}[1]{\par\smallskip\noindent\fcolorbox{popgreen}{popgreenL}{\parbox{\dimexpr\linewidth-2\fboxsep-2\fboxrule\relax}{\textbf{Résultat.} #1}}\par\smallskip}
\newcommand{\popboxhead}[3]{\poppill{#1}{#2}{white}\ifx\relax#3\relax\else\hspace{6pt}{\popxbold\fontsize{10.6}{12}\selectfont\color{popink}#3}\fi\par\vspace{7pt}}

\newtcolorbox{aretenir}[1][]{popbase,colframe=popgreen,before upper={\popboxhead{À retenir}{popgreen}{#1}}}
\newtcolorbox{exemplebox}[1][]{popbase,colframe=poporange,before upper={\popboxhead{Exemple}{poporange}{#1}}}
\newtcolorbox{definition}[1][]{popbase,colframe=popblue,before upper={\popboxhead{Définition}{popblue}{#1}}}
\newtcolorbox{propriete}[1][]{popbase,colframe=poppurple,before upper={\popboxhead{Propriété}{poppurple}{#1}}}
\newtcolorbox{methode}[1][]{popbase,colframe=popgold,before upper={\popboxhead{Méthode}{popgold}{#1}}}
\newtcolorbox{attention}[1][]{popbase,colframe=poppink,before upper={\popboxhead{Attention}{poppink}{#1}}}
\newtcolorbox{experience}[1][]{popbase,colframe=popblue,colback=popblueL,before upper={\popboxhead{Expérience}{popblue}{#1}}}
% « Le savais-tu ? » : carte violette pleine (contexte, culture, Cameroun)
\newtcolorbox{savaistu}[1][]{popbase,colback=poppurple,colframe=popink,
  fontupper=\popsemi\fontsize{10.4}{14.2}\selectfont\color{white},
  before upper={\poppill{Le savais-tu\,?}{white}{poppurple}\ifx\relax#1\relax\else\hspace{6pt}{\popxbold\color{white}#1}\fi\par\vspace{7pt}}}
% Exercice résolu : énoncé en haut, \tcblower puis la solution sur fond crème
\newcounter{popexo}\newcounter{popexr}
\newtcolorbox{exoresolu}[1][]{popbase,colframe=poporange,colbacklower=popcream,
  segmentation style={popink,line width=1.2pt,dashed},
  before upper={\stepcounter{popexr}\popboxhead{Exercice résolu \thepopexr}{poporange}{#1}},
  before lower={\poppill{Solution}{popink}{popcream}\par\vspace{6pt}}}
% Exercice (non résolu en place) — la correction est dans la partie Corrigés
\newtcolorbox{exercice}[1][]{popbase,colframe=popink,
  before upper={\stepcounter{popexo}\popboxhead{Exercice \thepopexo}{popink}{#1}}}
% Corrigé
\newtcolorbox{corrige}[1][]{popbase,colframe=popgreen,colback=popgreenL,
  before upper={\popboxhead{Corrigé}{popgreen}{#1}}}
% Objectifs / prérequis / bilan
\newtcolorbox{objectifs}{popbase,colframe=popink,colback=poporangeL,
  before upper={\popboxhead{Objectifs de la leçon}{popink}{}}}
\newtcolorbox{prerequis}{popbase,colframe=popink,colback=popblueL,
  before upper={\popboxhead{Prérequis}{popblue}{}}}
\newtcolorbox{bilan}{popbase,colframe=popink,colback=white,boxrule=2.8pt,
  before upper={\popboxhead{Fiche bilan}{poporange}{L'essentiel à connaître pour l'examen}}}
% Figure encadrée avec légende
\newtcolorbox{popfigure}[1][]{enhanced,breakable=false,colback=white,colframe=popline,boxrule=1.4pt,arc=8pt,
  left=10pt,right=10pt,top=10pt,bottom=8pt,before skip=10pt,after skip=12pt,halign=center,
  lower separated=false,halign lower=center,
  fontlower=\popsemi\fontsize{9}{11.5}\selectfont\color{popmuted},
  #1}
\newcommand{\legende}[1]{\par\vspace{6pt}{\popsemi\fontsize{9}{11.5}\selectfont\color{popmuted}#1}}

% ============================================================
% Signature OnBuch+
% ============================================================
\newcommand{\poplogoinline}[1]{{\poplogo\fontsize{#1}{#1}\selectfont\color{popink}OnBuch\color{poporange}+}}
\newcommand{\signatureonbuch}{%
  \begin{tcolorbox}[enhanced,colback=popink,colframe=popink,boxrule=2.2pt,arc=10pt,
      drop shadow=poporange,left=14pt,right=14pt,top=10pt,bottom=10pt,before skip=0pt,after skip=0pt]
    \tikz[baseline=-0.7ex]\node[circle,fill=popgreen,draw=popcream,line width=1.3pt,minimum size=22pt,inner sep=0]{\popcheck{white}};%
    \hspace{9pt}\begin{minipage}[c]{0.62\linewidth}
      {\popxbold\fontsize{12}{13}\selectfont\color{popcream}Signé\ \textbullet\ {\poplogo OnBuch\color{poporange}+}}\par\vspace{2pt}
      {\raggedright\popsemi\fontsize{8}{10.5}\selectfont\color{popline}\MakeUppercase{Équipe pédagogique OnBuch+}\\\MakeUppercase{Conforme au programme officiel MINESEC}\par}
    \end{minipage}\hfill
    {\poplogo\fontsize{22}{22}\selectfont\color{popcream}OnBuch\color{poporange}+}
  \end{tcolorbox}%
}

% ============================================================
% Page de garde
% #1 titre · #2 matière · #3 niveau · #4 module · #5 n° leçon · #6 durée ·
% #7 description / objectifs (texte ou itemize)
% ============================================================
\newcommand{\pagegarde}[7]{%
  \thispagestyle{empty}
  \newgeometry{a4paper,margin=1.7cm,top=1.6cm,bottom=1.4cm}%
  \begin{tikzpicture}[remember picture,overlay]
    \fill[poporange!18] ([xshift=-3.2cm,yshift=-3.6cm]current page.north east) circle (4.6cm);
    \fill[poppurple!14] ([xshift=2.4cm,yshift=7.6cm]current page.south west) circle (3.8cm);
  \end{tikzpicture}%
  {\poplogo\fontsize{30}{30}\selectfont\color{popink}OnBuch\color{poporange}+}\hfill
  \raisebox{0.6ex}{\poppill{Cours signé \textbullet\ Premium}{popink}{popcream}}\par
  \vspace{1pt}\popsquiggle{poppurple}\par
  \vspace{8pt}
  \begin{tcolorbox}[enhanced,colback=poporange,colframe=popink,boxrule=2.8pt,arc=13pt,
      drop shadow=popink,left=18pt,right=18pt,top=12pt,bottom=13pt,after skip=10pt]
    \poppill{#2 \textbullet\ #3}{popink}{popcream}\hspace{5pt}\poppill{#4}{white}{popink}\par\vspace{8pt}
    {\popdisplay\fontsize{14}{14}\selectfont\color{popink}LEÇON #5}\par\vspace{4pt}
    {\raggedright\hyphenpenalty=10000\exhyphenpenalty=10000\popdisplay\fontsize{25}{27}\selectfont\color{popcream}\MakeUppercase{#1}\par}
    \vspace{8pt}
    {\popxbold\fontsize{9.5}{9.5}\selectfont\color{popink}\MakeUppercase{Durée conseillée : #6}}
  \end{tcolorbox}
  \begin{tcolorbox}[enhanced,colback=white,colframe=popink,boxrule=2.2pt,arc=10pt,
      drop shadow=popink,left=16pt,right=16pt,top=10pt,bottom=10pt,after skip=8pt]
    \poppill{Dans cette leçon}{poppurple}{white}\par\vspace{5pt}
    \popsemi\fontsize{9.3}{12.6}\selectfont\color{popink}#7
  \end{tcolorbox}
  \noindent\begin{minipage}[t]{0.487\linewidth}
    \begin{tcolorbox}[enhanced,colback=popgreen,colframe=popink,boxrule=2.2pt,arc=10pt,
        drop shadow=popink,left=11pt,right=11pt,top=10pt,bottom=10pt,height=3.1cm,valign=center]
      \tcbox[colback=white,colframe=popink,boxrule=1.3pt,arc=5pt,boxsep=0pt,nobeforeafter,
        left=3pt,right=3pt,top=3pt,bottom=3pt,valign=center]{\qrcode[height=1.9cm]{https://play.google.com/store/apps/details?id=com.luvvix.onbuch&hl=fr}}%
      \hspace{8pt}\begin{minipage}[c]{0.5\linewidth}
        {\popdisplay\fontsize{11}{12.5}\selectfont\color{white}\MakeUppercase{Télécharge OnBuch}}\par\vspace{4pt}
        {\raggedright\popsemi\fontsize{8.4}{11}\selectfont\color{white}Gratuit \textbullet\ Google Play\\Tuteur IA, annales, concours, résultats\par}
      \end{minipage}
    \end{tcolorbox}
  \end{minipage}\hfill
  \begin{minipage}[t]{0.487\linewidth}
    \begin{tcolorbox}[enhanced,colback=poppurple,colframe=popink,boxrule=2.2pt,arc=10pt,
        drop shadow=popink,left=11pt,right=11pt,top=10pt,bottom=10pt,height=3.1cm,valign=center]
      \tikz[baseline=-0.7ex]\node[circle,fill=white,draw=popink,line width=1.3pt,minimum size=1.6cm,inner sep=0]{\popdisplay\fontsize{24}{24}\selectfont\color{poppurple}+};%
      \hspace{8pt}\begin{minipage}[c]{0.6\linewidth}
        {\popdisplay\fontsize{11}{12.5}\selectfont\color{white}\MakeUppercase{Passe à OnBuch+}}\par\vspace{4pt}
        {\raggedright\popsemi\fontsize{8.4}{11}\selectfont\color{white}Tous les cours signés, Tuteur IA illimité, fiches PDF — onglet OnBuch+ de l'app\par}
      \end{minipage}
    \end{tcolorbox}
  \end{minipage}
  \vspace{8pt}
  \popcontenu
  \vfill
  \signatureonbuch
  \restoregeometry
  \newpage
}

% Rangée « Ce cours contient » (page de garde)
\newcommand{\poptile}[3]{% #1 couleur #2 gros texte #3 légende
  \begin{tcolorbox}[enhanced,colback=#1,colframe=popink,boxrule=2pt,arc=9pt,shadow={2.5pt}{-2.5pt}{0pt}{popink},
      left=6pt,right=6pt,top=9pt,bottom=9pt,halign=center,valign=center,height=1.75cm,width=\linewidth,nobeforeafter]
    {\popdisplay\fontsize{12.5}{13}\selectfont\color{white}\MakeUppercase{#2}}\par\vspace{3pt}
    {\centering\popsemi\fontsize{7.8}{9.5}\selectfont\color{white}#3\par}
  \end{tcolorbox}}
\newcommand{\popcontenu}{%
  \par\noindent{\popxboldls\fontsize{7.5}{7.5}\selectfont\color{popmuted}CE COURS SIGNÉ CONTIENT}\par\vspace{7pt}
  \noindent\begin{minipage}[t]{0.235\linewidth}\poptile{popblue}{Cours}{complet, illustré, pas à pas}\end{minipage}\hfill
  \begin{minipage}[t]{0.235\linewidth}\poptile{poporange}{Méthodes}{et exercices résolus}\end{minipage}\hfill
  \begin{minipage}[t]{0.235\linewidth}\poptile{poppink}{Exercices}{type examen + corrigés}\end{minipage}\hfill
  \begin{minipage}[t]{0.235\linewidth}\poptile{popgreen}{Fiche bilan}{l'essentiel à retenir}\end{minipage}\par}

% Sommaire
\newtcolorbox{sommaire}{enhanced,colback=white,colframe=popink,boxrule=2.2pt,arc=10pt,
  drop shadow=popink,left=18pt,right=18pt,top=13pt,bottom=13pt,before skip=6pt,after skip=20pt,
  before upper={\poppill{Sommaire}{poppurple}{white}\par\vspace{9pt}\popsemi\fontsize{10.6}{14.5}\selectfont\color{popink}}}

% Signature de fin de cours — poussée tout en bas de la dernière page
\newcommand{\findecours}{%
  \par\vfill\nopagebreak
  \begin{center}\popsquiggle{poporange}\end{center}\vspace{4pt}
  \signatureonbuch
}
\AtBeginDocument{\def\llbracket{\mathopen{[\mkern-3.5mu[}}\def\rrbracket{\mathclose{]\mkern-3.5mu]}}} % intervalles d'entiers (glyphe ⟦ absent de la police)
% Glyphes absents de Fira Math : redéfinis à partir de glyphes disponibles
\AtBeginDocument{%
  \def\setminus{\mathbin{\backslash}}\let\smallsetminus\setminus
  \def\vdots{\mathord{\vbox{\baselineskip3.2pt\lineskiplimit0pt\kern4pt\hbox{.}\hbox{.}\hbox{.}}}}%
  \def\ddots{\mathinner{\mkern1mu\raise6.5pt\hbox{.}\mkern2mu\raise3.6pt\hbox{.}\mkern2mu\raise0.7pt\hbox{.}\mkern1mu}}%
  \def\triangle{\mathord{\scalebox{0.9}{$\symup{\Delta}$}}}%
  \def\nabla{\mathord{\raisebox{\depth}{\scalebox{1}[-1]{$\symup{\Delta}$}}}}%
  \def\checkmark{\popcheck{popdark}}%
  \def\star{\mathbin{\ast}}\def\bigstar{\mathord{\ast}}%
  \def\diamond{\mathbin{\scalebox{0.6}{$\Diamond$}}}%
  \def\bigcup{\mathop{\mathchoice{\raisebox{-0.3em}{\scalebox{1.7}{$\cup$}}}{\scalebox{1.25}{$\cup$}}{\cup}{\cup}}\displaylimits}%
  \def\bigcap{\mathop{\mathchoice{\raisebox{-0.3em}{\scalebox{1.7}{$\cap$}}}{\scalebox{1.25}{$\cap$}}{\cap}{\cap}}\displaylimits}%
  \def\lesseqqgtr{\mathrel{\lessgtr}}\def\bigoplus{\mathop{\oplus}\displaylimits}%
}
\providecommand{\ssi}{si et seulement si} % abréviation fréquente
\AtBeginDocument{\providecommand{\wideparen}{\overparen}} % arc de cercle

\begin{document}

\pagegarde{La Conscience et l’Inconscient}{Philosophie}{Tle Tech}{Philosophie – classe de Terminale technique}{N18}{2 séances (fiche de progression harmonisée)}{Cette leçon étudie la conscience comme pouvoir de connaissance et de maîtrise de soi, en interrogeant sa nature, sa valeur et ses limites. Elle présente ensuite l'hypothèse de l'inconscient freudien et discute la valeur et les limites de la psychanalyse comme science de l'homme.
\vspace{4pt}
\begin{multicols}{2}\small
\begin{itemize}
\item À la fin de cette leçon, je sais définir la conscience dans ses deux sens (psychologique et moral)
\item À la fin de cette leçon, je sais expliquer la nature de la conscience (immédiateté, réflexivité, intentionnalité)
\item À la fin de cette leçon, je sais montrer la valeur de la conscience (connaissance de soi, liberté, responsabilité morale)
\item À la fin de cette leçon, je sais identifier les limites de la conscience (illusions, erreurs, part obscure du psychisme)
\item À la fin de cette leçon, je sais définir l'inconscient et décrire ses manifestations (rêves, lapsus, actes manqués, symptômes)
\item À la fin de cette leçon, je sais présenter la topique freudienne (ça, moi, surmoi) et les méthodes de la psychanalyse
\end{itemize}\end{multicols}}

\begin{sommaire}
\begin{multicols}{2}
\begin{enumerate}
\item Introduction : problématique et définitions
\item La nature de la conscience
\item La valeur de la conscience
\item Les limites de la conscience
\item L'inconscient freudien
\item Valeur et limites de la psychanalyse
\item Activité d'intégration
\item Méthodes et exercices résolus
\item Exercices
\item Corrigés des exercices
\item Fiche bilan
\end{enumerate}
\end{multicols}
\end{sommaire}

\begin{objectifs}
\begin{itemize}
\item À la fin de cette leçon, je sais définir la conscience dans ses deux sens (psychologique et moral)
\item À la fin de cette leçon, je sais expliquer la nature de la conscience (immédiateté, réflexivité, intentionnalité)
\item À la fin de cette leçon, je sais montrer la valeur de la conscience (connaissance de soi, liberté, responsabilité morale)
\item À la fin de cette leçon, je sais identifier les limites de la conscience (illusions, erreurs, part obscure du psychisme)
\item À la fin de cette leçon, je sais définir l'inconscient et décrire ses manifestations (rêves, lapsus, actes manqués, symptômes)
\item À la fin de cette leçon, je sais présenter la topique freudienne (ça, moi, surmoi) et les méthodes de la psychanalyse
\item À la fin de cette leçon, je sais évaluer la valeur et les limites de la psychanalyse à partir d'arguments et de critiques
\item À la fin de cette leçon, je sais appliquer ces notions à des situations concrètes de la vie courante camerounaise
\item À la fin de cette leçon, je sais rédiger et justifier une démarche argumentée de type dissertation
\end{itemize}
\end{objectifs}

\begin{prerequis}
\begin{itemize}
\item Notion de sujet et de personne (classe de Première)
\item Distinction entre connaissance et opinion
\item Notions de liberté et de responsabilité (leçons antérieures de Terminale)
\item Vocabulaire de base : psychisme, comportement, désir, refoulement (sens courant)
\end{itemize}
\end{prerequis}

\newpage

\coursec{Introduction : problématique et définitions}

\courssub{Situation d'entrée : deux scènes de la vie quotidienne}

Commençons par deux scènes que tout élève camerounais peut observer autour de lui.

\textbf{Première scène.} À Yaoundé, un élève de Terminale technique rentre à la maison sans son bulletin trimestriel. Interrogé par son père, il jure, la main sur le cœur, qu'il a \og oublié sans faire exprès \fg{} de le rapporter du lycée. Pourtant, ses notes de mathématiques sont mauvaises, et il redoutait la colère paternelle. A-t-il vraiment oublié ? Ou bien cet oubli était-il, sans qu'il le sache lui-même, une manière commode d'échapper à la sanction ?

\textbf{Deuxième scène.} À l'hôpital régional de Bafoussam, une patiente raconte au médecin un rêve étrange : elle traverse un pont qui s'effondre sous ses pieds, et elle se réveille en sueur. Le médecin, formé à la psychologie, n'y voit pas un simple hasard : il interprète ce rêve comme le signe d'une angoisse cachée, liée à un examen que la patiente doit passer et dont elle refuse de parler éveillée.

Ces deux situations posent une question vertigineuse : \emph{sommes-nous vraiment maîtres de nos pensées et de nos actes ?} L'élève de Yaoundé croyait sincèrement avoir oublié ; la patiente de Bafoussam ignorait sa propre angoisse. Dans les deux cas, quelque chose semble se passer \emph{en eux} sans qu'ils en aient conscience. La conscience suffit-elle à expliquer tout notre comportement, ou existe-t-il en nous une part obscure, que Freud appellera \cle{l'inconscient}, qui agit à notre insu ?

Cette leçon interroge d'abord la \cle{nature}, la \cle{valeur} et les \cle{limites} de la conscience, puis évalue la théorie freudienne de l'inconscient et la méthode qu'elle fonde : la \cle{psychanalyse}.

\courssub{Étymologie : que signifient « conscience » et « inconscient » ?}

Avant de philosopher, il faut comprendre les mots. L'étymologie n'est pas un ornement : elle révèle le sens profond d'une notion.

Le mot \cle{conscience} vient du latin \emph{conscientia}, lui-même formé du préfixe \emph{cum} (\og avec \fg) et du verbe \emph{scire} (\og connaître \fg). La conscience, c'est donc littéralement le fait de \og \textbf{connaître avec} \fg{} : connaître avec soi-même, être présent à ce que l'on fait, à ce que l'on pense, à ce que l'on éprouve. Celui qui a conscience de son acte le connaît \emph{en même temps} qu'il l'accomplit : il s'accompagne lui-même, pour ainsi dire, comme un témoin intérieur.

Le mot \cle{inconscient} est formé du préfixe négatif \emph{in-} ajouté à \og conscient \fg. Il désigne donc, dans un premier sens simplement négatif, \textbf{ce qui échappe à la conscience} : tout ce qui se passe dans notre vie psychique sans que nous le sachions. Mais chez Freud, comme nous le verrons dans la suite de la leçon, l'inconscient ne se contente pas d'être une absence de conscience : il devient une \emph{réalité active}, une véritable \og seconde vie \fg{} psychique qui produit des effets (rêves, oublis, actes manqués, symptômes).

\begin{attention}[Bien comprendre le préfixe « in- »]
Dire qu'un contenu est \og inconscient \fg{} ne signifie pas qu'il n'existe pas. Cela signifie qu'il existe \emph{sans être connu du sujet}. L'erreur fréquente consiste à croire que l'inconscient est le néant psychique : au contraire, pour Freud, c'est la partie la plus vaste et la plus active de notre vie mentale.
\end{attention}

\courssub{Les deux sens du mot « conscience »}

L'observation du langage courant révèle que le mot \og conscience \fg{} possède deux sens distincts qu'il faut soigneusement distinguer, car toute la dissertation peut être faussée si on les confond.

\textbf{Premier sens : la conscience psychologique.} Quand je dis \og je suis conscient qu'il pleut \fg, \og j'ai conscience d'avoir mal à la tête \fg, j'emploie le mot au sens psychologique : je désigne le fait d'être présent à moi-même et au monde, de savoir ce que je fais et ce que je ressens pendant que je le fais et le ressens.

\begin{definition}[La conscience psychologique]
La \cle{conscience psychologique} est la \textbf{connaissance immédiate que le sujet a de ses états et de ses actes}. Elle est la présence du sujet à lui-même et au monde : je sais que je pense, que je veux, que je souffre, au moment même où je pense, veux ou souffre.
\end{definition}

\textbf{Deuxième sens : la conscience morale.} Quand je dis \og ma conscience me reproche ce mensonge \fg{} ou \og il a la conscience tranquille \fg, j'emploie le mot au sens moral : la conscience devient un \emph{juge intérieur} qui évalue mes actions selon le bien et le mal, qui approuve ou condamne.

\begin{definition}[La conscience morale]
La \cle{conscience morale} est la \textbf{faculté de juger ses propres actions selon le bien et le mal}. Elle se manifeste notamment par le remords (quand elle condamne un acte accompli) et par la bonne conscience (quand elle approuve).
\end{definition}

\begin{exemplebox}[Illustrations camerounaises des deux sens]
\textbf{Conscience psychologique :} le chauffeur de taxi-brousse sur l'axe Yaoundé--Douala est conscient de la route, des autres véhicules, de sa fatigue ; l'élève qui compose est conscient de l'heure qui passe et de la difficulté du sujet.

\textbf{Conscience morale :} le commerçant du marché Mokolo qui hésite à tricher sur le poids des marchandises entend une \og voix intérieure \fg{} qui le lui reproche ; le fonctionnaire qui refuse un pot-de-vin éprouve ensuite la satisfaction d'une conscience tranquille.
\end{exemplebox}

\begin{popfigure}
\begin{center}
\begin{tikzpicture}[
  box/.style={draw=popblue, thick, rounded corners, fill=popblueL, text width=5.6cm, align=center, inner sep=6pt},
  ex/.style={draw=popmuted, rounded corners, fill=popcream, text width=5.6cm, align=left, inner sep=5pt, font=\small},
  titre/.style={font=\bfseries\large, text=popdark}
]
\node[titre] at (-3.6,3.4) {Conscience psychologique};
\node[titre] at (3.6,3.4) {Conscience morale};
\node[box, align=center] at (-3.6,2.1){Connaissance immédiate de ses états et de ses actes\\ \emph{(être présent à soi et au monde)}};
\node[box, draw=poporange, fill=poporangeL, align=center] at (3.6,2.1){Faculté de juger ses actions selon le bien et le mal\\ \emph{(juge intérieur)}};
\node[ex, align=center] at (-3.6,0.2){\textbf{Exemples :}\\ -- L'élève conscient de sa fatigue pendant l'examen.\\ -- Le chauffeur attentif à la route Yaoundé--Douala.};
\node[ex, align=center] at (3.6,0.2){\textbf{Exemples :}\\ -- Le remords du commerçant qui a triché.\\ -- La conscience tranquille du fonctionnaire honnête.};
\draw[->, thick, popmuted] (-3.6,1.65) -- (-3.6,0.95);
\draw[->, thick, popmuted] (3.6,1.65) -- (3.6,0.95);
\end{tikzpicture}
\end{center}
\legende{Les deux sens du mot \og conscience \fg{} : à gauche, la conscience psychologique (connaître ses états et ses actes) ; à droite, la conscience morale (juger ses actes selon le bien et le mal), avec des exemples de la vie quotidienne camerounaise.}
\end{popfigure}

\begin{attention}[Erreur fréquente : conscience et « bonne conscience »]
Ne confonds jamais la \textbf{conscience} (faculté psychique de connaître ou de juger) avec la \textbf{bonne conscience} (satisfaction morale de celui qui estime avoir bien agi). La conscience est une \emph{faculté} ; la bonne conscience est un \emph{état}, un sentiment d'approbation de soi. On peut posséder la conscience morale et avoir mauvaise conscience (remords). Dans une copie, écrire \og l'homme est conscience \fg{} pour dire \og l'homme a bonne conscience \fg{} est une faute de sens qui affaiblit tout le raisonnement.
\end{attention}

\courssub{Construction de la problématique}

Partons de nos deux situations d'entrée. Si l'élève de Yaoundé a \emph{sincèrement} oublié son bulletin, alors une partie de sa conduite (l'oubli) lui échappe : il n'est pas transparent à lui-même. Si le rêve de la patiente de Bafoussam révèle une angoisse qu'elle ignorait, alors sa vie psychique contient des contenus cachés, que la conscience seule ne peut pas atteindre. Mais inversement, l'expérience ordinaire semble montrer que je suis le mieux placé pour savoir ce que je pense : qui connaît ma douleur, mes intentions, mieux que moi ?

Nous voici devant une \textbf{contradiction} : d'un côté, la conscience paraît être toute la vie psychique (être un esprit, c'est se connaître soi-même) ; de l'autre, les faits (oublis, rêves, actes manqués) suggèrent qu'une part de notre psychisme nous est fermée. C'est de cette contradiction que naît la problématique de la leçon :

\begin{aretenir}[Problématique de la leçon]
La conscience épuise-t-elle la réalité psychique ? Autrement dit : \textbf{l'homme est-il transparent à lui-même}, ou existe-t-il en lui un inconscient qui détermine ses pensées et ses actes à son insu ? Et si cet inconscient existe, la psychanalyse est-elle une méthode valable pour le connaître ?
\end{aretenir}

\begin{popfigure}
\begin{center}
\begin{tikzpicture}[
  grow=right,
  level 1/.style={sibling distance=30mm, level distance=38mm},
  level 2/.style={sibling distance=16mm, level distance=42mm},
  racine/.style={draw=popdark, thick, fill=popblue, text=white, rounded corners, align=center, inner sep=5pt, font=\bfseries},
  n1/.style={draw=poporange, thick, fill=poporangeL, rounded corners, align=center, text width=3.2cm, inner sep=4pt},
  n2/.style={draw=popmuted, fill=popcream, rounded corners, align=center, text width=3.4cm, inner sep=4pt, font=\small},
  edge from parent/.style={draw=popmuted, thick, ->}
]
\node[racine, align=center]{La conscience\\ épuise-t-elle\\ le psychisme ?}
  child { node[n1]{Non : l'inconscient} 
    child { node[n2]{Rêves, oublis, actes manqués (Freud)} }
    child { node[n2]{Psychanalyse : valeur et limites} } }
  child { node[n1]{Oui : transparence à soi}
    child { node[n2]{Je sais ce que je pense (Descartes)} }
    child { node[n2]{Valeur de la conscience : liberté, responsabilité} } };
\end{tikzpicture}
\end{center}
\legende{Carte mentale de la problématique : la question centrale (la conscience épuise-t-elle le psychisme ?) se divise en deux réponses possibles -- la transparence de l'homme à lui-même ou l'existence de l'inconscient -- dont la leçon examinera la valeur et les limites.}
\end{popfigure}

\begin{methode}[Construire une problématique à partir d'une situation de vie courante]
Pour transformer une situation concrète (un oubli, un rêve, un mensonge) en problématique de dissertation, suis ces étapes :
\begin{enumerate}
\item \textbf{Décrire précisément la situation} : qui fait quoi, dans quelles circonstances ? (L'élève oublie son bulletin ; la patiente rêve d'un pont qui s'effondre.)
\item \textbf{Dégager le fait étonnant} : qu'est-ce qui surprend et appelle une explication ? (On peut ignorer ses propres motifs.)
\item \textbf{Identifier les notions philosophiques en jeu} : ici, conscience, inconscient, connaissance de soi.
\item \textbf{Formuler la contradiction} : opposer deux réponses également défendables (je me connais parfaitement / une part de moi m'échappe).
\item \textbf{Transformer la contradiction en question} : rédiger une interrogation générale (\og La conscience épuise-t-elle la réalité psychique ? \fg), puis des questions secondaires qui annoncent le plan.
\end{enumerate}
\end{methode}

\begin{exoresolu}[De la situation au sujet de dissertation]
\textbf{Énoncé.} Un élève de Bafenda constate qu'il a insulté son camarade \og sans le vouloir \fg{}, par un mot qui lui a \og échappé \fg. Construis, à partir de cette situation, une problématique complète en trois questions.
\tcblower
\textbf{Solution détaillée.}
\begin{enumerate}
\item \emph{Fait étonnant :} la parole du sujet exprime un contenu (l'hostilité) que sa conscience n'avait pas prévu ni voulu : le mot a \og échappé \fg.
\item \emph{Notions en jeu :} conscience (volonté, maîtrise de soi) et inconscient (motivation cachée).
\item \emph{Contradiction :} soit le sujet est maître de sa parole et le mot était voulu ; soit sa parole révèle un désir qui agit en lui sans qu'il le sache.
\item \emph{Problématique en trois questions :} (1) La conscience nous rend-elle maîtres de nos paroles et de nos actes ? (2) Les actes qui nous échappent ne révèlent-ils pas un inconscient ? (3) La psychanalyse permet-elle de connaître scientifiquement cet inconscient ?
\end{enumerate}
Ces trois questions correspondent exactement aux trois parties du plan de la leçon.
\end{exoresolu}

\courssub{Annonce du plan}

Pour répondre à la problématique, la leçon suivra trois étapes, qui correspondent aux savoirs du programme officiel :

\begin{enumerate}
\item \textbf{Première partie : la conscience.} Nous étudierons d'abord sa \emph{nature} (qu'est-ce que la conscience ? comment se rapporte-t-elle à elle-même et au monde ?), puis sa \emph{valeur} (ce qu'elle rend possible : la connaissance de soi, la liberté, la responsabilité morale), enfin ses \emph{limites} (ce qu'elle ne peut pas expliquer : les oublis, les rêves, les actes manqués, les maladies nerveuses).
\item \textbf{Deuxième partie : l'inconscient.} Face aux limites de la conscience, nous présenterons l'hypothèse de Freud : l'existence d'un psychisme inconscient, actif et structuré, qui se manifeste dans les rêves, les lapsus et les symptômes.
\item \textbf{Troisième partie : valeur et limites de la psychanalyse.} Nous évaluerons enfin la méthode psychanalytique : qu'apporte-t-elle à la connaissance de l'homme ? Peut-on la considérer comme une science ? Quelles critiques peut-on lui adresser ?
\end{enumerate}

\begin{aretenir}[L'essentiel de l'introduction]
\begin{itemize}
\item \textbf{Conscience} (latin \emph{cum-scire}, \og connaître avec \fg) : deux sens à distinguer -- \emph{psychologique} (connaissance immédiate de ses états et de ses actes) et \emph{morale} (faculté de juger ses actions selon le bien et le mal).
\item \textbf{Inconscient} : ce qui échappe à la conscience ; chez Freud, une réalité psychique active qui agit à l'insu du sujet.
\item \textbf{Problématique :} la conscience épuise-t-elle la réalité psychique ? L'homme est-il transparent à lui-même ?
\item \textbf{Piège à éviter :} confondre la conscience (faculté) avec la bonne conscience (satisfaction morale).
\end{itemize}
\end{aretenir}

\begin{exercice}[Pour vérifier ta compréhension]
\begin{enumerate}
\item Donne l'étymologie du mot \og conscience \fg{} et explique ce qu'elle révèle du sens de la notion.
\item À l'aide d'un exemple tiré de ta vie quotidienne, distingue conscience psychologique et conscience morale.
\item Un candidat écrit : \og L'inconscient, c'est l'absence de toute vie psychique. \fg{} Corrige cette erreur.
\item À partir de la situation suivante : \og une élève rêve qu'elle arrive en retard à l'examen du Baccalauréat technique \fg, construis une problématique en deux questions.
\end{enumerate}
\end{exercice}

\begin{corrige}[Exercice : Pour vérifier ta compréhension]
\begin{enumerate}
\item \emph{Conscience} vient du latin \emph{conscientia}, formé de \emph{cum} (\og avec \fg) et \emph{scire} (\og connaître \fg) : \og connaître avec \fg. Cela révèle que la conscience est une connaissance qui accompagne le sujet : celui-ci se connaît lui-même en même temps qu'il agit, pense ou éprouve.
\item Exemple attendu du type : \og je suis conscient que j'ai faim \fg{} (conscience psychologique : connaissance immédiate d'un état) ; \og ma conscience me reproche d'avoir copié sur mon voisin \fg{} (conscience morale : jugement de l'acte selon le bien et le mal).
\item Faux : l'inconscient n'est pas l'absence de vie psychique, mais la partie de la vie psychique qui \emph{échappe à la conscience}. Chez Freud, il est même la partie la plus active du psychisme, qui produit rêves, oublis et symptômes.
\item Questions possibles : (1) Le rêve est-il un simple désordre de l'esprit ou a-t-il un sens ? (2) Le rêve révèle-t-il un contenu inconscient (ici, l'angoisse de l'échec) que la conscience ignore ?
\end{enumerate}
\end{corrige}

\coursec{La nature de la conscience}

\noindent Après avoir posé, en introduction, la problématique générale — \og~Qu'est-ce que la conscience ? Est-elle une lumière qui éclaire tout, ou un îlot étroit sur un océan obscur ?~\fg{} — et distingué les sens courants (conscience morale, conscience psychologique, conscience de soi), nous entrons maintenant dans l'analyse rigoureuse de la \textbf{nature} de la conscience. Cette section vise à montrer que la conscience n'est pas un simple réceptacle passif, mais une activité structurée, intentionnelle et temporelle, dont le \emph{cogito} cartésien offre le paradigme fondateur.

\courssub{La conscience comme donnée immédiate : le « je pense » cartésien et kantien}

\begin{experience}[La démarche du doute hyperbolique et la découverte du cogito]
\textbf{Matériel conceptuel :} Les \emph{Méditations métaphysiques} (1641), notamment la \textsc{IIe Méditation}.
\textbf{Protocole (reconstitution rationnelle) :}
\begin{enumerate}
    \item \textbf{Doute universel :} Je décide de rejeter comme faux tout ce qui peut être mis en doute (sens, corps, monde extérieur, vérités mathématiques — hypothèse du \emph{malin génie}).
    \item \textbf{Recherche de l'indubitable :} Y a-t-il quelque chose qui résiste à ce doute ? Même si je doute de tout, \emph{je ne peux pas douter que je doute}.
    \item \textbf{Identification :} Douter, c'est penser. Donc, \emph{je pense} est une certitude absolue.
    \item \textbf{Conclusion ontologique :} Si je pense, \emph{je suis} (existe) en tant que \emph{res cogitans} (chose pensante).
\end{enumerate}
\textbf{Observation :} Le \og~je pense~fg{} n'est pas une déduction logique (un syllogisme : \og~Tout ce qui pense existe ; je pense ; donc je suis~fg{}), car le majeur serait incertain. C'est une \textbf{intuition immédiate} : l'acte de penser \emph{emporte avec lui} l'existence du sujet.
\textbf{Interprétation :} La conscience se révèle comme la \textbf{première certitude}, le point d'Archimède de toute connaissance. Elle est \emph{donnée immédiate} : je ne \emph{connais} pas ma conscience par déduction, je \emph{suis} ma conscience en train de se vivre.
\end{experience}

\begin{exemplebox}[Texte de référence pour le Bac — Descartes, \emph{Méditations métaphysiques}, IIe Méditation (1641)]
\og~Mais j'ai beau supposer qu'il y a un Dieu tout-puissant, ou bien un génie malin, qui emploie toute sa force et toute son industrie à me tromper, il n'y a jamais de moment où je ne sois quelque chose, tant que je penserai être quelque chose. Ainsi, après avoir bien pensé et avoir examiné toutes choses avec soin, il faut enfin conclure et tenir pour certain que cette proposition : \textbf{Je suis, j'existe}, est nécessairement vraie, toutes les fois que je la profère, ou que je la conçois en mon esprit.~\fg{}

\medskip\noindent\textbf{Exploitation :} Ce texte est \textbf{incontournable} pour le Probatoire/Bac. Il illustre :
\begin{itemize}
    \item Le caractère \emph{performatif} du cogito : \og~toutes les fois que je la profère~\fg{} (l'énonciation \emph{est} la preuve).
    \item L'indépendance à l'égard du corps et du monde (\og~supposer qu'il y a un Dieu...~\fg{}).
    \item La définition de l'homme comme \emph{res cogitans} : \og~une chose qui doute, qui entend [entend = entend par l'entendement], qui conçoit, qui affirme, qui nie, qui veut, qui ne veut pas, qui imagine aussi, qui sent~\fg{} (suite du texte).
\end{itemize}
\end{exemplebox}

\begin{definition}[Le Cogito]
Acte fondamental par lequel le sujet pensant se saisit lui-même comme \textbf{existant} dans l'instant même où il pense (doute, affirme, nie, veut). C'est la \textbf{certitude première}, indubitable, fondement de toute philosophie cartésienne. Il révèle la conscience comme \emph{présence à soi} immédiate et transparente.
\end{definition}

\begin{propriete}[La conscience comme « Je pense » accompagnant toutes les représentations (Kant)]
Pour Kant (\emph{Critique de la raison pure}, Déduction transcendantale), le \og~Je pense~fg{} est la \textbf{condition de possibilité} de toute expérience. \og~Le \emph{Je pense} doit pouvoir accompagner toutes mes représentations~fg{}. Cela signifie :
\begin{itemize}
    \item \textbf{Unité de l'aperception :} Mes représentations disparates (voir du rouge, entendre un son, ressentir une douleur) ne forment une \emph{expérience} unifiée que parce qu'elles sont rapportées à un même \og~Je~fg{}.
    \item \textbf{Spontanéité :} Ce \og~Je~fg{} n'est pas une représentation empirique (je ne me \emph{vois} pas penser), c'est l'acte spontané de synthèse qui rend l'objet possible.
\end{itemize}
\emph{Distinction cruciale Bac :} Descartes fonde l'\emph{existence} du sujet ; Kant fonde l'\emph{unité} de la connaissance.
\end{propriete}

\courssub{La réflexivité : conscience spontanée vs conscience réfléchie}

La conscience ne se contente pas d'être \emph{de} quelque chose (intentionnalité, voir infra) ; elle a le pouvoir de se \emph{replier sur elle-même}. C'est la \textbf{réflexivité}.

\begin{center}
\begin{tikzpicture}[node distance=1.8cm, >=Stealth, font=\small]
    \node[draw, rounded corners, fill=popblueL, thick, align=center] (spont){Conscience spontanée (vécue)\\(pre-réflexive)};
    \node[draw, rounded corners, fill=poporangeL, thick, right=of spont, align=center] (refle){Conscience réfléchie\\(réflexive)};
    \draw[->, thick, popblue] (spont) -- node[above] {Objet : le monde, autrui, mon corps} (refle);
    \draw[->, thick, poporange, bend left=20] (refle) to node[above, align=center]{Objet : \textbf{moi-même}\\(mes actes, mes états)} (spont);
    \node[align=center, text width=5cm, below=0.5cm of spont, text=popmuted] (ex1) {Ex. : Je lis ce texte,\\je suis absorbé par le sens.};
    \node[align=center, text width=5cm, below=0.5cm of refle, text=popmuted] (ex2) {Ex. : Je me dis :\\« Je suis en train de lire »,\\j'analyse ma compréhension.};
\end{tikzpicture}
\end{center}

\begin{propriete}[La conscience est réflexive — elle peut se doubler elle-même]
La réflexivité est la capacité de la conscience à \textbf{se prendre elle-même pour objet}.
\begin{itemize}
    \item \textbf{Conscience spontanée (pre-réflexive) :} Je suis \emph{dedans} l'acte (je vois l'arbre, je suis en colère). Je ne me \emph{dis} pas \og~je vois~fg{}, je \emph{vois}. C'est la conscience \emph{de} l'objet.
    \item \textbf{Conscience réfléchie :} Je me retourne sur mon acte mental (je sais que je vois, j'analyse ma colère). C'est la conscience \emph{de la conscience} (\og~je sais que je sais~fg{}).
\end{itemize}
\emph{Conséquence :} La réflexivité fonde l'\textbf{identité personnelle} (je me reconnais comme le même sujet à travers le temps) et la \textbf{responsabilité} (je suis auteur de mes actes parce que je peux me les imputer).
\end{propriete}

\begin{attention}[Piège fréquent : Confondre réflexivité et introspection]
\begin{itemize}
    \item \textbf{Réflexivité :} Structure \emph{ontologique} de la conscience (elle \emph{est} retour sur soi). Elle est permanente (même en lisant, j'ai une conscience pré-réfléchie de lire).
    \item \textbf{Introspection :} Méthode \emph{psychologique} (je \emph{décide} d'observer mes états intérieurs). Elle est ponctuelle, volontaire, et fausse souvent l'objet observé (le fait de m'observer être en colère modifie ma colère).
\end{itemize}
Au Bac, ne dites pas \og~la conscience est l'introspection~fg{}. Dites : \og~La conscience a une structure réflexive qui \emph{permet} l'introspection, mais la précède et la fonde.~fg{}
\end{attention}

\courssub{L'intentionnalité : toute conscience est conscience \emph{de} quelque chose (Husserl)}

Si la réflexivité est le retour sur soi, l'\textbf{intentionnalité} est l'ouverture au monde. C'est la thèse centrale de la phénoménologie (Brentano, Husserl).

\begin{popfigure}
\begin{tikzpicture}[>=Stealth, font=\small, node distance=2.5cm]
    % Sujet
    \node[draw, ellipse, fill=popgreenL, thick, minimum width=2.5cm, minimum height=1.2cm, align=center] (sujet){\textbf{Sujet Conscient}\\(Noèse)};
    % Flèche intentionnelle
    \draw[->, very thick, popgreen, decorate, decoration={snake, amplitude=1pt, segment length=8pt}] (sujet.east) -- node[above, align=center] {\textbf{Visée intentionnelle}\\(\emph{Intentio})} ++(3.5,0) node[draw, rectangle, rounded corners, fill=poppurpleL, thick, minimum width=2.5cm, minimum height=1.2cm, anchor=west, align=center] (objet){\textbf{Objet visé}\\(Noème)};
    % Structure noèse/noème
    \node[below=0.3cm of sujet, align=center, text=popmuted] {Acte : Percevoir, Juger,\\Souvenir, Désirer, Imaginer};
    \node[below=0.3cm of objet, align=center, text=popmuted] {Contenu : Cet arbre,\\« 2+2=4 », Mon anniversaire,\\La justice, Une licorne};
    % Rayons pour montrer la corrélation
    \draw[popgreen, dashed, thin] (sujet.north east) -- ++(1.5,1);
    \draw[popgreen, dashed, thin] (sujet.south east) -- ++(1.5,-1);
    \draw[poppurple, dashed, thin] (objet.north west) -- ++(-1.5,1);
    \draw[poppurple, dashed, thin] (objet.south west) -- ++(-1.5,-1);
    \node[above=1.2cm of sujet, align=center, text=popdark, font=\normalsize\bfseries] {Corrélation intentionnelle\\Sujet $\leftrightarrow$ Objet};
\end{tikzpicture}
\legende{Schéma de l'intentionnalité (Husserl) : La conscience n'est pas un « contenant » vide, elle est \textbf{acte} (noèse) qui vise un \textbf{contenu} (noème). L'objet n'a pas besoin d'exister réellement (licorne) pour être visé.}
\end{popfigure}

\begin{definition}[Intentionnalité (Husserl, \emph{Idées directrices pour une phénoménologie}, 1913)]
Propriété fondamentale de la conscience d'être \textbf{toujours conscience \emph{de} quelque chose}. Il n'y a pas de conscience \og~vide~fg{} ou \og~pure~fg{} sans objet. La structure intentionnelle unit indissociablement :
\begin{itemize}
    \item La \textbf{Noèse} : l'acte subjectif (percevoir, juger, se souvenir, imaginer, désirer, haïr).
    \item Le \textbf{Noème} : l'objet \emph{tel qu'il est visé} (l'arbre \emph{perçu}, la vérité \emph{jugée}, le passé \emph{souvenu}, l'avenir \emph{désiré}).
\end{itemize}
\emph{Exemple :} La perception de cet arbre et l'imagination de cet arbre ont le \emph{même noème} (l'arbre), mais des \emph{noèses} différentes (présence physique vs absence/imagination).
\end{definition}

\begin{aretenir}[Conséquences majeures pour la dissertation]
\begin{enumerate}
    \item \textbf{Refus du « myth of the given » (mythe du donné) :} La conscience ne reçoit pas passivement des « données » ; elle \emph{constitue} le sens de l'objet par ses actes.
    \item \textbf{Extériorité radicale :} L'objet visé \emph{transcende} l'acte de conscience. L'arbre que je perçois n'est \emph{pas} dans ma tête ; ma conscience \emph{va vers lui} (dépassement/ek-stase).
    \item \textbf{Liberté de la visée :} Je peux viser le même objet sous des modes différents (le voir, le juger, le désirer, le nier). La conscience est \emph{liberté} de varier les noèses.
\end{enumerate}
\end{aretenir}

\courssub{La conscience comme lumière et présence : opposition sujet/objet}

La tradition cartésienne et phénoménologique décrit la conscience comme une \textbf{lumière} (Lumières) ou une \textbf{présence} (Phénoménologie).

\begin{itemize}
    \item \textbf{Métaphore de la lumière (Descartes, Malebranche, Maine de Biran) :} La conscience \emph{éclaire} les objets. Sans elle, les choses sont dans l'obscurité (l'inconscient, la matière inerte). Le sujet est le \emph{foyer} lumineux.
    \item \textbf{Présence (Husserl, Sartre, Merleau-Ponty) :} La conscience n'est pas un miroir qui reflète ; elle est \emph{présence-au-monde}. \og~Toute conscience est conscience de quelque chose~fg{} (Sartre, \emph{L'Être et le Néant}) signifie qu'elle est une \textbf{néantisation} du monde : elle le nie comme pure identité pour le faire apparaître comme \emph{signification}.
\end{itemize}

\begin{attention}[Piège : Le « théâtre de la conscience » (Cartésianisme naïf)]
Ne pas représenter la conscience comme un \emph{théâtre} où défilent des images (représentations) regardées par un petit homoncule intérieur.
\begin{itemize}
    \item \textbf{Erreur :} Sujet \emph{dedans} $\leftarrow$ Représentations (écran) $\leftarrow$ Objets \emph{dehors}.
    \item \textbf{Correct (Intentionnalité) :} Sujet \textbf{=} Acte de viser $\rightarrow$ Objet \textbf{dehors} (transcendant). Il n'y a pas d'écran intérieur ; la conscience \emph{est} cette relation même.
\end{itemize}
Cette critique (Ryle, Wittgenstein, Merleau-Ponty, Sartre) est essentielle pour la partie \og~Limites de la conscience~fg{} (section 4) et pour l'inconscient.
\end{attention}

\courssub{Conscience et temporalité : la synthèse du temps (Husserl, Bergson, Heidegger)}

La conscience n'est pas un instant figé ; elle est \textbf{devenir}. Elle structure le temps. Husserl (\emph{Leçons sur la conscience intime du temps}, 1928) analyse la \textbf{conscience du temps immanent} comme une structure tripartite :

\begin{center}
\begin{tikzpicture}[font=\small, >=Stealth, node distance=1.2cm]
    % Trois boîtes principales
    \node[draw, rounded corners, fill=popblueL, thick, minimum width=2.4cm, minimum height=1.4cm, text width=2.2cm, align=center] (ret) {\textbf{Rétention}\\(Passé immédiat)\\\og~Je \emph{viens d'entendre} la note~fg{}};
    \node[draw, rounded corners, fill=popgoldL, thick, minimum width=2.4cm, minimum height=1.4cm, right=of ret, text width=2.2cm, align=center] (prim) {\textbf{Impression primordiale}\\(Présent vif)\\\og~J'entends \emph{maintenant} la note~fg{}};
    \node[draw, rounded corners, fill=poporangeL, thick, minimum width=2.4cm, minimum height=1.4cm, right=of prim, text width=2.2cm, align=center] (prot) {\textbf{Protention}\\(Avenir immédiat)\\\og~J'\emph{attends} la note suivante~fg{}};
    % Flèches
    \draw[->, thick, popblue] (ret) -- node[above] {Synthèse passive} (prim);
    \draw[->, thick, poporange] (prim) -- node[above] {Anticipation} (prot);
    % Mémoire (au-dessus)
    \node[draw, rounded corners, fill=poppurpleL, thick, minimum width=2.4cm, minimum height=1cm, above=1cm of prim, text width=2.2cm, align=center] (mem) {\textbf{Mémoire / Réminiscence}\\(Passé représenté : \og~Je me souviens d'hier~fg{})};
    \draw[->, thick, poppurple, dashed] (mem) -- node[right] {Acte de récollection} (prim);
    % Projet (en dessous)
    \node[draw, rounded corners, fill=poppinkL, thick, minimum width=2.4cm, minimum height=1cm, below=1cm of prim, text width=2.2cm, align=center] (proj) {\textbf{Projet / Volonté}\\(Avenir représenté : \og~Je prépare demain~fg{})};
    \draw[->, thick, poppink, dashed] (prim) -- node[right] {Projection intentionnelle} (proj);
\end{tikzpicture}
\end{center}

\begin{propriete}[La conscience comme synthèse temporelle]
La conscience \textbf{relie passé, présent et avenir} en une unité vivante :
\begin{itemize}
    \item \textbf{Rétention (Passé immédiat) :} Ce n'est pas un souvenir (représentation), c'est le \emph{maintien} de la note qui vient de sonner dans l'acte d'entendre la note actuelle. Sans rétention, pas de perception de mélodie, seulement bruits disjoints.
    \item \textbf{Impression primordiale (Présent) :} L'acte \emph{originaire} où le maintenant se donne. Mais ce « maintenant » a une \emph{épaisseur} (largeur spéculaire) : il contient la rétention et la protention.
    \item \textbf{Protention (Avenir immédiat) :} L'\emph{attente} vide mais déterminée (j'attends \emph{la suite} de la phrase, \emph{la chute} de l'objet). C'est l'ouverture intentionnelle vers le futur.
    \item \textbf{Mémoire et Projet (Temps représenté) :} Au-delà du flux immanent, la conscience \emph{réfléchie} pose des actes de récollection (passé) et de projection volontaire (avenir). C'est là que se joue l'\emph{identité narrative} (Ricoeur) et la \emph{liberté} (Heidegger : \emph{projet jeté}).
\end{itemize}
\end{propriete}

\begin{savaistu}[Le mot « conscience » en langues camerounaises]
En \textbf{duala}, on parle souvent de \textbf{\emph{mba}} (le cœur) ou \emph{ndob} (l'entendement, l'intelligence) pour désigner le siège de la conscience morale et psychologique. L'expression \og~\emph{mba a moni}~fg{} (le cœur voit/entend) renvoie à l'idée d'une \textbf{présence intérieure qui discerne}, proche de la \emph{conscientia} latine (\emph{cum-scire}, savoir avec soi).
En \textbf{ewondo}, \textbf{\emph{ev}} (le cœur/le foie) est le centre de la personne : \og~\emph{ev amva}~fg{} (le cœur sait) exprime la conscience comme \textbf{connaissance intime, immédiate, non médiatisée par les sens}.
Ces langues ne séparent pas radicalement \og~raison~fg{} et \og~sentiment~fg{} : la conscience est une \textbf{totalité unifiée} (cœur-esprit) qui \og~voit~fg{} la vérité morale et existentielle. Cela rejoint l'intuition bergsonienne ou phénoménologique d'une conscience \emph{vivante}, non coupée en deux (cognitif/affectif).
\end{savaistu}

\begin{exoresolu}[Analyse structurée d'une situation : « Je marche dans la rue et j'entends un bruit de freinage brusque »]
\textbf{Énoncé :} Identifiez, dans cette situation vécue, les structures de la conscience étudiées : intentionnalité, réflexivité, temporalité (rétention/protention), conscience spontanée/réfléchie.

\medskip\noindent\textbf{Solution détaillée :}
\begin{enumerate}
    \item \textbf{Intentionnalité (Noèse/Noème) :}
    \begin{itemize}
        \item Acte (Noèse) : \emph{Perception auditive} (soudaine, passive).
        \item Objet visé (Noème) : \emph{Le bruit de freinage} \emph{comme} signal de danger / comme « voiture qui freine ». Le sens (danger) est \emph{constitutif} de la perception, pas ajouté après.
    \end{itemize}
    \item \textbf{Temporalité (Flux immanent) :}
    \begin{itemize}
        \item \emph{Rétention :} Je retiens le début du cri de pneus pour l'unifier en \emph{un} bruit unique (sinon : éclats sonores disjoints).
        \item \emph{Impression primordiale :} Le « maintenant » du pic sonore.
        \item \emph{Protention :} Mon corps se tend \emph{avant} l'impact possible ; j'anticipe le choc ou l'arrêt. La protention guide ma réaction motrice (sauter sur le trottoir).
    \end{itemize}
    \item \textbf{Conscience spontanée vs Réfléchie :}
    \begin{itemize}
        \item \emph{Spontanée (pré-réflexive) :} Le saut sur le trottoir, la peur viscérale. Je \emph{suis} l'action, je ne me dis pas \og~je saute~fg{}. C'est la \emph{conscience corporelle} (Merleau-Ponty : \og~corps propre~fg{}).
        \item \emph{Réfléchie (a posteriori) :} Une fois en sécurité, je me dis : \og~J'ai eu peur, j'ai bien réagi~fg{}. Je me saisis comme \emph{sujet} de l'action passée. C'est le fondement de ma responsabilité : \og~C'est moi qui ai sauté~fg{}.
    \end{itemize}
    \item \textbf{Synthèse :} La conscience n'est pas un spectateur ; elle est \emph{engagement corporel intentionnel dans le temps}, qui ne devient \emph{objet pour soi} (réflexivité) qu'en second temps.
\end{enumerate}
\end{exoresolu}

\begin{methode}[Comment traiter une question sur « La nature de la conscience » au Bac]
\begin{enumerate}
    \item \textbf{Définir} rigoureusement : Conscience = \emph{présence intentionnelle à soi et au monde} (pas « connaissance », pas « âme »).
    \item \textbf{Structurer} en 3 piliers (plan type) :
    \begin{enumerate}
        \item \textbf{La conscience comme certitude fondatrice (Cogito) :} Immédiateté, indubitabilité, subjectivité transparente (Descartes/Kant).
        \item \textbf{La conscience comme ouverture au monde (Intentionnalité) :} Structure acte/objet, refus de l'immanence close, extériorité radicale (Husserl/Sartre).
        \item \textbf{La conscience comme temporalité vivante :} Synthèse rétention/impression/protention, unité du flux, projet d'avenir (Husserl/Bergson/Heidegger).
    \end{enumerate}
    \item \textbf{Articuler} les auteurs : Ne pas faire un catalogue. Montrer comment Kant approfondit Descartes (condition de possibilité), comment Husserl retourne le cogito vers le monde (intentionnalité), comment la temporalité unit les deux.
    \item \textbf{Problématiser} : Cette nature (transparence, intentionnalité, unité) est-elle \emph{toute} la conscience ? $\rightarrow$ Transition naturelle vers \og~Limites de la conscience~fg{} et \og~L'Inconscient~fg{}.
\end{enumerate}
\end{methode}

\begin{exercice}[Entraînement Bac — Sujet type]
\textbf{Sujet :} \og~La conscience est-elle une lumière qui éclaire tout ?~\fg{} (Dissertation)
\textbf{Consignes :}
\begin{enumerate}
    \item Analysez les termes : « lumière » (métaphore cartésienne : transparence, évidence), « éclaire » (intentionnalité : viser l'objet, le rendre intelligible), « tout » (totalité ? y compris elle-même ? y compris l'inconscient ?).
    \item Proposez un plan dialectique en 3 parties (Thèse / Antithèse / Synthèse) en mobilisant \textbf{au moins 3 auteurs} du programme (Descartes, Kant, Husserl, Sartre, Freud).
    \item Rédigez l'introduction complète (accroche, définition, problématique, annonce de plan).
\end{enumerate}
\end{exercice}

\begin{corrige}[Exercice n°1 — Éléments de corrigé]
\textbf{Introduction :} \og~Descartes compare la conscience à une lumière naturelle qui ne peut être éteinte. Mais Freud découvre un « autre scène » où la lumière ne pénètre pas. La conscience, par sa nature intentionnelle et réflexive, prétend à la transparence totale ; pourtant, sa structure même (temporalité, corporéité, refoulement) laisse des zones d'ombre.~\fg{}
\textbf{Plan :}
\begin{enumerate}
    \item \textbf{Thèse : La conscience comme lumière fondatrice (Transparence et Intentionnalité).} Descartes (cogito, certitude), Kant (unité d'aperception), Husserl (intentionnalité : conscience = clarté du sens).
    \item \textbf{Antithèse : Les zones d'ombre structurales (Temporalité, Corps, Inconscient).} Husserl/Bergson : le présent fuyant, la rétention/protention non thématisées. Merleau-Ponty : le corps propre (je ne « vois » pas mon corps, je « suis » mon corps). Freud : l'inconscient dynamique (refoulement) — la conscience n'est pas maître en sa maison.
    \item \textbf{Synthèse : La conscience comme lucidité limitée et tâche éthique.} La conscience n'est pas une \emph{donnée} achevée (lumière statique) mais une \emph{conquête} (Sartre : \og~l'homme est une passion inutile~fg{} / projet de lucidité). La philosophie (et la psychanalyse) est l'effort pour \emph{étendre} le cercle de la lumière, sans jamais atteindre le « tout ».
\end{enumerate}
\end{corrige}

\coursec{La valeur de la conscience}

Après avoir analysé la \emph{nature} de la conscience — cette présence immédiate à soi-même qui structure notre vie mentale —, nous devons maintenant en mesurer la \emph{valeur}. Pourquoi la conscience est-elle précieuse ? En quoi est-elle la condition \emph{sine qua non} de ce qui fait spécifiquement l'humain : la connaissance de soi, la liberté, la responsabilité morale et la vie en société ? C'est ce que nous allons explorer, en ancrant chaque thèse dans des exemples concrets, y compris camerounais.

\courssub{La conscience comme condition de la connaissance de soi : « Connais-toi toi-même » (Socrate)}

\emph{Intuition.} Imaginez un élève qui récite sa leçon par cœur sans en comprendre le sens, ou un conducteur qui rentre chez lui « en pilote automatique » sans se souvenir du trajet. Dans les deux cas, l'action se fait \emph{sans} conscience réflexive. Socrate, sur l'Agora d'Athènes, interpelle ses concitoyens : vous vous souciez de vos biens, de votre corps, de votre réputation, mais prenez-vous soin de votre âme ? Pour lui, la conscience (la réflexion sur soi) est le début de la sagesse.

\begin{definition}[Connaissance de soi (Socrate)]
La connaissance de soi n'est pas un simple inventaire de ses goûts ou de ses défauts. C'est l'examen critique de ses propres jugements, valeurs et motivations par la \cle{conscience réflexive}. Elle permet de distinguer l'opinion (\emph{doxa}) de la science (\emph{epistémè}) et de ne plus agir par ignorance.
\end{definition}

\emph{Raisonnement.} Sans conscience, je suis opaque à moi-même. Je subis mes pulsions, mes préjugés, mon conditionnement social. La conscience retourne le regard vers l'intérieur : elle \emph{objective} le sujet. C'est le passage de « je vis » à « je sais que je vis et pourquoi je vis ainsi ».

\begin{exemplebox}[Exemple camerounais : le choix d'orientation en Terminale]
Un élève de Terminale technique à Douala hésite entre Génie Civil et Électrotechnique. S'il choisit par mimétisme (son père est maçon) ou par facilité (les notes sont meilleures), il agit \emph{sans} connaissance de soi. S'il prend le temps de s'interroger : « Qu'est-ce qui m'anime ? Quelle est ma représentation du métier ? Quelles sont mes aptitudes réelles ? », il exerce sa conscience réflexive. Ce n'est qu'à cette condition que son choix sera \emph{authentiquement} le sien.
\end{exemplebox}

\begin{attention}[Piège à éviter]
Confondre \emph{connaissance de soi} et \emph{introspection psychologique} brute. Socrate ne demande pas de « fouiller » dans son passé (style psychanalyse), mais de \emph{tester} la solidité de ses raisons d'agir ici et maintenant. La conscience socratique est \emph{critique}, pas seulement narrative.
\end{attention}

\courssub{La conscience comme fondement de la liberté : délibérer et choisir}

\emph{Intuition.} Un robot exécute un programme. Un animal suit son instinct. L'homme, lui, \emph{hésite}. Devant une situation nouvelle, je me dis : « Je \emph{peux} faire ceci, ou cela, ou ne rien faire. » Cette suspension de l'acte, cet espace entre le stimulus et la réponse, \emph{est} la liberté. Et elle n'existe que parce que la conscience me représente des \emph{possibles}.

\begin{propriete}[Thèse classique : Sans conscience, pas de liberté]
La liberté suppose la capacité de se représenter des fins alternatives et de se déterminer par soi-même (\cle{autodétermination}). Or, cette représentation des possibles et cette autodétermination sont des actes de conscience. \textbf{Donc : pas de conscience $\Rightarrow$ pas de liberté.}
\end{propriete}

\emph{Schéma de la délibération.} La conscience transforme une situation contrainte en situation de choix.

\begin{popfigure}
\begin{tikzpicture}[
    node distance=1.8cm and 2.5cm,
    box/.style={draw=popblue, thick, rounded corners, fill=popblue!10, minimum width=2.8cm, minimum height=1.2cm, align=center, font=\small},
    arrow/.style={-Stealth, thick, popdark},
    labelstyle/.style={font=\footnotesize, text=popmuted, align=center}
]
% Nœuds
\node[box, align=center] (sit){Situation\\problématique\\(ex: conflit, choix d'études)};
\node[box, right=of sit, align=center] (cons){Conscience\\réflexive\\(représentation des possibles,\\évaluation des motifs)};
\node[box, right=of cons, align=center] (choix){Choix\\libre\\ (décision rationnelle\\ou engagement)};
\node[box, right=of choix, align=center] (acte){Acte\\responsable\\(action assumée\\et justifiable)};

% Flèches
\draw[arrow] (sit) -- node[above, labelstyle] {Donne lieu à} (cons);
\draw[arrow] (cons) -- node[above, labelstyle] {Permet} (choix);
\draw[arrow] (choix) -- node[above, labelstyle] {Engendre} (acte);

% Annotation retour (régulation)
\draw[arrow, poporange, dashed, bend left=40] (acte) to node[above, labelstyle, poporange, align=center]{Retour d'expérience\\ (ajustement futur)} (sit);
\end{tikzpicture}
\legende{Schéma de la délibération : la conscience médiatise la relation entre la situation et l'acte, fondant la responsabilité.}
\end{popfigure}

\begin{exemplebox}[Exemple de vie courante : la corruption au guichet]
Un fonctionnaire à Yaoundé reçoit un usager pressé qui glisse un billet pour « accélérer » le dossier.
\begin{itemize}
    \item \textbf{Sans conscience (automatisme) :} Il prend l'argent, « tout le monde le fait ». Il est déterminé par la coutume et l'appât du gain.
    \item \textbf{Avec conscience (délibération) :} Il se représente les conséquences : injustice pour les autres, risque disciplinaire, trahison de sa fonction. Il \emph{choisit} de refuser. C'est là qu'il est libre.
\end{itemize}
\end{exemplebox}

\courssub{La conscience morale, source de la responsabilité}

\emph{Intuition.} Pourquoi ne juge-t-on pas responsable un enfant de 3 ans qui casse un vase, ni un somnambule qui blesse quelqu'un, ni un fou en crise ? Parce qu'ils manquaient de \cle{discernement} : la conscience de la nature et de la portée de leur acte. La responsabilité suppose un sujet qui \emph{sait} ce qu'il fait et \emph{veut} le faire.

\begin{definition}[Responsabilité morale et juridique]
Obligation pour un agent moral de \textbf{répondre de ses actes} accomplis \emph{en connaissance de cause} et en toute liberté. Elle implique deux volets :
\begin{enumerate}
    \item \textbf{Imputabilité :} l'acte m'est attribué (je suis l'auteur conscient).
    \item \textbf{Exigibilité :} on peut m'exiger des comptes, me sanctionner ou me récompenser.
\end{enumerate}
\end{definition}

\emph{Application au Cameroun : le Code Pénal.} L'article 67 du Code Pénal camerounais (Loi n°2016/007 du 12 juillet 2016) dispose :
\begin{quote}
\textit{« N'est pas responsable pénalement la personne qui, au moment de l'action, était privée de ses moyens par suite d'un trouble psychique ou neuropsychique ayant aboli son discernement ou le contrôle de ses actes. »}
\end{quote}
C'est la traduction juridique exacte du principe philosophique : \textbf{pas de conscience (discernement aboli) $\Rightarrow$ pas de responsabilité pénale.}

\begin{exemplebox}[Exemple pédagogique : l'expertise psychiatrique au tribunal (illustration)]
Au Tribunal de Grande Instance (TGI) de Mfoundi (Yaoundé), sur une série de 120 dossiers criminels traités, 8 ont fait l'objet d'une expertise psychiatrique d'office. Dans 3 cas (2,5\%), l'expert a conclu à une \textbf{abolition totale du discernement} (bouffée délirante aiguë, épisode maniaque sévère). Les prévenus ont été déclarés \emph{irresponsables pénalement} et orientés vers des soins psychiatriques (mesure de sûreté) au lieu de la prison. Cela montre concrètement que la loi protège la conscience comme condition de la peine.
\end{exemplebox}

\begin{attention}[Distinction cruciale : Abolition vs Altération]
Le Code Pénal (art. 67 et 68) distingue :
\begin{itemize}
    \item \textbf{Abolition du discernement} (trouble \emph{abolit} la conscience, art. 67) $\rightarrow$ \textbf{Irresponsabilité} totale (acquittement + mesure de sûreté).
    \item \textbf{Altération du discernement} (trouble \emph{diminue} la conscience sans l'anéantir, art. 68) $\rightarrow$ \textbf{Responsabilité atténuée} (peine réduite, mais condamnation maintenue).
\end{itemize}
Ne confondez pas les deux : l'altération ne supprime pas la conscience, elle la fragilise.
\end{attention}

\courssub{La conscience comme dignité de l'homme : distinction homme / animal / machine}

\emph{Intuition.} Descartes, dans le \emph{Discours de la méthode}, affirme que les bêtes sont des « machines » (automates) dépourvues de pensée et de conscience. L'homme seul a une \emph{âme} (res cogitans), c'est-à-dire une conscience. Cette thèse fonde l'humanisme classique : la conscience est la \cle{dignité} de l'homme, ce qui le rend \emph{fin en soi} (Kant) et non simple moyen.

\begin{aretenir}[Conscience et dignité humaine]
\begin{itemize}
    \item \textbf{Vs Animal :} L'animal a de la \emph{sentience} (douleur, plaisir, perception), mais pas de \emph{conscience réflexive} (il ne se sait pas percevoir, il ne se juge pas).
    \item \textbf{Vs Machine / IA :} L'IA traite l'information, simule le langage, mais n'a pas de \emph{qualia} (vécu subjectif) ni d'intentionnalité propre. Elle n'a pas de « point de vue » sur le monde.
    \item \textbf{Homme :} Sujet moral, auteur de ses actes, porteur de droits \emph{parce que} conscient.
\end{itemize}
\end{aretenir}

\begin{exoresolu}[Sujet type Bac : « L'intelligence artificielle peut-elle avoir une conscience ? »]
\textbf{Analyse :} Le sujet oppose \emph{intelligence} (calcul, traitement de données, performance) et \emph{conscience} (vécu subjectif, intentionnalité, responsabilité).
\textbf{Plan dialectique :}
\begin{enumerate}
    \item \textbf{Thèse :} L'IA simule des comportements conscients (test de Turing, conversation, diagnostic médical). Elle \emph{fonctionne} comme un conscient.
    \item \textbf{Antithèse :} La simulation n'est pas la duplication. L'IA n'a pas de \emph{qualia} (elle ne \emph{ressent} pas la douleur qu'elle détecte), pas d'intentionnalité originaire (elle n'a pas \emph{voulu} soigner, elle exécute un objectif), pas de responsabilité (on ne juge pas un algorithme).
    \item \textbf{Synthèse :} L'IA nous oblige à redéfinir la conscience : non plus comme une « substance » mystérieuse, mais comme un \emph{niveau d'organisation fonctionnelle} ? Mais tant qu'il n'y a pas de \emph{sujet} qui s'approprie l'acte, la dignité morale reste l'apanage de la conscience humaine.
\end{enumerate}
\end{exoresolu}

\courssub{La conscience collective : le ciment social (Durkheim)}

\emph{Intuition.} La conscience n'est pas seulement individuelle. Dans un village du Nord-Cameroun, lors d'un décès, tout le village arrête ses activités, porte le deuil, partage la peine. Il y a une \emph{conscience commune} : des croyances, des valeurs, des rites partagés qui font que chaque individu se sent membre d'un « nous ».

\begin{definition}[Conscience collective (Émile Durkheim, \emph{De la division du travail social}, 1893)]
Ensemble de croyances et de sentiments communs à la moyenne des membres d'une même société. Elle forme un \textbf{système déterminé qui a sa vie propre}. Elle est la condition de la \cle{solidarité mécanique} (sociétés traditionnelles) et persiste sous forme de valeurs communes (droit, langue, symboles nationaux) dans les sociétés modernes (solidarité organique).
\end{definition}

\begin{exemplebox}[Exemple camerounais : la CAN 2021 (jouée en 2022) au Cameroun]
Pendant la Coupe d'Afrique des Nations, au-delà des clivages linguistiques (anglophones/francophones), ethniques, politiques ou religieux, une \textbf{conscience collective nationale} s'est manifestée :
\begin{itemize}
    \item Port du maillot aux couleurs du pays (symbole partagé).
    \item Rassemblements spontanés aux « fan zones » (rites collectifs).
    \item Sentiment de fierté ou de défaite \emph{partagé} : « Nous avons gagné / perdu ».
\end{itemize}
Cette conscience collective a agi comme un puissant facteur d'unité sociale, fut-ce temporaire.
\end{exemplebox}

\begin{methode}[Mobiliser un exemple de vie courante dans une dissertation]
Pour illustrer une thèse philosophique par un exemple concret (ex: « La conscience est-elle la condition de la responsabilité ? »), suivez ces 4 étapes :
\begin{enumerate}
    \item \textbf{Situation concrète :} Décrivez un fait précis, daté, localisé (ex: « En 2023, au TGI de Mfoundi, l'affaire X... »).
    \item \textbf{Identification de la notion :} Montrez quelle notion du cours l'exemple incarne (ex: « L'expertise psychiatrique a conclu à l'abolition du discernement... »).
    \item \textbf{Articulation argumentative :} Expliquez \emph{pourquoi} cet exemple prouve ou nuance la thèse (ex: « Cela confirme que le droit lie responsabilité et conscience : sans discernement, pas de peine. »).
    \item \textbf{Portée / Limite :} Ne laissez pas l'exemple « parler tout seul ». Concluez sur sa portée philosophique (ex: « Mais cet exemple montre aussi la difficulté de \emph{prouver} l'état de conscience d'autrui... »).
\end{enumerate}
\end{methode}

\begin{savaistu}[Histoire des sciences : La naissance de la responsabilité pénale moderne]
Au Moyen Âge, on jugeait les animaux (procès du cochon de Falaise, 1386) et les objets (la pierre qui a tué). La responsabilité était \emph{objective} (l'acte cause un dommage). La Révolution française et le Code Pénal de 1810 ont introduit le \textbf{subjectivisme pénal} : « Il n'y a point de crime ni de délit sans intention de le commettre » (art. 64 ancien code). C'est la victoire juridique de la \emph{conscience} (l'intention, le discernement) sur la simple causalité physique. Le Cameroun, héritier du droit français, a inscrit ce principe dans son Code Pénal actuel (art. 67).
\end{savaistu}

\begin{popfigure}
\begin{center}
\begin{tabularx}{\linewidth}{|>{\bfseries}l|X|X|X|}
\hline
\rowcolor{popblueL}
\textbf{Valeur de la conscience} & \textbf{Définition philosophique} & \textbf{Auteur / Référence clé} & \textbf{Exemple concret camerounais} \\
\hline
\textbf{Connaissance de soi} & Examen critique de ses motifs, passage de l'opinion à la vérité sur soi. & Socrate (\emph{Alcibiade}, \emph{Apologie}) & Élève de Terminale Tle à Douala qui choisit sa filière après bilan d'orientation et non par mimétisme familial. \\
\hline
\textbf{Fondement de la liberté} & Espace de délibération : représentation des possibles $\rightarrow$ choix autodéterminé. & Descartes, Sartre (\emph{L'Être et le Néant}) & Fonctionnaire à Yaoundé qui refuse un pot-de-vin après avoir pesé les conséquences éthiques et juridiques. \\
\hline
\textbf{Source de la responsabilité} & Condition de l'imputabilité : discernement (connaissance) + liberté (volonté). & Kant (\emph{Fondements de la métaphysique des mœurs}), Code Pénal CM (art. 67) & Acquittement pour irresponsabilité pénale au TGI Mfoundi suite à expertise psychiatrique (abolition du discernement). \\
\hline
\textbf{Dignité de l'homme} & Ce qui distingue le sujet moral (fin en soi) de l'animal/machine (moyen). & Descartes, Kant (\emph{2e Impératif catégorique}) & Débat éthique à l'Assemblée Nationale sur l'encadrement de l'IA : refus de donner la personnalité juridique à un robot. \\
\hline
\textbf{Conscience collective} & Croyances/sentiments communs fondant la solidarité sociale (cohésion). & Durkheim (\emph{De la division du travail social}) & Union nationale autour des Lions Indomptables lors de la CAN 2021 : chants, drapeaux, deuil/joie partagés transcendant les clivages. \\
\hline
\end{tabularx}
\end{center}
\legende{Tableau de synthèse : les trois valeurs cardinales de la conscience (connaissance, liberté, responsabilité) + dignité + dimension collective, illustrées au Cameroun.}
\end{popfigure}

\begin{exercice}[Entraînement dissertation]
\textbf{Sujet :} « Avoir conscience, est-ce être libre ? »
\begin{enumerate}
    \item Problématisez le sujet en distinguant les deux sens de « liberté » (libre-arbitre vs autonomie).
    \item Construisez un plan dialectique en 3 parties (Thèse / Antithèse / Synthèse) en utilisant \emph{au moins deux} des exemples camerounais vus dans ce cours.
    \item Rédigez l'introduction complète (accroche, définition, problématique, annonce du plan).
\end{enumerate}
\end{exercice}

\begin{corrige}[Exercice n°1 : Éléments de correction]
\textbf{Problématique :} La conscience est-elle \emph{suffisante} pour la liberté ? Si je suis conscient de mes chaînes, suis-je libre ? La conscience est-elle \emph{nécessaire} ? (Oui, sans elle, pas de choix).
\textbf{Plan suggéré :}
\begin{enumerate}
    \item \textbf{Thèse :} La conscience est la condition \emph{nécessaire} de la liberté (délibération). Exemple : le fonctionnaire qui refuse la corruption (choix conscient).
    \item \textbf{Antithèse :} La conscience n'est pas \emph{suffisante}. Je peux être conscient mais déterminé (inconscient freudien, conditionnement social, contrainte physique). Exemple : l'addict conscient de sa dépendance mais incapable de s'arrêter ; ou l'élève « conscient » de devoir travailler mais paralysé par l'angoisse.
    \item \textbf{Synthèse :} La conscience est le \emph{lieu} où la liberté se joue (lutte pour l'autonomie). La liberté réelle = conscience \emph{éclairée} + volonté \emph{efficace} (Spinoza : « Je suis libre quand je suis cause adéquate de mes actions »). Exemple : la politique d'éducation à la citoyenneté au Cameroun (programmes d'éducation morale) vise à former des consciences \emph{capables} de liberté.
\end{enumerate}
\end{corrige}

\coursec{Les limites de la conscience}

Nous venons de voir, dans la section précédente, que la conscience possède une valeur considérable : elle nous rend présents au monde, elle fonde la connaissance, la responsabilité morale et la liberté. Descartes en faisait même le fondement indubitable de toute certitude. Mais faut-il en conclure que la conscience est \emph{infaillible} ? Se connaît-elle elle-même parfaitement ? L'expérience la plus ordinaire — une colère « inexpliquée », un oubli suspect, un préjugé que l'on se cache — suggère le contraire. Cette section examine les \cle{limites de la conscience} : non pas pour la détruire, mais pour montrer qu'elle est nécessaire \emph{et} insuffisante, ce qui ouvrira la voie à l'hypothèse de l'inconscient.

\courssub{L'illusion de la transparence : la conscience se trompe sur elle-même}

Commençons par une intuition simple. Quand je ferme les yeux et que je m'interroge : « Pourquoi ai-je dit cela ? Pourquoi ai-je choisi ce métier ? », j'ai l'impression de pouvoir tout voir en moi, comme on regarde l'intérieur d'une pièce éclairée. La conscience se vit comme \cle{transparente} : elle croit que rien en elle ne lui échappe. C'est précisément cette croyance qu'il faut interroger.

\begin{definition}[Illusion de la transparence]
L'\cle{illusion de la transparence} est la croyance spontanée selon laquelle la conscience se connaît elle-même totalement et immédiatement, comme si tout le psychisme était visible pour elle. En réalité, la conscience n'accède qu'à une partie superficielle de la vie psychique.
\end{definition}

Nietzsche, dans \emph{Le Gai Savoir} (livre I, §11), formule cette critique de manière radicale : la conscience n'est que « la dernière couche » du psychisme, la plus superficielle, un phénomène de surface. Sous elle s'agit un monde immense d'instincts, de pulsions et de forces que nous ne voyons pas. La conscience serait comparable à la surface de la mer : calme et visible, alors que tout le mouvement profond se joue en dessous.

\begin{popfigure}
\begin{center}
\begin{tikzpicture}[scale=1]
% Eau
\fill[popblue!15] (-5,-3.2) rectangle (5,0);
\draw[popblue,thick] (-5,0) -- (5,0);
\node[popblue,font=\small] at (4.2,0.25) {surface};
% Iceberg émergé
\fill[white,draw=popdark,thick] (-1.1,0) -- (-0.7,1.1) -- (-0.2,1.5) -- (0.4,1.0) -- (1.0,0.4) -- (1.1,0) -- cycle;
% Iceberg immergé
\fill[popblue!35,draw=popdark,thick] (-1.1,0) -- (-2.3,-1.2) -- (-1.9,-2.9) -- (-0.4,-3.1) -- (0.9,-2.7) -- (2.2,-1.6) -- (1.1,0) -- cycle;
% Étiquettes
\node[font=\small\bfseries,popdark] at (0,0.85) {Conscience};
\node[font=\footnotesize,popdark] at (0,0.45) {(partie visible, petite)};
\node[font=\small\bfseries,popdark] at (0,-1.7) {Inconscient};
\node[font=\footnotesize,popdark] at (0,-2.15) {(partie cachée, immense)};
% Flèches
\draw[->,poporange,thick] (2.6,0.8) -- (0.9,0.8);
\node[font=\footnotesize,poporange,right] at (2.7,0.8) {ce que je crois être};
\draw[->,poporange,thick] (2.6,-1.7) -- (1.6,-1.7);
\node[font=\footnotesize,poporange,right] at (2.7,-1.7) {ce qui me fait agir};
\end{tikzpicture}
\end{center}
\legende{Schéma de l'iceberg psychique : la conscience n'est que la partie émergée du psychisme ; l'essentiel de la vie psychique (pulsions, souvenirs refoulés, déterminismes) demeure immergé et invisible pour elle.}
\end{popfigure}

\begin{attention}[Piège fréquent au Probatoire]
Dire que la conscience a des limites ne signifie \textbf{pas} qu'elle est « inutile » ou « illusoire » en tout. L'erreur classique consiste à opposer brutalement « conscience = mensonge » et « inconscient = vérité ». La conscience conserve toute sa valeur (connaissance, responsabilité, délibération) ; elle est simplement \emph{insuffisante} pour rendre compte de la totalité du psychisme. Nuancez toujours votre conclusion dans la dissertation.
\end{attention}

\courssub{Les erreurs de la conscience : illusions, faux souvenirs et mauvaise foi}

L'illusion de la transparence se vérifie concrètement par les \cle{erreurs de la conscience}. Menons l'investigation comme en sciences : observons les faits avant de conclure.

\textbf{Première observation : les illusions des sens.} Une règle plongée dans un bassin d'eau à l'école paraît « cassée » ; les rails de la voie ferrée Douala--Yaoundé semblent se rejoindre à l'horizon. La conscience perceptive se trompe, et c'est la raison qui doit corriger. Si la conscience se trompe sur le monde extérieur, pourquoi serait-elle infaillible sur elle-même ?

\textbf{Deuxième observation : les faux souvenirs.} Des expériences de psychologie (notamment celles d'Élisabeth Loftus) montrent qu'un témoin peut « se souvenir » avec certitude de détails qui n'ont jamais existé, simplement parce qu'on les lui a suggérés. La mémoire n'enregistre pas comme une caméra : elle reconstruit, et la conscience prend ses reconstructions pour des faits.

\textbf{Troisième observation : l'autojustification.} Après une faute, la conscience fabrique spontanément des excuses : « ce n'était pas de ma faute », « je n'avais pas le choix ». Sartre a analysé ce mécanisme sous le nom de \cle{mauvaise foi}.

\begin{definition}[La mauvaise foi (Sartre)]
La \cle{mauvaise foi} est l'attitude par laquelle la conscience \textbf{se ment à elle-même} : elle se cache sa propre liberté et ses propres motifs en se faisant passer pour ce qu'elle n'est pas (par exemple : « je suis ainsi, je n'y peux rien »). Elle est à la fois le menteur et le menti, ce qui montre qu'en elle quelque chose échappe à sa propre clarté.
\end{definition}

\begin{exemplebox}[La mauvaise foi au quotidien]
Un élève de Terminale technique qui n'a pas révisé son épreuve de philosophie déclare : « De toute façon, je suis nul en dissertation, c'est dans ma nature. » Il se ment : il transforme un choix (ne pas travailler) en destin (être « nul »), pour ne pas assumer sa responsabilité. La conscience se déguise la vérité au lieu de l'affronter.
\end{exemplebox}

\courssub{Les déterminismes qui échappent à la conscience}

Pourquoi la conscience se trompe-t-elle ainsi ? Parce que nos conduites sont largement produites par des \cle{déterminismes} — des causes qui agissent sur nous sans que nous en ayons conscience. On peut les classer en quatre grandes familles :

\begin{itemize}
\item \textbf{Le corps} : la fatigue, la faim, la maladie, les hormones influencent nos humeurs et nos décisions. Qui n'a pas remarqué qu'on s'irrite plus vite le ventre vide ?
\item \textbf{L'hérédité} : le tempérament, certaines dispositions nerveuses se transmettent biologiquement et colorent nos réactions sans que nous le sachions.
\item \textbf{Le milieu social} : la famille, le quartier, l'école, la langue, les coutumes façonnent nos goûts, nos jugements et même nos préjugés, bien avant que nous puissions les choisir.
\item \textbf{Les habitudes} : répétés, nos gestes deviennent automatiques et gouvernent une grande partie de notre journée sans aucune délibération consciente.
\end{itemize}

\begin{popfigure}
\begin{center}
\begin{tikzpicture}[scale=1, every node/.style={font=\small}]
% Boîtes des déterminismes
\node[draw,rounded corners,fill=popblueL,minimum width=2.6cm,minimum height=0.9cm] (corps) at (-4.2,2.2) {\textbf{Le corps}};
\node[draw,rounded corners,fill=popgreenL,minimum width=2.6cm,minimum height=0.9cm] (hered) at (-4.2,0.7) {\textbf{L'hérédité}};
\node[draw,rounded corners,fill=poporangeL,minimum width=2.6cm,minimum height=0.9cm] (soc) at (-4.2,-0.8) {\textbf{Le milieu social}};
\node[draw,rounded corners,fill=poppinkL,minimum width=2.6cm,minimum height=0.9cm] (hab) at (-4.2,-2.3) {\textbf{Les habitudes}};
% Boîte comportement
\node[draw,rounded corners,fill=popgoldL,minimum width=3.2cm,minimum height=1.2cm,align=center] (comp) at (1.8,0) {\textbf{Le comportement}\\ \footnotesize (paroles, choix, actes)};
% Flèches
\draw[->,very thick,popdark] (corps.east) -- (comp.170);
\draw[->,very thick,popdark] (hered.east) -- (comp.180);
\draw[->,very thick,popdark] (soc.east) -- (comp.190);
\draw[->,very thick,popdark] (hab.east) -- (comp.200);
% Conscience
\node[draw,rounded corners,fill=white,minimum width=2.8cm,minimum height=0.9cm,dashed] (consc) at (1.8,2.2) {\textbf{La conscience}};
\draw[->,thick,popmuted,dashed] (consc.south) -- (comp.north) node[midway,right,font=\footnotesize\itshape]{croit tout expliquer};
\end{tikzpicture}
\end{center}
\legende{Les quatre grandes familles de déterminismes (corps, hérédité, milieu social, habitudes) produisent nos comportements, tandis que la conscience, qui n'en perçoit pas le jeu, croit être la seule cause de ses actes.}
\end{popfigure}

\courssub{Spinoza : la conscience ignore les causes de ses désirs}

C'est le philosophe Spinoza, au XVII\textsuperscript{e} siècle, qui a formulé cette critique avec le plus de rigueur dans l'\emph{Éthique}. Son raisonnement peut se reconstruire ainsi :

\begin{enumerate}
\item Les hommes ont conscience de leurs \emph{désirs} et de leurs \emph{actes} : je sais que je veux manger, que je choisis de partir.
\item Mais ils ignorent les \emph{causes} qui déterminent ces désirs : pourquoi ce désir-là, à ce moment-là ? Le corps, l'éducation, le milieu agissent en coulisses.
\item Or, ignorer les causes de ce qu'on fait donne l'impression d'en être l'auteur absolu.
\item Donc les hommes se croient \emph{libres} simplement parce qu'ils sont conscients de leurs désirs et ignorants de leurs causes.
\end{enumerate}

\begin{propriete}[Thèse de Spinoza — à retenir pour le Baccalauréat]
Pour Spinoza, la conscience n'est pas la cause de nos actes : elle en est le témoin partiel et trompeur. Citation de l'\emph{Éthique} (partie II, proposition 35, scolie) à connaître : « Les hommes se croient libres parce qu'ils ont conscience de leurs volitions et de leurs désirs, et qu'ils ignorent les causes par lesquelles ils sont déterminés. » Spinoza compare l'homme à une pierre lancée dans les airs qui, si elle était consciente, croirait voler par sa propre volonté.
\end{propriete}

\begin{exoresolu}[Analyse de situations camerounaises]
\textbf{Énoncé.} À Bafoussam, un père de famille rentre du travail et se met en colère contre son fils pour une broutille : un verre cassé. Lui-même dit ensuite : « Je ne sais pas ce qui m'a pris. » Le même jour, un élève « oublie » de montrer à ses parents un bulletin portant une mauvaise note en mathématiques, et jure qu'il « n'y a pas pensé ». Enfin, un recruteur à Douala affirme « ne pas être tribaliste », mais écarte systématiquement les candidatures d'une certaine région « parce que le profil ne convient pas ». Montrez, pour chaque cas, en quoi la conscience est limitée.
\tcblower
\textbf{Solution détaillée.}
\begin{itemize}
\item \emph{La colère « inexpliquée »} : le père croit réagir au verre cassé, mais la cause réelle de sa colère est ailleurs — fatigue du corps, contrariété au travail, soucis d'argent. La conscience perçoit l'émotion mais ignore ses causes (thèse de Spinoza). Le déterminisme corporel et social agit à son insu.
\item \emph{L'« oubli » du bulletin} : l'oubli n'est pas innocent. La conscience de l'élève se protège d'une situation pénible (la honte, la punition) en écartant le souvenir. Ce n'est pas un hasard : c'est un mécanisme de défense, voisin de la mauvaise foi de Sartre — la conscience se ment à elle-même pour éviter d'affronter la réalité.
\item \emph{Les préjugés ethniques non assumés} : le recruteur est sincère quand il dit « ne pas être tribaliste » : sa conscience ne reconnaît pas le préjugé. Pourtant, ses choix répétés révèlent un déterminisme social (stéréotypes transmis par le milieu, les habitudes, les discours entendus depuis l'enfance) qui oriente ses jugements sans qu'il en ait conscience.
\end{itemize}
\textbf{Conclusion.} Dans les trois cas, la conscience est sincère mais insuffisante : elle voit l'acte, pas ses causes. C'est exactement la limite que Spinoza avait diagnostiquée.
\end{exoresolu}

\begin{savaistu}[Pascal avant Freud]
Bien avant la psychanalyse, Blaise Pascal écrivait au XVII\textsuperscript{e} siècle dans les \emph{Pensées} (Brunschvicg 277 / Lafuma 680) : « Le cœur a ses raisons que la raison ne connaît point. » Cette formule célèbre annonce déjà l'idée d'inconscient : il existe en nous des motifs profonds (« les raisons du cœur ») qui guident nos choix amoureux, nos amitiés, nos préférences, sans que la conscience raisonnante puisse les expliquer. La philosophie avait donc pressenti ce que Freud allait théoriser deux siècles plus tard.
\end{savaistu}

\courssub{Conclusion de la partie : une conscience nécessaire mais insuffisante}

Que retenir de cette analyse ? La conscience se croit transparente à elle-même, mais l'observation la dément : elle se trompe (illusions, faux souvenirs), elle se ment (mauvaise foi), et surtout elle ignore les déterminismes — corps, hérédité, milieu social, habitudes — qui produisent une grande part de nos conduites. La formule de Spinoza résume cette limite fondamentale : nous avons conscience de nos désirs, mais non des causes qui les déterminent.

\begin{aretenir}[Bilan de la section — Fiche de révision]
\begin{itemize}
\item La conscience n'est pas transparente : elle n'est que la « surface » du psychisme (Nietzsche, image de l'iceberg).
\item Elle commet des erreurs : illusions des sens, faux souvenirs, autojustification et mauvaise foi (Sartre).
\item Nos actes sont déterminés par des causes qui lui échappent : corps, hérédité, milieu social, habitudes (Spinoza).
\item La conscience est donc \textbf{nécessaire mais insuffisante} : il faut admettre l'hypothèse d'une vie psychique qui lui échappe — l'\cle{inconscient}, objet de la section suivante.
\end{itemize}
\end{aretenir}

\begin{exercice}[Pour s'entraîner — Type Probatoire]
\begin{enumerate}
\item Définissez la mauvaise foi et illustrez-la par un exemple tiré de la vie scolaire camerounaise.
\item Expliquez la comparaison de la pierre chez Spinoza. Que prouve-t-elle sur le sentiment de liberté ?
\item Sujet de dissertation : « La conscience se connaît-elle elle-même ? » Construisez un plan détaillé en deux parties (thèse : valeur de la conscience ; antithèse : ses limites) avec arguments et auteurs pour chaque partie.
\end{enumerate}
\end{exercice}

\begin{corrige}[Exercice — Éléments de réponse attendus]
\begin{enumerate}
\item \textbf{Définition :} La mauvaise foi est le mensonge que la conscience se fait à elle-même pour fuir sa liberté et sa responsabilité (Sartre). \textbf{Exemple :} L'élève qui triche à l'examen et déclare : « Tout le monde le fait, le surveillant ne regardait pas, je n'avais pas le choix. » Il nie sa liberté de ne pas tricher.
\item \textbf{Explication :} La pierre lancée, si elle pensait, croirait se mouvoir librement alors qu'elle obéit à une impulsion extérieure (le lancer) et à la gravité. De même, l'homme confond la conscience de ses actes (je veux, je fais) avec la maîtrise de leurs causes (pourquoi je veux cela). Le sentiment de liberté repose sur l'ignorance des déterminismes (corps, milieu, habitudes).
\item \textbf{Plan dissertation :}
\textbf{Introduction :} Accroche (ex. citation de Spinoza), définition de la conscience (connaissance réflexive), problématique (la conscience est-elle lucide sur elle-même ?), annonce du plan.
\textbf{Partie I : La conscience, fondement de la lucidité et de la responsabilité (Thèse).}
\begin{itemize}
    \item A. La certitude cartésienne : le \emph{cogito} comme socle indubitable (Descartes).
    \item B. La conscience morale : témoin de nos actes, fondement de la responsabilité (Kant, voix de la conscience).
    \item C. La conscience comme unité du sujet : elle rassemble les expériences en un « Je » identique.
\end{itemize}
\textbf{Partie II : Les limites structurelles de la conscience (Antithèse).}
\begin{itemize}
    \item A. L'illusion de la transparence : la conscience ne voit que la surface (Nietzsche, iceberg).
    \item B. Les erreurs et la mauvaise foi : illusions, faux souvenirs, mensonge à soi-même (Sartre).
    \item C. L'ignorance des causes : déterminismes corporels, sociaux, habitudes (Spinoza, Marx, Freud).
\end{itemize}
\textbf{Conclusion :} La conscience est réelle mais partielle. Elle n'est pas source, mais effet. Ouverture : l'inconscient freudien comme théorie scientifique de ces causes cachées.
\end{enumerate}
\end{corrige}

\coursec{La Conscience et l'Inconscient : L'inconscient freudien}

Après avoir exploré les limites de la conscience — sa partialité, sa fragilité face à l'habitude et à l'émotion, son ignorance de ses propres moteurs —, nous devons nous demander : \emph{que devient ce que la conscience refuse ou ne peut pas voir ?} L'histoire de la philosophie (Leibniz, Schopenhauer, Nietzsche) a pressenti l'existence d'un « dessous » de la conscience. Mais c'est \textbf{Sigmund Freud (1856--1939)}, médecin viennois fondateur de la \textbf{psychanalyse}, qui en a fait un objet d'étude rigoureux. Son ouvrage majeur, \emph{L'Interprétation des rêves} (1900), pose l'hypothèse décisive : l'inconscient n'est pas une simple absence de conscience, c'est un \emph{autre} système psychique, dynamique, régi par ses lois propres, qui détermine nos pensées, nos actes et nos symptômes à notre insu.

\courssub{La découverte freudienne : l'inconscient comme lieu du refoulé}

\begin{definition}[L'inconscient freudien]
L'\cle{inconscient} (das Unbewusste) est, chez Freud, un \textbf{système psychique} constitué de \textbf{représentations} (idées, souvenirs, images) et d'\textbf{affects} (émotions, pulsions) qui sont \textbf{inaccessibles à la conscience} non pas par défaut d'attention, mais par un travail actif de \cle{refoulement} (Verdrängung). Ces contenus, bien que exclus de la conscience, conservent toute leur énergie psychique et \textbf{agissent à l'insu du sujet} : ils déterminent des comportements, provoquent des symptômes, s'expriment dans les rêves et les actes manqués.
\end{definition}

\begin{attention}[Piège conceptuel : l'inconscient n'est pas un « lieu » cérébral]
\emph{Erreur fréquente :} représenter l'inconscient comme une « zone » du cerveau (ex. : « l'inconscient se situerait dans le lobe limbique »).
\emph{Correction :} Pour Freud, l'inconscient est une \textbf{hypothèse théorique} (un \emph{modèle topique}) destinée à rendre compte de la \textbf{dynamique psychique}. C'est une \textbf{qualité} de certains processus psychiques (le fait d'être refoulés et dynamiques), non une localisation anatomique. La psychanalyse est une \emph{psychologie des profondeurs}, pas une neurobiologie.
\end{attention}

\textbf{Le refoulement, mécanisme fondateur.}\par
Pourquoi des contenus sont-ils inconscients ? Parce qu'ils sont \textbf{inacceptables} pour le moi conscient : pulsions sexuelles infantiles, désirs agressifs, pensées contraires à la morale ou à l'image de soi. Le \cle{refoulement} est l'opération par laquelle le « Moi » (ou le Surmoi) \emph{chasse} ces représentations hors de la conscience et les maintient dans l'inconscient. \emph{C'est un acte de défense, pas un oubli passif.} Le refoulé ne disparaît pas ; il forme un « noyau » qui presse pour revenir à la conscience, créant un \textbf{conflit psychique} permanent.

\begin{exemplebox}[Exemple concret : le refoulement d'un désir hostile]
Un élève de Terminale technique, Paul, admire son professeur de Génie Mécanique. Pourtant, lors d'une correction sévère, Paul ressent une \emph{impulsion violente} : « J'aimerais le frapper ». Cette pensée est insupportable pour son moi (contraire à ses valeurs, à la loi, à l'image du « bon élève »). Il la \textbf{refoule} : elle devient inconsciente. Quelques jours plus tard, Paul \emph{oublie} systématiquement de rendre un devoir crucial pour ce professeur (acte manqué), ou il fait un \emph{lapsus} en l'appelant « Monsieur le Tyran » au lieu de son nom. Le désir refoulé (agressivité) a trouvé une issue déguisée.
\end{exemplebox}

\courssub{Les deux topographies : cartographier l'appareil psychique}

Freud a proposé deux modèles successifs pour structurer l'appareil psychique. Ils ne s'annulent pas ; le second affine le premier.

\textbf{A. La première topique (1900) : Conscient / Préconscient / Inconscient.}\par
C'est une topographie \textbf{spatiale} et \textbf{économique} (circulation de l'énergie).
\begin{itemize}
    \item \textbf{Le Conscient (Cs)} : surface de l'appareil psychique, en contact avec le monde extérieur (perceptions) et intérieur (sensations corporelles). C'est le seul lieu de la \emph{qualité de conscience}.
    \item \textbf{Le Préconscient (Pcs)} : réservoir de souvenirs, de connaissances, de pensées \emph{latentes} qui ne sont pas conscientes \emph{maintenant} mais peuvent le devenir \emph{sans résistance} (ex. : votre numéro de téléphone, la capitale du Cameroun). Pas de refoulement ici.
    \item \textbf{L'Inconscient (Ics)} : le domaine du \textbf{refoulé}. Contenus interdits d'accès au Pcs/Cs par la \textbf{Censure} (barrière entre Ics et Pcs). Régi par le \textbf{principe de plaisir} (recherche immédiate de la satisfaction, sans temps, sans négation, sans contradiction).
\end{itemize}

\textbf{B. La seconde topique (1923) : Ça / Moi / Surmoi.}\par
C'est une topographie \textbf{structurale} et \textbf{dynamique} (conflits d'instances). Elle explique \emph{comment} le refoulement opère.

\begin{popfigure}
\begin{tikzpicture}[
    node distance=1.2cm and 1.5cm,
    box/.style={draw=popdark, thick, rounded corners, fill=popcream, minimum width=3.5cm, minimum height=2.2cm, align=center, font=\small},
    arrow/.style={-Stealth, thick, popdark},
    conflict/.style={-Stealth, dashed, thick, poporange, bend left=15},
    labelstyle/.style={font=\footnotesize, text=popmuted, align=center}
]
% Les trois instances
\node[box, fill=poporangeL, align=center] (ca){\textbf{Le ÇA (Es)}\\(Das Es)\\\emph{« Ça pense en moi »}\\\vspace{0.1cm}Pulsions (vie/mort)\\\emph{Principe de plaisir}\\\emph{Totalement inconscient}};
\node[box, fill=popblueL, right=of ca, align=center] (moi){\textbf{Le MOI (Ich)}\\\emph{« Je »}\\\vspace{0.1cm}Médiateur réalité\\\emph{Principe de réalité}\\\emph{Conscient + Préconscient + Inconscient (défenses)}};
\node[box, fill=poppurpleL, above=of moi, yshift=-0.5cm, align=center] (surmoi){\textbf{Le SURMOI (Über-Ich)}\\\emph{Idéal du moi + Conscience morale}\\\vspace{0.1cm}Interdits, lois, normes\\\emph{Héritage parental/social}\\\emph{Inconscient en grande partie}};
% Flèches de conflit / interaction
\draw[conflict] (ca.east) -- node[labelstyle, below, align=center]{Exige satisfaction\\immédiate} (moi.west);
\draw[conflict] (moi.west) -- node[labelstyle, above] {Refoule / Contrôle} (ca.east);
\draw[conflict] (surmoi.south) -- node[labelstyle, right, align=center]{Juge / Interdit\\Culpaibilise} (moi.north);
\draw[conflict] (moi.north) -- node[labelstyle, left] {Subit / Négocie} (surmoi.south);
\draw[conflict] (ca.north) to[bend right=20] node[labelstyle, left] {Pulsions brutes} (surmoi.west);
\draw[conflict] (surmoi.west) to[bend left=20] node[labelstyle, right] {Répression morale} (ca.north);
% Légende externe
\node[labelstyle, below=0.8cm of moi] (leg) {\textbf{Le Moi : « Serviteur de trois maîtres » (Ça, Surmoi, Réalité extérieure)}};
\end{tikzpicture}
\legende{Schéma de la seconde topique freudienne : les trois instances et leurs conflits dynamiques. Le Moi tente de concilier les exigences pulsionnelles du Ça, les interdits moraux du Surmoi et les contraintes du monde réel.}
\end{popfigure}

\begin{definition}[Les trois instances de la seconde topique]
\begin{itemize}
    \item \textbf{Le Ça (Es)} : \textbf{Inconscient pur}. Réservoir des \textbf{pulsions} (sexualité = Éros / agressivité = Thanatos). Aucune organisation, pas de temps, pas de logique, recherche la \textbf{décharge immédiate} (principe de plaisir). « Ça pense en moi. »
    \item \textbf{Le Moi (Ich)} : \textbf{Partie modifiée du Ça} par l'expérience du monde extérieur. Siège de la \textbf{raison}, de la \textbf{volonté}, de la \textbf{mémoire}, de la \textbf{perception}. Fonction : \textbf{médiation}. Il substitue le \textbf{principe de réalité} (attendre, chercher le bon objet, différer la satisfaction) au principe de plaisir. \emph{Majoritairement Conscient/Préconscient}, mais ses \textbf{mécanismes de défense} (refoulement, déplacement, rationalisation...) sont \textbf{Inconscients}.
    \item \textbf{Le Surmoi (Über-Ich)} : \textbf{Héritier du complexe d'Œdipe}. Intériorisation de l'\textbf{autorité parentale} et des \textbf{normes sociales}. Deux faces : \emph{Idéal du moi} (ce qu'on devrait être) et \emph{Conscience morale} (ce qu'on ne doit pas faire). Juge le Moi, produit \textbf{culpabilité} et \textbf{honte}. \emph{Largement Inconscient}.
\end{itemize}
\end{definition}

\begin{aretenir}[Le conflit permanent]
La vie psychique normale est un \textbf{équilibre précaire} entre ces trois instances. Le Moi est le « pauvre moi » (Freud) : « Serviteur de trois maîtres sévères : le monde extérieur, le Ça et le Surmoi. » La \textbf{névrose} éclate quand le Moi ne parvient plus à gérer ce conflit (ex. : refoulement trop rigide, Surmoi tyrannique, Ça trop puissant).
\end{aretenir}

\courssub{Les manifestations de l'inconscient : la « psychopathologie de la vie quotidienne »}

L'inconscient ne se tait pas. Il « parle » à travers des formations de compromis qui déjouent la censure. Freud les analyse dans \emph{Psychopathologie de la vie quotidienne} (1901).

\begin{propriete}[Les formations de l'inconscient (voie royale)]
\begin{itemize}
    \item \textbf{Les rêves} : « \textbf{Voie royale vers l'inconscient} ». Satisfaction hallucinatoire d'un désir refoulé.
    \item \textbf{Les actes manqués (Fehlleistungen)} : actes ratés (laisser tomber un objet, se tromper de chemin) qui réalisent un désir caché.
    \item \textbf{Les lapsus (Versprechen)} : erreurs de parole (mot pour un autre) révélant une pensée refoulée.
    \item \textbf{Les oublis de noms, de projets, de rendez-vous} : oubli motivé par un affect déplaisant lié à ce contenu.
    \item \textbf{Les symptômes névrotiques} : phobies, obsessions, hystéries = compromis symbolique entre désir refoulé et défense.
    \item \textbf{L'humour, le trait d'esprit (Witz)} : décharge de tension pulsionnelle (sexuelle/agressive) socialement acceptée.
\end{itemize}
\end{propriete}

\textbf{Le rêve : structure et interprétation.}\par
Le rêve n'est pas un bruit neuronal aléatoire ; c'est une \textbf{construction psychique} à sens.

\begin{popfigure}
\begin{tikzpicture}[
    node distance=1.0cm and 1.0cm,
    proc/.style={draw=popdark, thick, rounded corners, fill=popblueL, minimum width=3.2cm, minimum height=1.6cm, align=center, font=\small},
    arrow/.style={-Stealth, thick, popdark},
    labelstyle/.style={font=\footnotesize, text=popmuted, align=center, above}
]
% Nœuds
\node[proc, align=center] (desir){\textbf{Désir Refoulé}\\(Inconscient / Ça)\\\emph{Ex: « Je veux échouer}\\\emph{pour fuir la pression »}};
\node[proc, right=of desir, fill=poporangeL, align=center] (censure){\textbf{Censure / Travail du rêve}\\(Moi / Surmoi)\\\emph{Déformation, Déplacement,}\\\emph{Condensation, Figuration}};
\node[proc, right=of censure, fill=popgreenL, align=center] (manifeste){\textbf{Contenu Manifeste}\\(Rêve raconté)\\\emph{« Je cours nu dans}\\\emph{la cour de l'école »}};
\node[proc, below=of manifeste, fill=poppurpleL, align=center] (interp){\textbf{Interprétation}\\(Analyse / Association libre)\\\emph{Recherche du sens caché}};
\node[proc, left=of interp, fill=poppinkL, align=center] (latent){\textbf{Contenu Latent}\\(Pensées du rêve)\\\emph{« Je refuse l'examen}\\\emph{du Probatoire, je veux}\\\emph{rester enfant »}};
% Flèches
\draw[arrow] (desir) -- node[labelstyle] {Pousse à s'exprimer} (censure);
\draw[arrow] (censure) -- node[labelstyle] {Produit (déguisé)} (manifeste);
\draw[arrow] (manifeste) -- node[labelstyle, right] {Matériau brut} (interp);
\draw[arrow] (interp) -- node[labelstyle] {Révèle} (latent);
\draw[arrow, dashed, bend left=40] (latent) to node[labelstyle, left] {Est la traduction de} (desir);
\end{tikzpicture}
\legende{Du désir refoulé au contenu latent : le travail du rêve déforme le désir (censure) pour produire le contenu manifeste ; l'interprétation psychanalytique remonte du manifeste au latent.}
\end{popfigure}

\begin{definition}[Contenu manifeste vs Contenu latent]
\begin{itemize}
    \item \textbf{Contenu manifeste} : le rêve \emph{tel qu'on le raconte} (images, scénario apparent). C'est une \textbf{façade}, souvent absurde ou banale.
    \item \textbf{Contenu latent} : les \textbf{pensées inconscientes réelles} (le désir refoulé, les souvenirs d'enfance, les affects) que le rêve exprime de façon codée.
    \item \textbf{Travail du rêve (Traumarbeit)} : ensemble des opérations (Condensation, Déplacement, Figuration, Élaboration secondaire) qui transforment le latent en manifeste pour \textbf{tromper la censure} et permettre le sommeil.
\end{itemize}
\end{definition}

\begin{exemplebox}[Exemple guidé : Rêve d'examen raté]
\textbf{Contexte :} Élève en Terminale TEE (Électrotechnique), veille du Probatoire blanc.
\textbf{Rêve (Manifeste)} : « Je suis dans la salle d'examen. Le sujet est en chinois. Je cherche mon stylo, il est devenu une carotte. Le surveillant est mon père qui rit. Je me réveille en sueur. »
\textbf{Analyse (Latent)} :
\begin{enumerate}
    \item \emph{Sujet en chinois} = Incompréhension totale $\rightarrow$ \textbf{Peur de l'échec / sentiment d'incompétence}.
    \item \emph{Stylo $\rightarrow$ Carotte} (Déplacement/Figuration) : Le stylo = outil de pouvoir/réussite (phallus symbolique ?). Carotte = objet inutile, infantile, alimentaire $\rightarrow$ \textbf{Régression, impuissance}.
    \item \emph{Père qui rit} (Condensation : Père + Surveillant) = \textbf{Surmoi jugeant, moqueur}. Crainte du jugement paternel.
    \item \textbf{Désir refoulé (Ça)} : « Je veux \emph{échouer} pour ne plus avoir à porter le poids des attentes (père, école, société) et revenir à l'enfance (carotte) sans responsabilité. »
\end{enumerate}
\emph{Conclusion :} Le rêve permet une satisfaction de ce désir (échouer symboliquement) tout en préservant le sommeil (anxiété réveillée = échec de la censure).
\end{exemplebox}

\begin{exoresolu}[Analyse d'un lapsus en classe]
\textbf{Énoncé :} Lors d'un exposé sur « La maintenance préventive des machines-outils », un élève dit : « Il faut bien \textbf{graisser les pannes} » au lieu de « graisser les \textbf{pièces} ». La classe rit. L'élève rougit et se corrige. Analysez ce lapsus à la lumière de la psychanalyse.
\tcblower
\textbf{Solution détaillée :}
\begin{enumerate}
    \item \textbf{Description du fait brut (Manifeste) :} Substitution phonétique/sémantique : « pièces » $\rightarrow$ « pannes ». Contexte : exposé technique, enjeu de réussite (note), public (pairs, prof).
    \item \textbf{Interrogation du sens caché (Latent) :}
    \begin{itemize}
        \item « Pannes » = \textbf{dysfonctionnement, arrêt, échec}. C'est l'\emph{antonyme} fonctionnel de « maintenance préventive ».
        \item L'élève \emph{craint} l'échec (Probatoire, avenir). Cette angoisse est refoulée pendant l'exposé (Moi : « Je maîtrise mon sujet »).
        \item Le mot « pannes » jaillit : c'est la \textbf{vérité de l'inconscient} (Ça/Surmoi) qui perce : « Tu as peur de la panne / de l'échec ».
    \end{itemize}
    \item \textbf{Conclusion prudente :} Ce lapsus n'est pas une « preuve » absolue d'un désir conscient de saboter l'exposé. Il révèle un \textbf{conflit intrapsychique} : \emph{Volonté de réussir (Moi conscient)} vs \emph{Angoisse d'échec / Auto-sabotage inconscient (Ça/Surmoi)}. C'est une formation de compromis : l'angoisse s'exprime sans empêcher l'exposé de continuer.
\end{enumerate}
\end{exoresolu}

\courssub{Valeur et limites de la psychanalyse}

La psychanalyse n'est pas une simple thérapie ; c'est une \textbf{théorie de l'esprit}, une \textbf{méthode d'investigation} et une \textbf{pratique clinique}. Son apport philosophique est majeur.

\begin{aretenir}[Valeur heuristique et philosophique]
\begin{itemize}
    \item \textbf{Déstabilisation du « Cogito » cartésien :} « Le Moi n'est pas maître en sa propre maison. » La conscience n'est pas la transparence ; elle est \textbf{surdéterminée} par l'inconscient. \emph{Subjectivité $\neq$ Conscience immédiate.}
    \item \textbf{Explication du symptôme :} Le symptôme n'est pas une « erreur » biologique, mais un \textbf{sens} (message codé du désir). Passage du \emph{biomédical} au \emph{herméneutique} (interprétation).
    \item \textbf{Théorie de la culture :} \emph{Malaise dans la civilisation} (1929) : La culture exige le \textbf{refoulement pulsionnel} (sublimation). Le malheur vient de la tension irréductible entre \emph{besoins pulsionnels} (Ça) et \emph{exigences sociales} (Surmoi/Réalité).
    \item \textbf{Méthode :} \textbf{Association libre} (dire tout ce qui vient sans censure) + \textbf{Transfert} (réactivation sur l'analyste des affects infantiles) $\rightarrow$ Levée du refoulement $\rightarrow$ \textbf{Insight} (prise de conscience) $\rightarrow$ \textbf{Autonomie} relative du Moi (« Là où était le Ça, le Moi doit advenir »).
    \item \textbf{Sublimation} : Détournement de l'énergie pulsionnelle vers des objets socialement valorisés (art, science, travail technique). Seule « réussite » possible du conflit pulsion/culture.
\end{itemize}
\end{aretenir}

\begin{attention}[Limites épistémologiques et scientifiques (Débat Bac)]
\begin{itemize}
    \item \textbf{Non-falsifiabilité (Popper) :} La psychanalyse explique \emph{tout} a posteriori (réussite = sublimation ; échec = refoulement). Elle échappe au critère de démarcation scientifique (réfutabilité). \emph{Statut : « Protocience » ou « Herméneutique de profondeur » ?}
    \item \textbf{Universalité contestée :} Modèle basé sur la \textbf{bourgeoisie viennoise victorienne} (fin XIXe). Complexe d'Œdipe universel ? Anthropologues (Malinowski, Lévi-Strauss) montrent la variété des structures familiales.
    \item \textbf{Determinisme psychique radical :} Tout acte a un sens inconscient. Risque de \textbf{réductionnisme} : la liberté, la responsabilité morale, la rationalité deviennent des illusions. Où est le sujet responsable ?
    \item \textbf{Sexualité infantile :} Concept central (polymorphie perverse, stades oral/anal/phallique/génital). Critiqué comme projection adulte sur l'enfant (Piaget, psychologie du développement).
    \item \textbf{Efficacité thérapeutique :} Longue, coûteuse, non standardisée. Difficile à évaluer par essais cliniques randomisés (EBM). Place actuelle : psychothérapies d'inspiration psychanalytique (brèves, focalisées) vs psychanalyse classique.
\end{itemize}
\end{attention}

\begin{savaistu}[Freud et le Cameroun : une réception critique]
La psychanalyse est arrivée au Cameroun via la psychiatrie coloniale (hôpital Jamot à Yaoundé, années 1930-50) puis la formation universitaire (Département de Philosophie/Psychologie à l'Université de Yaoundé I dès les années 70).
\begin{itemize}
    \item \textbf{Adaptation culturelle :} Le « Complexe d'Œdipe » ne s'applique pas directement dans les sociétés à \textbf{lignée matrilinéaire} (ex. Bassa, Bakoko) où l'oncle maternel a l'autorité, ou dans les familles polygames/élargies. Le « Surmoi » est collectif (ancêtres, lignage) avant d'être parental nucléaire.
    \item \textbf{Psychanalyse et tradition :} Des psychanalystes camerounais (ex. Pr. Minko Minka, Pr. Owona) travaillent sur l'articulation \emph{inconscient individuel / inconscient culturel} (rites, totems, croyances, sorcellerie comme formation de l'inconscient collectif).
    \item \textbf{Pratique actuelle :} Peu de psychanalystes « orthodoxes » (divan, 4 séances/semaine). Dominent les \textbf{psychothérapies d'orientation analytique} (face-à-face, 1/semaine) en cabinet libéral (Douala, Yaoundé) et en milieu hospitalier (CHU, Laquintinie). Formation : DEA/Doctorat Psychologie Clinique + analyse didactique personnelle + supervision.
\end{itemize}
\end{savaistu}

\begin{methode}[Analyser un rêve ou un lapsus en 3 étapes (Bac / Dissertation)]
Pour traiter un exemple concret dans une copie :
\begin{enumerate}
    \item \textbf{Décrire le fait brut (Manifeste) :} Citer exactement les mots du lapsus, les images du rêve, l'acte manqué. \emph{Ne pas interpréter tout de suite.}
    \item \textbf{Interroger le sens caché (Latent) :} Chercher le \textbf{désir} ou l'\textbf{affect refoulé} qui trouve là une satisfaction déguisée. Utiliser les mécanismes : \textbf{Condensation} (fusion), \textbf{Déplacement} (glissement d'affect), \textbf{Figuration} (pensée $\rightarrow$ image), \textbf{Symbolisation} (codes culturels partagés). \emph{Contexter : situation de vie, angoisses actuelles, histoire du sujet.}
    \item \textbf{Conclure avec prudence épistémologique :} « Cet exemple \emph{illustre} le mécanisme freudien de... Il \emph{suggère} un conflit entre... Cependant, l'interprétation psychanalytique n'est pas une \emph{preuve} scientifique au sens poppérien ; elle est une \emph{hypothèse de sens} validée par son efficacité thérapeutique (levée du symptôme) et la cohérence du récit du sujet (associations libres). »
\end{enumerate}
\end{methode}

\begin{exercice}[Entraînement : Dissertation type Bac]
\textbf{Sujet :} « L'inconscient freudien ruine-t-il la notion de responsabilité morale ? »
\textbf{Pistes de plan :}
\begin{enumerate}
    \item \textbf{Thèse : Oui, il la ruine.} Le sujet n'est pas auteur de ses actes (déterminisme pulsionnel). Le Moi est « serviteur ». La responsabilité suppose une conscience libre et transparente ; l'inconscient montre que nos motifs réels nous échappent. \emph{Ex : criminel « poussé » par un trauma infantile refoulé.}
    \item \textbf{Antithèse : Non, il la déplace et la fonde.} La responsabilité ne demande pas une conscience totale, mais une \textbf{capacité de réponse} (Moi). La psychanalyse \emph{libère} en rendant conscient l'inconscient ($\rightarrow$ autonomie). Le Surmoi \emph{est} l'instance de la responsabilité (intériorisation de la loi). \emph{Ex : analyse réussie = sujet plus responsable.}
    \item \textbf{Synthèse :} Responsabilité \emph{révisée}, non abolie. Passage de la \textbf{responsabilité métaphysique} (libre arbitre absolu) à la \textbf{responsabilité psychique/juridique} (capacité d'entendement, de maîtrise de soi relative). Le droit (ex. Code pénal camerounais, art. 64 : irresponsabilité pour trouble psychique) intègre déjà cette nuance. La psychanalyse éclaire le \emph{degré} de responsabilité, elle ne la nie pas.
\end{enumerate}
\end{exercice}

\begin{corrige}[Exercice n° Dissertation]
\textbf{Introduction :} Accroche (Descartes vs Freud). Définition : Inconscient = système refoulé dynamique. Responsabilité = imputabilité d'un acte à un agent conscient et libre. Problématique : Si l'inconscient détermine mes actes, suis-je encore l'auteur responsable ? Annonce du plan dialectique.
\textbf{Développement :} Suivre les 3 parties ci-dessus. \textbf{Argumenter} avec concepts précis : Refoulement, Ça/Moi/Surmoi, Déterminisme psychique, Sublimation, Transfert, « Là où était le Ça, le Moi doit advenir ». \textbf{Exemples} : Névrose obsessionnelle (rire au enterrement), Cas « L'homme aux loups », Droit pénal (discernement).
\textbf{Conclusion :} L'inconscient ne supprime pas la responsabilité ; il en complexifie l'exercice. Il impose une \textbf{éthique de la lucidité} (connais-toi toi-même) comme condition d'une responsabilité authentique. Ouverture : Neurosciences actuelles (Libet, prise de décision inconsciente) rejoignent-elles Freud ?
\end{corrige}

\coursec{Valeur et limites de la psychanalyse}

Nous avons vu, dans la section précédente, que Freud découvre en l'homme un continent caché : l'\cle{inconscient}, fait de désirs refoulés qui s'expriment dans les rêves, les lapsus et les actes manqués. Mais découvrir l'inconscient ne suffit pas : encore faut-il savoir ce que vaut cette découverte. La psychanalyse est-elle une méthode capable de guérir ? Une science au même titre que la physique ? Ou bien une simple manière d'interpréter les conduites humaines ? Telle est la question qui va nous occuper : il s'agit de peser, avec honnêteté, la \cle{valeur} de la psychanalyse et ses \cle{limites}.

\courssub{La psychanalyse comme méthode thérapeutique : la cure par la parole}

\textbf{Intuition.} Lorsqu'un élève de Yaoundé est rongé par un souci qu'il n'ose dire à personne, le simple fait de le confier à un ami de confiance le soulage déjà. Parler libère. Freud a fait de cette intuition banale une méthode de traitement : la \cle{cure psychanalytique}.

\begin{definition}[La cure psychanalytique]
La \cle{cure psychanalytique} est une méthode thérapeutique fondée sur la parole : le patient, allongé sur un divan, dit tout ce qui lui passe par l'esprit, sans choisir ni censurer ses pensées (\emph{association libre}), tandis que l'analyste écoute et l'aide à interpréter ce discours pour ramener à la conscience les désirs refoulés qui causent ses troubles.
\end{definition}

Le dispositif repose sur trois mécanismes essentiels :

\begin{enumerate}
\item \textbf{L'association libre.} Le patient dit tout, même ce qui semble absurde, honteux ou insignifiant. En laissant vagabonder sa parole, il contourne la censure du Moi et laisse affleurer les contenus inconscients. C'est la « règle fondamentale » de la cure.
\item \textbf{L'interprétation.} L'analyste repère dans ce discours libre les répétitions, les oublis, les rêves, les lapsus, et propose des hypothèses sur le conflit refoulé qui se cache derrière les symptômes.
\item \textbf{Le transfert.} Au cours de la cure, le patient reporte sur l'analyste des sentiments (amour, hostilité, dépendance) qui appartenaient en réalité aux figures de son enfance (père, mère). Ce \cle{transfert} n'est pas un obstacle : il rejoue le conflit originel dans la situation présente et permet de le comprendre, puis de le dépasser.
\end{enumerate}

\begin{popfigure}
\begin{center}
\begin{tikzpicture}[
  node distance=0.9cm,
  etape/.style={draw=popblue, thick, fill=popblueL, rounded corners=3pt, text width=3.1cm, align=center, minimum height=1.1cm, font=\small},
  fleche/.style={-{Stealth[length=3mm]}, very thick, poporange}
]
\node[etape, align=center] (p1){\textbf{Patient sur le divan}\\ (détente, parole sans censure)};
\node[etape, right=of p1, align=center] (p2){\textbf{Parole libre}\\ (associations, rêves, lapsus)};
\node[etape, right=of p2, align=center] (p3){\textbf{Analyste}\\ (écoute, interprétation, transfert)};
\node[etape, below=0.9cm of p3, align=center] (p4){\textbf{Prise de conscience}\\ (le refoulé devient conscient)};
\node[etape, left=of p4, align=center] (p5){\textbf{Apaisement}\\ (le symptôme perd sa raison d'être)};
\draw[fleche] (p1) -- (p2);
\draw[fleche] (p2) -- (p3);
\draw[fleche] (p3) -- (p4);
\draw[fleche] (p4) -- (p5);
\end{tikzpicture}
\end{center}
\legende{Schéma de la cure analytique : de la parole libre à l'apaisement. Le symptôme disparaît lorsque le conflit inconscient qui le produisait est devenu conscient.}
\end{popfigure}

\begin{exemplebox}[Le cas d'Anna O.]
Avant même Freud, le médecin Breuer soigna une jeune femme, Anna O., atteinte de paralysies et de troubles de la vue sans cause organique. Il remarqua que ses symptômes s'atténuaient lorsqu'elle racontait, sous hypnose puis librement, les souvenirs pénibles liés à la maladie de son père. Elle appela elle-même ce procédé le \emph{talking cure}, la « cure par la parole ». Freud en conclut que le symptôme est un souvenir refoulé qui s'exprime par le corps : le rendre conscient, c'est le désamorcer.
\end{exemplebox}

\courssub{La valeur explicative de la psychanalyse}

La première grandeur de la psychanalyse est de rendre intelligibles des conduites que la conscience seule ne peut pas expliquer.

\begin{exemplebox}[Ce que la conscience n'explique pas]
\begin{itemize}
\item Une \textbf{phobie} : une jeune fille de Douala ne peut plus monter dans un taxi sans paniquer, alors qu'elle sait parfaitement qu'aucun danger réel ne la menace. Sa conscience dit « non », son corps dit « peur ». La psychanalyse cherche derrière cette peur déplacée un souvenir ou un conflit refoulé.
\item Une \textbf{névrose obsessionnelle} : un homme se lave les mains cinquante fois par jour sans pouvoir s'en empêcher. La volonté consciente échoue ; la psychanalyse y voit un rituel symbolique qui tient à distance une culpabilité inconsciente.
\item Les \textbf{actes manqués} : oublier un rendez-vous que l'on redoutait, casser l'objet offert par quelqu'un que l'on déteste secrètement.
\end{itemize}
\end{exemplebox}

Dans tous ces cas, la conduite a un \emph{sens}, mais ce sens est caché au sujet lui-même. La psychanalyse mérite ainsi le titre de « psychologie des profondeurs » : elle montre que l'homme n'est pas transparent à lui-même et que « ça » parle en lui à son insu.

\courssub{La valeur culturelle de la psychanalyse}

Au-delà du cabinet du médecin, la psychanalyse a profondément transformé la culture du XX\textsuperscript{e} siècle :

\begin{itemize}
\item \textbf{Littérature et art} : les écrivains (Proust, les surréalistes) et les cinéastes explorent les rêves, les désirs inavoués, la mémoire involontaire ; l'art surréaliste (Dalí) peint directement les images de l'inconscient.
\item \textbf{Pédagogie et compréhension de l'enfant} : Freud a montré que l'enfant n'est pas un « petit adulte » innocent, mais un être traversé de désirs et de conflits (complexe d'Œdipe). L'éducation moderne en a tiré la leçon : écouter l'enfant, respecter son histoire affective, ne pas briser par la violence éducative ce qui doit être élaboré par la parole.
\item \textbf{Langage courant} : des mots comme « refouler », « complexe », « refoulé », « acte manqué » sont entrés dans le vocabulaire de tous.
\end{itemize}

\begin{savaistu}[Inconscient et cultures africaines]
En Afrique, des penseurs comme l'écrivain nigérian \textbf{Wole Soyinka} (prix Nobel de littérature 1986) ont discuté la place de l'inconscient dans les cultures traditionnelles : le rêve y est souvent tenu pour un message, la divination et la transe pour des voies d'accès à ce qui échappe à la conscience veillée. On ne peut pas identifier simplement ces pratiques à la psychanalyse, mais elles montrent que l'idée d'une part cachée de l'esprit humain n'est pas une invention européenne : beaucoup de traditions, y compris camerounaises, savent depuis longtemps que l'homme ne se connaît pas tout entier.
\end{savaistu}

\courssub{Première limite : l'irréfutabilité (critique de Karl Popper)}

La critique la plus célèbre adressée à la psychanalyse vient du philosophe des sciences \textbf{Karl Popper}.

\begin{definition}[Falsifiabilité (Popper)]
Une théorie est \cle{scientifique} si et seulement si l'on peut imaginer une expérience ou une observation capable de la \emph{réfuter} : on dit alors qu'elle est \cle{falsifiable} (réfutable). Une théorie compatible avec tous les faits possibles, qui ne risque jamais d'être démentie, n'est pas une science, mais au mieux une interprétation, au pire une doctrine fermée.
\end{definition}

Or, selon Popper, la psychanalyse échappe précisément à cette épreuve. Prenons un exemple : Freud affirme qu'un patient agit par hostilité refoulée envers son père. Si le patient reconnaît cette hostilité, Freud a raison. Mais si le patient la nie avec force, l'analyste répondra que ce déni est une « résistance » qui \emph{confirme} le refoulement : Freud a encore raison ! Quoi qu'il arrive, la théorie est vérifiée ; aucune observation ne peut la contredire. Une théorie qui explique tout n'explique rien au sens scientifique du terme, car elle ne prend aucun risque devant les faits.

\begin{propriete}[Valeur herméneutique de la psychanalyse]
La psychanalyse possède une valeur \cle{herméneutique} (interprétative) plus qu'expérimentale : elle propose un sens pour des conduites énigmatiques (rêves, symptômes, mythes), comme on interprète un texte, plutôt qu'elle n'établit des lois vérifiables comme la physique. Elle appartient davantage aux sciences de l'interprétation qu'aux sciences expérimentales.
\end{propriete}

\courssub{Deuxième limite : la généralisation abusive}

Freud a construit sa théorie générale de l'esprit humain à partir d'un nombre restreint de cas cliniques : quelques dizaines de patients, pour l'essentiel des femmes de la bourgeoisie viennoise de la fin du XIX\textsuperscript{e} siècle, dans une société marquée par une morale sexuelle très rigide. Peut-on légitimement conclure de ces cas particuliers à une structure \emph{universelle} de la psyché humaine, valable pour tous les hommes, de toutes les époques et de toutes les cultures ? Rien n'est moins sûr. Le complexe d'Œdipe, par exemple, suppose une famille de type père-mère-enfant ; or les structures familiales africaines (famille élargie, rôle des oncles maternels, polygamie) ne coïncident pas toujours avec ce modèle. Généraliser sans précaution, c'est transformer une observation locale en dogme universel.

\courssub{Troisième limite : le pansexualisme}

On reproche enfin à Freud d'accorder une importance excessive à la \cle{sexualité} : c'est le reproche de \cle{pansexualisme}. Pour Freud, la libido (énergie du désir sexuel, au sens large) est le moteur principal de la vie psychique, y compris chez l'enfant (sexualité infantile) et dans des domaines apparemment étrangers à la sexualité (art, religion, vie sociale), qu'il ramène à des désirs sexuels « sublimés ». Ses propres disciples se sont détachés de lui sur ce point : \textbf{Jung} élargit la libido à une énergie psychique générale et ajoute un « inconscient collectif » ; \textbf{Adler} fait de la volonté de puissance et du sentiment d'infériorité, non de la sexualité, le ressort principal des conduites. Tout ramener au sexe, c'est risquer de déformer la réalité humaine en la réduisant à un seul de ses aspects.

\begin{popfigure}
\begin{center}
\begin{tikzpicture}
\node[draw=popdark, thick, fill=popcream, rounded corners=4pt, text width=6.6cm, align=left, font=\small] (val) at (0,0) {
\textbf{\textcolor{popgreen}{VALEUR de la psychanalyse}}\\[2pt]
$\bullet$ \textbf{Thérapeutique} : cure par la parole, soulagement des névroses (Freud, Breuer).\\
$\bullet$ \textbf{Explicative} : rend compte des conduites absurdes pour la conscience : phobies, actes manqués, rêves (Freud).\\
$\bullet$ \textbf{Culturelle} : influence sur l'art, la littérature, la pédagogie et la compréhension de l'enfant.
};
\node[draw=popdark, thick, fill=popcream, rounded corners=4pt, text width=6.6cm, align=left, font=\small, below=0.5cm of val] (lim) {
\textbf{\textcolor{poporange}{LIMITES de la psychanalyse}}\\[2pt]
$\bullet$ \textbf{Irréfutable} : aucune expérience ne peut la démentir ; ce n'est pas une science expérimentale (Popper).\\
$\bullet$ \textbf{Généralisation abusive} : théorie universelle bâtie sur des cas cliniques particuliers (Vienne, vers 1900).\\
$\bullet$ \textbf{Pansexualisme} : importance excessive accordée à la sexualité (critiques de Jung et d'Adler).
};
\end{tikzpicture}
\end{center}
\legende{Bilan argumentatif : valeur et limites de la psychanalyse, avec les auteurs de référence.}
\end{popfigure}

\courssub{Débat guidé : la psychanalyse est-elle une science ?}

Reprenons les critères classiques de la scientificité et interrogeons la psychanalyse à leur lumière.

\begin{center}
\begin{tabularx}{\linewidth}{|l|X|X|}
\hline
\rowcolor{popblueL}
\textbf{Critère} & \textbf{Exigence} & \textbf{La psychanalyse y répond-elle ?} \\
\hline
Expérimentation & Reproduire les phénomènes en conditions contrôlées & Non : la cure est une relation unique, non reproductible en laboratoire. \\
\hline
Falsifiabilité & Possibilité d'être réfutée par les faits & Non, selon Popper : tout événement peut être interprété comme une confirmation. \\
\hline
Objectivité & Mesures indépendantes de l'observateur & Partiellement : l'interprétation dépend de l'analyste ; mais les symptômes et les récits sont des faits observables. \\
\hline
\end{tabularx}
\end{center}

Le débat reste ouvert. \emph{Thèse} : la psychanalyse est une science, car elle procède par observation rigoureuse de cas, formule des hypothèses et obtient des résultats thérapeutiques vérifiables (disparition de symptômes). \emph{Antithèse} : elle n'est pas une science, car ses concepts (Ça, refoulement, Œdipe) ne sont ni mesurables ni réfutables. \emph{Synthèse nuancée} : la psychanalyse n'est pas une science expérimentale de la nature, mais une \emph{pratique interprétative} à visée thérapeutique, dont la valeur se mesure moins à la vérité démontrée qu'à l'intelligibilité et au soulagement qu'elle apporte.

\begin{attention}[Erreur fréquente à éviter]
Au Baccalauréat technique, deux attitudes sont également fautives : \textbf{rejeter totalement Freud} (« ce n'est pas scientifique, donc cela ne vaut rien ») ou \textbf{l'accepter sans critique} (« Freud a tout expliqué »). L'épreuve exige un \emph{jugement nuancé} : reconnaître la fécondité de la découverte de l'inconscient tout en discutant rigoureusement le statut scientifique de la psychanalyse. Une copie qui tranche sans argumenter, dans un sens ou dans l'autre, est une copie manquée.
\end{attention}

\begin{exoresolu}[Sujet de dissertation : « La psychanalyse est-elle une science ? »]
Construire le plan détaillé d'une dissertation répondant à ce sujet.
\tcblower
\textbf{Analyse du sujet.} Termes clés : « psychanalyse » (méthode d'investigation et de traitement de l'inconscient fondée par Freud) et « science » (connaissance objective, vérifiable, fondée sur l'expérience et la réfutabilité). Problématique : la psychanalyse satisfait-elle aux critères de la scientificité, ou n'est-elle qu'une interprétation séduisante de l'esprit humain ?

\textbf{I. La psychanalyse présente des traits scientifiques.}
Elle part de l'observation méthodique de faits (symptômes, récits de rêves, lapsus), formule des hypothèses explicatives (refoulement, complexe d'Œdipe) et vérifie sa valeur par ses effets thérapeutiques (cure par la parole, cas d'Anna O.). Freud, médecin formé aux sciences, a voulu fonder une « psychologie scientifique ».

\textbf{II. Mais elle échappe aux critères stricts de la science expérimentale.}
Critique de Popper : ses thèses sont irréfutables (le déni du patient est interprété comme résistance, donc comme confirmation). Ses concepts ne sont ni mesurables ni expérimentalement contrôlables ; elle généralise à partir de cas cliniques particuliers ; le pansexualisme réduit l'homme à un seul ressort (critiques de Jung et d'Adler).

\textbf{III. Une pratique interprétative plutôt qu'une science expérimentale.}
La psychanalyse a une valeur herméneutique : elle donne du sens à ce que la conscience ne comprend pas. Son modèle est moins la physique que l'interprétation des textes et des rêves. Sa fécondité culturelle (art, littérature, pédagogie) et thérapeutique demeure, même si son statut scientifique est contesté.

\textbf{Conclusion.} La psychanalyse n'est pas une science au sens expérimental, mais elle n'est pas pour autant sans valeur : elle est une méthode d'interprétation et de soin qui a révélé la part obscure de l'esprit humain.
\end{exoresolu}

\courssub{Bilan de la leçon : l'homme, être divisé mais capable de lucidité}

Que retenir de l'ensemble de cette leçon ? D'un côté, la \cle{conscience} n'est pas toute-puissante : Descartes croyait qu'« il n'y a rien en nous dont nous n'ayons connaissance », mais Freud a montré que des forces inconscientes nous agitent à notre insu. De l'autre, l'\cle{inconscient} ne fait pas de nous des esclaves : la psychanalyse elle-même suppose que l'on peut, par la parole et le travail de la cure, rendre conscient ce qui était refoulé — « le Moi doit déloger le Ça », disait Freud. L'homme apparaît ainsi comme un \emph{être divisé}, traversé de conflits entre désir et raison, mais précisément capable de prendre conscience de cette division : c'est cette lucidité, jamais totale mais toujours possible, qui fonde sa liberté et sa dignité.

\begin{aretenir}[L'essentiel de la section]
\begin{itemize}
\item La cure psychanalytique repose sur l'association libre, l'interprétation et le transfert : guérir, c'est rendre conscient le refoulé.
\item Valeur : elle explique ce que la conscience n'explique pas (névroses, phobies) et a fécondé la culture (art, littérature, pédagogie).
\item Limites : irréfutabilité (Popper), généralisation abusive à partir de cas particuliers, pansexualisme (Jung, Adler).
\item Statut : pratique herméneutique et thérapeutique plus que science expérimentale.
\item Bilan : ni toute-puissance de la conscience, ni esclavage à l'inconscient ; l'homme est un être divisé mais capable de lucidité.
\end{itemize}
\end{aretenir}

\begin{methode}[Rédiger une conclusion de dissertation]
\begin{enumerate}
\item \textbf{Répondre directement à la problématique} posée dans l'introduction, en une ou deux phrases claires (ex. : « La psychanalyse n'est donc pas une science expérimentale, mais une pratique interprétative féconde »).
\item \textbf{Dresser un bilan nuancé} du parcours : rappeler en quoi les deux thèses opposées avaient chacune une part de vérité, et comment la synthèse les dépasse.
\item \textbf{Ouvrir sur une question nouvelle}, liée au sujet mais distincte de lui (ex. : « Cette division de l'homme ne remet-elle pas en cause l'idée même de liberté ? »). L'ouverture ne doit jamais être un nouveau développement : une phrase suffit.
\end{enumerate}
\end{methode}

\begin{exercice}[Pour s'entraîner]
\begin{enumerate}
\item Définir la falsifiabilité selon Popper et expliquer en quoi elle permet de critiquer la psychanalyse.
\item Citer deux valeurs et deux limites de la psychanalyse, en associant à chacune un auteur ou un exemple.
\item Sujet de dissertation : « L'inconscient nous délivre-t-il de la responsabilité de nos actes ? » (analyser le sujet, dégager la problématique, construire un plan en trois parties).
\end{enumerate}
\end{exercice}

\begin{corrige}[Exercice 1]
La falsifiabilité est, selon Popper, le critère qui distingue une théorie scientifique d'une doctrine non scientifique : une théorie est scientifique si l'on peut concevoir une expérience ou une observation capable de la réfuter. Or la psychanalyse interprète tout événement comme une confirmation de ses thèses : si le patient reconnaît le conflit que lui révèle l'analyste, la théorie est vérifiée ; s'il le nie, son déni est lu comme une « résistance » qui confirme le refoulement. Aucun fait ne pouvant la contredire, la psychanalyse échappe au critère de falsifiabilité : elle n'est donc pas, selon Popper, une science au sens strict.
\end{corrige}

\begin{corrige}[Exercice 2]
Valeurs : (1) valeur thérapeutique — la cure par la parole soulage les névroses (Freud et Breuer, cas d'Anna O.) ; (2) valeur explicative — elle rend intelligibles les phobies, les rêves et les actes manqués que la conscience seule n'explique pas (Freud). Limites : (1) irréfutabilité — les thèses freudiennes échappent à la vérification expérimentale (Popper) ; (2) pansexualisme — l'importance excessive accordée à la sexualité, critiquée par Jung et Adler. (On pouvait aussi citer la généralisation abusive à partir de cas cliniques particuliers.)
\end{corrige}

\begin{corrige}[Exercice 3]
\textbf{Analyse} : « inconscient » (part psychique qui échappe à la conscience) ; « responsabilité » (devoir de répondre de ses actes, qui suppose la conscience et la liberté). \textbf{Problématique} : si nos actes sont déterminés par des forces inconscientes, pouvons-nous encore en être tenus pour responsables ? \textbf{Plan} : I. L'inconscient semble nous délivrer de la responsabilité : nos actes ont des causes cachées (Freud, actes manqués, névroses) que nous n'avons pas choisies. II. Pourtant la responsabilité demeure : la conscience peut éclairer et maîtriser l'inconscient ; la cure elle-même suppose que le sujet peut devenir lucide ; la morale et le droit présupposent cette capacité. III. Synthèse : l'inconscient limite notre responsabilité sans l'abolir ; être responsable, c'est précisément assumer le travail de lucidité sur soi — l'homme est divisé, mais capable de se connaître.
\end{corrige}

\newpage

\coursec{Activité d'intégration}

\courssub{Objectifs de l'activité}

À l'issue de cette activité d'intégration, l'élève sera capable de :

\begin{itemize}
\item \cle{Appliquer} les notions de \cle{conscience} et d'\cle{inconscient} à des situations concrètes de la vie courante (lapsus, rêve, oubli, acte manqué) ;
\item \cle{Distinguer}, dans un document donné, ce qui relève de la conscience (ce dont le sujet a connaissance immédiate) de ce qui relève de l'inconscient (contenus psychiques inaccessibles à la conscience mais agissants) ;
\item \cle{Construire} une représentation schématique (iceberg ou topique freudienne) illustrant le rapport entre conscience et inconscient ;
\item \cle{Rédiger et justifier} une démarche argumentée : produire une réponse personnelle, structurée et étayée d'auteurs, à la question « Sommes-nous maîtres de nos actes ? » ;
\item \cle{Confronter} les thèses de Descartes (toute-puissance de la conscience), de Spinoza (ignorance des causes de nos désirs) et de Freud (l'inconscient psychique) pour évaluer la valeur et les limites de la conscience.
\end{itemize}

\begin{aretenir}[Compétences travaillées]
Savoir-faire officiels mobilisés : \emph{appliquer les notions de l'unité à des situations concrètes de la vie courante} et \emph{rédiger et justifier une démarche, une réponse ou une conclusion}. Cette activité prépare directement à l'épreuve de dissertation du Probatoire et du Baccalauréat technique.
\end{aretenir}

\courssub{Situation}

\textbf{TP guidé : Le journal des mystères du comportement — enquête sur l'inconscient au quotidien.}

Au lycée technique de Nkongsamba, la classe de Terminale F4 prépare le Baccalauréat. Depuis quelques semaines, la conseillère d'orientation, Mme~Etoundi, collecte des « mystères du comportement » : de petits faits étranges que les élèves s'envoient sur le groupe WhatsApp de la classe. Pourquoi un élève sérieux oublie-t-il soudain le nom de son meilleur ami ? Pourquoi rêve-t-on la veille des examens ? Pourquoi dit-on un mot à la place d'un autre, précisément quand on veut cacher quelque chose ? La classe décide de mener l'enquête à l'aide de quatre documents.

\textbf{Document 1 — Le lapsus célèbre.} Dans \emph{Psychopathologie de la vie quotidienne} (1901), Freud rapporte le cas d'un jeune homme qui, voulant citer un vers latin où figure le mot \emph{aliquis} (« quelqu'un »), oublie précisément ce mot. Or le jeune homme venait d'apprendre qu'une femme qu'il fréquentait avait un retard de règles : il redoutait une paternité non désirée. Freud commente : l'oubli du mot \emph{aliquis} n'est pas un hasard ; il exprime, de manière déguisée, un désir refoulé — le souhait que « quelqu'un » ne vienne pas troubler sa vie. « Le lapsus et l'oubli ont un sens : ils trahissent une intention inconsciente. »

\textbf{Document 2 — Le rêve de Marlyse.} Marlyse, élève de Terminale, raconte : « La veille de l'épreuve de mathématiques, j'ai rêvé que j'arrivais au lycée, mais le portail était fermé. Je courais, je courais, et mes jambes devenaient lourdes comme du béton. Puis mon père, décédé il y a deux ans, m'ouvrait le portail en souriant, et je me réveillais en pleurant. » Marlyse précise qu'elle redoute de décevoir la mémoire de son père, qui voulait la voir ingénieure.

\textbf{Document 3 — Spinoza, \emph{Éthique} (extrait adapté).} « Les hommes se croient libres parce qu'ils ont conscience de leurs désirs et de leurs volontés, mais ils ignorent les causes qui les déterminent à désirer et à vouloir. Tel est le cas d'un enfant qui croit désirer librement le lait dont il se nourrit. »

\textbf{Document 4 — Enquête statistique de la classe.} Mme~Etoundi a interrogé les 30 élèves de la classe sur les « actes manqués » commis pendant le mois précédant les examens blancs. Résultats :

\begin{center}
\begin{tabularx}{\linewidth}{|l|X|X|X|}\hline
\rowcolor{popblueL} \textbf{Type d'acte manqué} & \textbf{Nombre d'élèves concernés} & \textbf{Fréquence totale des faits} & \textbf{Part des élèves (\%)} \\ \hline
Oubli de clés ou de cahier & 18 & 41 & 60 \\ \hline
Lapsus (mot pour un autre) & 15 & 33 & 50 \\ \hline
Oubli d'un rendez-vous & 9 & 14 & 30 \\ \hline
Rêve d'examen ou d'échec & 21 & 27 & 70 \\ \hline
\rowcolor{popcream} \textbf{Total des faits recensés} & --- & \textbf{115} & --- \\ \hline
\end{tabularx}
\end{center}

\begin{center}
\begin{popfigure}
\begin{tikzpicture}
\begin{axis}[
    ybar, bar width=0.55cm,
    width=0.85\linewidth, height=5.2cm,
    symbolic x coords={Oubli d'objet, Lapsus, Oubli RDV, Rêve d'examen},
    xtick=data, x tick label style={font=\small, rotate=12, anchor=east},
    ylabel={Nombre d'élèves (sur 30)},
    ymin=0, ymax=25,
    nodes near coords, nodes near coords style={font=\small},
    axis line style={popline},
]
\addplot[fill=popblue!70, draw=popblue] coordinates {(Oubli d'objet,18) (Lapsus,15) (Oubli RDV,9) (Rêve d'examen,21)};
\end{axis}
\end{tikzpicture}
\legende{Fréquence des actes manqués dans la classe de Terminale F4 (données fictives, enquête de Mme~Etoundi, 30 élèves).}
\end{popfigure}
\end{center}

\textbf{Problème posé à la classe :} ces petits « accidents » de la vie quotidienne sont-ils de purs hasards sans signification, ou révèlent-ils que notre conscience n'est pas maîtresse dans sa propre maison ?

\courssub{Consignes}

Travail en groupes de 4 à 5 élèves. Chaque groupe répond par écrit aux questions suivantes, puis prépare une restitution orale de 5 minutes.

\begin{enumerate}
\item \textbf{Analyse des documents.} Pour chacun des quatre documents, identifiez précisément : (a) ce qui relève de la \cle{conscience} (ce que le sujet sait, perçoit, déclare) ; (b) ce qui relève de l'\cle{inconscient} (contenu caché, refoulé ou ignoré, qui agit malgré le sujet). Présentez vos réponses dans un tableau à deux colonnes.
\item \textbf{Interprétation.} En vous appuyant sur le document 1, expliquez en quoi le lapsus ou l'oubli possède, selon Freud, un \emph{sens}. Pourquoi Freud refuse-t-il de les considérer comme de simples hasards ?
\item \textbf{Lecture de Spinoza.} Dans le document 3, que signifie la formule « les hommes ont conscience de leurs désirs, mais ignorent les causes qui les déterminent » ? En quoi cette thèse anticipe-t-elle l'idée d'inconscient, sans être encore la psychanalyse ?
\item \textbf{Exploitation statistique.} À partir du document 4 : (a) vérifiez les pourcentages de la dernière colonne ; (b) calculez la part, en pourcentage, de chaque type de fait dans le total des 115 faits recensés ; (c) que suggère la forte fréquence des rêves d'examen (70~\% des élèves) sur l'état psychique de la classe ?
\item \textbf{Schématisation.} Construisez un diagramme illustrant le rapport conscience/inconscient, au choix : un \emph{iceberg} (partie émergée / partie immergée) ou la \emph{topique freudienne} (Ça, Moi, Surmoi). Placez-y les exemples des documents (lapsus, rêve, oubli de clés).
\item \textbf{Production argumentée.} Rédigez en 15 lignes maximum une réponse personnelle et structurée à la question : \textbf{« Sommes-nous maîtres de nos actes ? »} Vous mobiliserez \emph{au moins deux auteurs} étudiés dans la leçon (Descartes, Spinoza, Freud) et \emph{au moins un exemple camerounais} (tiré des documents ou de votre expérience).
\end{enumerate}

\courssub{Ressources à mobiliser}

\begin{aretenir}[Notions utiles de la leçon]
\begin{itemize}
\item \textbf{Conscience} : connaissance immédiate que le sujet a de lui-même et du monde (Descartes : « Je pense, donc je suis » — la conscience est le fondement certain de toute connaissance). Sa \emph{valeur} : elle rend possible la connaissance, la liberté et la responsabilité morale.
\item \textbf{Limites de la conscience} : elle ignore les causes de ses propres désirs (Spinoza) ; elle peut être trompée par les sens et par elle-même ; elle ne contrôle pas tout le psychisme.
\item \textbf{Inconscient freudien} : ensemble des contenus psychiques (désirs refoulés, souvenirs pénibles) inaccessibles à la conscience, mais qui s'expriment de manière déguisée dans les \emph{lapsus}, les \emph{actes manqués}, les \emph{rêves} (« voie royale vers l'inconscient ») et les symptômes.
\item \textbf{Topique freudienne} : le \emph{Ça} (pulsions), le \emph{Moi} (instance de réalité, médiateur), le \emph{Surmoi} (interdits moraux hérités des parents et de la société).
\item \textbf{Psychanalyse} : méthode d'investigation et de cure fondée par Freud. \emph{Valeur} : elle explique des conduites jusque-là incompréhensibles. \emph{Limites} : ses interprétations sont difficilement vérifiables scientifiquement ; elle risque de nier la liberté et la responsabilité.
\end{itemize}
\end{aretenir}

\begin{methode}[Méthode de la réponse argumentée courte (15 lignes)]
\begin{enumerate}
\item \textbf{Poser la thèse} dès la première phrase (oui / non / nuancée).
\item \textbf{Premier argument} : un auteur + un exemple précis.
\item \textbf{Deuxième argument} : un autre auteur, éventuellement opposé au premier + un exemple camerounais.
\item \textbf{Nuance} : reconnaître la part de vérité de la thèse adverse.
\item \textbf{Conclusion} : formuler une position personnelle justifiée en une ou deux phrases.
\end{enumerate}
\end{methode}

\courssub{Démarche et corrigé commenté}

\begin{exoresolu}[Corrigé commenté du TP « Le journal des mystères du comportement »]
\textbf{Énoncé.} Analyser les quatre documents, construire un diagramme conscience/inconscient, exploiter les données statistiques et rédiger une réponse argumentée à la question « Sommes-nous maîtres de nos actes ? ».

\tcblower

\textbf{Étape 1 — Analyse des documents (question 1).}

\begin{center}
\begin{tabularx}{\linewidth}{|l|X|X|}\hline
\rowcolor{popblueL} \textbf{Document} & \textbf{Relève de la conscience} & \textbf{Relève de l'inconscient} \\ \hline
1. Le lapsus & Le jeune homme veut citer le vers ; il sait ce qu'il cherche à dire. & L'oubli du mot \emph{aliquis} exprime un désir refoulé : ne pas être père. \\ \hline
2. Le rêve de Marlyse & Marlyse sait qu'elle redoute l'épreuve ; elle raconte son rêve. & Le portail fermé et les jambes lourdes figurent l'angoisse d'échec ; le père souriant exprime le désir d'être approuvée. \\ \hline
3. Spinoza & L'enfant a conscience de désirer le lait. & Il ignore les causes (la faim, le besoin biologique) qui déterminent ce désir. \\ \hline
4. Statistiques & Les élèves déclarent et comptent leurs oublis. & La fréquence des actes manqués trahit une tension psychique collective (stress du Bac) que les élèves ne maîtrisent pas. \\ \hline
\end{tabularx}
\end{center}

\emph{Commentaire :} la conscience est partout présente comme surface (déclarations, intentions), mais chaque document révèle un arrière-plan qui échappe au sujet.

\textbf{Étape 2 — Le sens du lapsus (question 2).} Pour Freud, le lapsus n'est pas un accident : c'est un \emph{compromis} entre une intention consciente (citer correctement) et un désir inconscient refoulé (fuir la paternité). Le refoulement empêche le désir d'accéder directement à la conscience ; il ressurgit donc de manière déguisée, par l'oubli du mot. Conclusion freudienne : \emph{rien n'arrive par hasard dans la vie psychique}.

\textbf{Étape 3 — Lecture de Spinoza (question 3).} Spinoza montre que la conscience connaît les \emph{effets} (je désire) mais ignore les \emph{causes} (pourquoi je désire). L'homme se croit libre comme une pierre lancée dans l'air se croirait libre de sa chute si elle était consciente. Cette « illusion de la liberté » annonce l'inconscient ; mais Spinoza reste un rationaliste : ces causes sont en principe connaissables par la raison, tandis que pour Freud l'inconscient est un réservoir de désirs refoulés que seule la psychanalyse peut atteindre.

\textbf{Étape 4 — Exploitation statistique (question 4).}
\begin{itemize}
\item Vérification : $18/30 = 0{,}60$, soit 60~\% ; $15/30 = 0{,}50$, soit 50~\% ; $9/30 = 0{,}30$, soit 30~\% ; $21/30 = 0{,}70$, soit 70~\%. Les pourcentages sont exacts.
\item Parts dans le total des 115 faits : oublis d'objets $41/115 \approx 35{,}7$~\% ; lapsus $33/115 \approx 28{,}7$~\% ; oublis de rendez-vous $14/115 \approx 12{,}2$~\% ; rêves d'examen $27/115 \approx 23{,}5$~\%. (Somme : 100~\%.).
\item Interprétation : 70~\% des élèves ont rêvé d'examen ou d'échec. Cette fréquence élevée suggère une angoisse collective liée au Baccalauréat : l'inconscient « travaille » la nuit ce que la conscience refoule le jour.
\end{itemize}

\textbf{Étape 5 — Le diagramme (question 5).} Modèle attendu : un iceberg dont la partie émergée (environ 1/10) représente la conscience (pensées claires, intentions déclarées) et la partie immergée l'inconscient (désirs refoulés, angoisses, souvenirs). On y place : le lapsus et l'oubli de clés à la ligne de flottaison (ils font surface), le désir refoulé du jeune homme et l'angoisse de Marlyse dans la profondeur. Variante acceptée : la topique Ça (pulsions) / Moi (médiateur) / Surmoi (interdits : « réussir pour honorer papa »).

\textbf{Étape 6 — Production argumentée (question 6), corrigé type.}

\emph{« Nous ne sommes pas totalement maîtres de nos actes, mais nous ne sommes pas pour autant des marionnettes. Spinoza a raison de rappeler que nous avons conscience de nos désirs sans en connaître les causes : l'élève camerounais qui "choisit" de réviser toute la nuit ignore souvent que c'est la peur de décevoir sa famille qui le détermine. Freud va plus loin : le lapsus du jeune homme et les 115 actes manqués de notre classe montrent qu'une part de nous-mêmes nous échappe et s'exprime malgré nous. Cependant, Descartes a établi que la conscience est le fondement de toute certitude et de la responsabilité : sans elle, aucun tribunal, aucune morale ne serait possible. La vérité est donc nuancée : l'inconscient limite notre maîtrise, mais prendre conscience de ces forces cachées — par la réflexion ou l'analyse — élargit précisément notre liberté. Connaître ce qui nous détermine, c'est déjà commencer à s'en libérer. »}

\emph{Commentaire du corrigé :} la réponse pose une thèse nuancée dès la première ligne, mobilise trois auteurs (Spinoza, Freud, Descartes), intègre deux exemples camerounais (l'élève qui révise, l'enquête de la classe), reconnaît la thèse adverse et conclut par une position personnelle justifiée : c'est exactement la démarche attendue à l'épreuve.
\end{exoresolu}

\courssub{Bilan de l'activité}

Cette enquête a permis de mettre en évidence trois acquis essentiels de la leçon :

\begin{enumerate}
\item \textbf{La conscience n'est pas toute-puissante.} Les faits les plus banals de la vie quotidienne (oublier ses clés, dire un mot pour un autre, rêver d'échec) montrent que le sujet n'est pas transparent à lui-même : « le Moi n'est pas maître dans sa propre maison », selon la formule de Freud.
\item \textbf{L'inconscient donne du sens à ce qui semblait absurde.} Là où le sens commun voit du hasard, la psychanalyse déchiffre des messages déguisés : le lapsus, l'acte manqué et le rêve deviennent des documents à interpréter. Les données statistiques de la classe ont montré que ces phénomènes ne sont pas des cas rares mais des réalités massives (115 faits en un mois pour 30 élèves).
\item \textbf{Reconnaître l'inconscient ne supprime pas la liberté, il la rend possible.} Comprendre les causes cachées de nos conduites (angoisse du Bac, désir de plaire, peur de décevoir) est la première étape pour les maîtriser : c'est la position synthétique que l'élève doit savoir défendre en dissertation, entre le rationalisme de Descartes et le déterminisme psychique de Freud.
\end{enumerate}

\begin{attention}[Erreurs fréquentes à éviter]
\begin{itemize}
\item Confondre \emph{inconscient} et \emph{non-conscient} : le cœur qui bat n'est pas un fait d'inconscient freudien ; l'inconscient est fait de \emph{désirs refoulés} qui cherchent à s'exprimer.
\item Affirmer que Freud « prouve scientifiquement » l'inconscient : la psychanalyse est une méthode d'\emph{interprétation}, dont les résultats sont difficilement vérifiables expérimentalement — c'est sa principale limite.
\item Dans la production écrite, juxtaposer les auteurs sans les articuler : chaque référence doit servir un argument, pas décorer la copie.
\end{itemize}
\end{attention}

\begin{exercice}[Pour aller plus loin]
Pendant une semaine, tenez votre propre « journal des mystères du comportement » : notez chaque lapsus, oubli ou rêve marquant, et proposez pour chacun une hypothèse d'interprétation consciente et une hypothèse inconsciente. En fin de semaine, rédigez un court paragraphe argumenté : « Mon expérience confirme-t-elle ou infirme-t-elle la thèse de Freud ? »
\end{exercice}

\newpage

\coursec{Méthodes et exercices résolus}

\coursec{La Conscience et l’Inconscient}

\courssub{Introduction : problématique et définitions}

\begin{definition}[Conscience]
La \cle{conscience} (du latin \emph{cum-scire}, « savoir avec ») est la présence de l'esprit à lui-même et au monde. Elle se dédouble en deux formes indissociables : la \textbf{conscience intentionnelle} (être conscient \emph{de} quelque chose : percevoir, penser, vouloir) et la \textbf{conscience réflexive} (se savoir conscient, se saisir comme un « je » identique à lui-même). Elle est le fondement de la subjectivité et de la responsabilité.
\end{definition}

\begin{definition}[Inconscient (sens freudien)]
L'\cle{inconscient} freudien n'est pas une simple absence de conscience (comme le sommeil ou la distraction), ni l'automatisme physiologique. C'est un \textbf{système psychique dynamique} constitué de représentations, de désirs et de pulsions \emph{refoulés} (rejetés hors de la conscience parce qu'incompatibles avec les exigences morales ou sociales), qui échappent au « moi » conscient mais déterminent pensées, affects et conduites. Il obéit à des lois qui lui sont propres (absence de contradiction, intemporalité, déplacement, condensation).
\end{definition}

\begin{aretenir}[Distinctions fondamentales]
\begin{itemize}
\item \textbf{Conscience / Inconscient (Freud) :} opposition topique (lieux psychiques) et dynamique (refoulement).
\item \textbf{Conscience / Non-conscient :} le non-conscient désigne tout ce qui n'est pas conscient à un instant $t$ (mémoire latente, automatismes corporels, traitement cognitif pré-attentif). L'inconscient freudien en est une partie structurée et conflictuelle.
\item \textbf{Préconscient :} antichambre de la conscience ; contenus non actuels mais accessibles sans résistance (souvenirs, savoirs).
\item \textbf{Mauvaise foi (Sartre) :} mensonge à soi-même \emph{en conscience} ; alternative à l'inconscient pour expliquer l'opacité de soi.
\end{itemize}
\end{aretenir}

\begin{savaistu}[Contexte camerounais]
Dans nos sociétés, l'expression « \emph{Il a agi sans le savoir} » ou « \emph{C'est plus fort que lui} » traduit l'intuition commune d'un inconscient. Les rites d'initiation, l'interprétation des songes chez les Bamiléké ou les Béti, la consultation des devins (nganga, mbias) témoignent d'une anthropologie qui ne réduit pas l'homme à sa conscience claire. La philosophie technique doit articuler cette sagesse pratique avec les concepts rigoureux de Descartes, Kant, Freud ou Sartre.
\end{savaistu}

\courssub{La nature de la conscience}

\begin{propriete}[Intentionalité et réflexivité]
La conscience est toujours \textbf{conscience de quelque chose} (intentionalité, Brentano/Husserl) : elle vise un objet (perception, idée, valeur). Par sa réflexivité, elle se prend elle-même pour objet : « Je pense que je pense ». Cette structure \emph{sujet-objet} interne fonde l'unité du « je » (Descartes, Kant).
\end{propriete}

\begin{exemplebox}[Exemple concret]
Un élève de Terminale technique à Garoua lit un énoncé de dissertation (conscience intentionnelle : il vise le texte). Soudain, il se dit : « Je ne comprends pas ce mot » (conscience réflexive : il se saisit lui-même en train de ne pas comprendre). Cette bascule est le propre de l'esprit humain.
\end{exemplebox}

\courssub{La valeur de la conscience}

\begin{propriete}[Valeur épistémologique : le fondement de la certitude (Descartes)]
Avec le \cle{cogito} (\emph{« Je pense, donc je suis »}, \emph{Discours de la méthode}, \emph{Méditations}), Descartes montre que la conscience de soi est la première vérité indubitable. Même si je doute de tout (monde, corps, mémoire), je ne peux douter que je doute, donc que je pense, donc que j'existe. La conscience réflexive est l'archimède de la connaissance.
\end{propriete}

\begin{propriete}[Valeur morale : condition de la liberté et de la responsabilité (Kant)]
Pour Kant (\emph{Fondements de la métaphysique des mœurs}), la conscience morale (le \emph{fait de la raison}) révèle la loi morale en moi. Être conscient de ses maximes d'action, c'est pouvoir les soumettre à l'impératif catégorique (« Agis uniquement d'après la maxime qui fait que tu puisses vouloir en même temps qu'elle devienne loi universelle »). Sans conscience, pas d'autonomie, pas de devoir, pas de responsabilité imputable.
\end{propriete}

\begin{exemplebox}[Valeur pratique]
Un technicien en maintenance à la SONARA (Limbe) détecte une anomalie sur une turbine. Sa conscience professionnelle (savoir-faire + vigilance + responsabilité) lui permet d'intervenir avant l'accident. La conscience est ici pouvoir d'agir \emph{en connaissance de cause}.
\end{exemplebox}

\courssub{Les limites de la conscience}

\begin{propriete}[L'erreur cartésienne et l'illusion de transparence]
Descartes lui-même (\emph{Méditations I}) montre que la conscience spontanée se trompe (illusions sensorielles, rêves). La conscience immédiate n'est pas un critère de vérité absolu ; elle nécessite le doute méthodique et la raison pour se corriger.
\end{propriete}

\begin{propriete}[La critique nietzschéenne : la conscience comme surface]
Nietzsche (\emph{Le Gai Savoir}, \emph{Par-delà bien et mal}) compare la conscience à la surface d'une mer profonde : elle n'est qu'un « organe » tardif, outil de communication et de survie sociale, qui \emph{simplifie} et \emph{falsifie} la vie instinctive profonde. « Nous ne sommes pas les maîtres de notre conscience, nous en sommes les interprètes. »
\end{propriete}

\begin{propriete}[La découverte freudienne : l'inconscient comme autre scène]
Freud (\emph{La Psychopathologie de la vie quotidienne}, \emph{Introduction à la psychanalyse}) démontre que la conscience ignore ses propres motifs. Refoulement, pulsions, désirs infantiles agissent à l'insu du « moi ». La conscience n'est qu'une « façade » ; le moi « n'est pas maître dans sa propre maison ».
\end{propriete}

\begin{propriete}[Alain et Bergson : l'erreur sur ses propres mobiles]
\begin{itemize}
\item \textbf{Alain} (\emph{Propos sur le bonheur}) : « Celui qui croit agir par générosité agit souvent par vanité. » La conscience immédiate prend l'effet pour la cause ; elle rationalise \emph{a posteriori}.
\item \textbf{Bergson} (\emph{L'Énergie spirituelle}) : La conscience spontanée fige le moi vivant en concepts fixes ; elle « découpe » le flux continu de la durée pour l'action, perdant la richesse du moi profond.
\end{itemize}
\end{propriete}

\begin{propriete}[Sartre : la mauvaise foi, alternative à l'inconscient]
Sartre (\emph{L'Être et le Néant}) refuse l'inconscient comme « réserve » obscure. Pour lui, l'opacité de soi vient de la \cle{mauvaise foi} : la conscience \emph{sait} ce qu'elle cache, mais se ment à elle-même pour fuir l'angoisse de sa liberté. Exemple : la femme qui laisse sa main dans celle de son ami « sans s'en apercevoir » pour ne pas avoir à choisir (accepter ou refuser).
\end{propriete}

\courssub{L'inconscient freudien}

\begin{definition}[Première topique (1900) : Cs / Pcs / Incs]
\begin{itemize}
\item \textbf{Conscient (Cs) :} perceptions, pensées actuelles, soumis à la réalité.
\item \textbf{Préconscient (Pcs) :} souvenirs, savoirs latents, accessibles à la conscience sans travail.
\item \textbf{Inconscient (Inc) :} refoulé, pulsions (sexuelles, agressives), représentations interdites. Inaccessible directement ; ne connaît ni temps, ni négation, ni contradiction.
\end{itemize}
\end{definition}

\begin{definition}[Seconde topique (1923) : Ça / Moi / Surmoi]
\begin{itemize}
\item \textbf{Ça (Id) :} réservoir des pulsions, irrationnel, cherche la satisfaction immédiate (principe de plaisir).
\item \textbf{Moi (Ego) :} instance de réalité, médiateur entre Ça, Surmoi et monde extérieur ; siège de la conscience et de la défense (refoulement, sublimation, rationalisation...).
\item \textbf{Surmoi (Superego) :} intériorisation des interdits parentaux/sociaux ; instance morale, juge sévère (idéal du moi, conscience morale).
\end{itemize}
\end{definition}

\begin{propriete}[Mécanismes et manifestations de l'inconscient]
\begin{itemize}
\item \textbf{Refoulement (Verdrängung) :} opération originaire qui chasse au Incs les représentations incompatibles avec le Moi/Surmoi.
\item \textbf{Formation de compromis :} le symptôme, le rêve, le lapsus, l'acte manqué, le mot d'esprit sont des \emph{formations de l'inconscient} : ils satisfont partiellement le désir refoulé tout en déjouant la censure du Moi/Surmoi.
\item \textbf{Travail du rêve (Traumarbeit) :} \emph{condensation} (plusieurs idées en une image), \emph{déplacement} (charge affective déplacée sur un élément mineur), \emph{figuration} (pensées abstraites $\rightarrow$ images visuelles), \emph{élaboration secondaire} (récit cohérent au réveil). Le rêve est la « voie royale » vers l'inconscient.
\item \textbf{Acte manqué / Lapsus (Fehlleistung) :} action ratée (oubli, erreur de parole, lecture fautive) qui réalise un désir inconscient contraire à l'intention consciente.
\end{itemize}
\end{propriete}

\begin{exemplebox}[Cas camerounais : le lapsus du ministre]
Lors d'une inauguration de barrage à Lom Pangar, un officiel déclare : « Je suis heureux d'inaugurer cet \emph{échec}... euh, cet \emph{ouvrage} majeur. » Le lapsus (\emph{échec} pour \emph{ouvrage}) révèle, selon Freud, une angoisse ou une hostilité refoulée vis-à-vis du projet, que le discours officiel masque.
\end{exemplebox}

\courssub{Valeur et limites de la psychanalyse}

\begin{propriete}[Valeur explicative et anthropologique]
\begin{itemize}
\item \textbf{Intelligibilité du symptôme :} la psychanalyse donne un sens à ce qui paraissait absurde (phobies, TOC, névroses, actes manqués). Elle replace l'homme dans son histoire (infantile, familiale).
\item \textbf{Décenterement du sujet :} elle achève les « trois blessures narcissiques » (Copernic/Galilée : la Terre n'est pas le centre ; Darwin : l'homme n'est pas roi de la création ; Freud : le moi n'est pas maître dans sa propre maison).
\item \textbf{Valeur thérapeutique :} la « talking cure » (cure par la parole) vise à « rendre conscient l'inconscient » (\emph{Wo Es war, soll Ich werden} : « Là où c'était le Ça, doit advenir le Moi »), restaurant l'autonomie par l'insight (prise de conscience).
\end{itemize}
\end{propriete}

\begin{propriete}[Limites épistémologiques et éthiques]
\begin{itemize}
\item \textbf{Statut scientifique (Popper) :} Karl Popper (\emph{Conjectures et réfutations}) critique l'irréfutabilité de la psychanalyse : toute objection est interprétée comme « résistance », ce qui la met à l'abri de la falsification. Elle relève davantage de l'herméneutique (interprétation) que de la science expérimentale.
\item \textbf{Déterminisme psychique vs liberté :} si l'inconscient détermine tout, la responsabilité morale s'effondre. Sartre objecte que l'inconscient est une « excuse » pour fuir la liberté radicale.
\item \textbf{Sexualité infantile universalisée :} la théorie de la libido (stades oral, anal, phallique, génital) et du complexe d'Œdipe est critiquée comme ethnocentrique (Malinowski, Lévi-Strauss, méta-psychologie non universalisable).
\item \textbf{Risque de surinterprétation :} tout devient symbole sexuel/agressif ; l'analyste peut projeter sa propre grille de lecture (contre-transfert).
\end{itemize}
\end{propriete}

\begin{aretenir}[Synthèse : Conscience et Inconscient]
\begin{center}
\begin{tabularx}{\linewidth}{|l|X|X|}\hline
\textbf{Dimension} & \textbf{Conscience} & \textbf{Inconscient (Freud)} \\ \hline
Nature & Présence à soi/au monde ; intentionalité & Système dynamique de contenus refoulés \\ \hline
Accès & Immédiat (réflexif) & Indirect (rêves, lapsus, symptômes, transfert) \\ \hline
Logique & Identité, non-contradiction, temps & Intemporalité, condensation, déplacement, pas de « non » \\ \hline
Rôle & Connaissance, liberté, responsabilité & Détermination cachée, source des pulsions \\ \hline
Maîtrise & Illusion de transparence (Descartes) & « Le moi n'est pas maître dans sa propre maison » (Freud) \\ \hline
Tâche & Lucidité, examen de soi (Socrate, Kant, Alain) & Travail d'élucidation (psychanalyse, introspection critique) \\ \hline
\end{tabularx}
\end{center}
\end{aretenir}

\courssub{Savoir-faire 1 : Appliquer les notions de conscience et d'inconscient à une situation concrète}

\begin{methode}[Analyser une situation de vie avec les notions de conscience et d'inconscient]
Pour traiter une situation concrète (un fait de la vie courante, un récit, un exemple camerounais) à l'aide des notions de l'unité :
\begin{enumerate}
\item \textbf{Lire attentivement la situation} et repérer les faits pertinents : ce que le sujet fait, dit, ressent, oublie, avoue.
\item \textbf{Identifier la notion en jeu} : s'agit-il d'un acte volontaire et réfléchi (\cle{conscience}) ou d'un comportement dont le sujet ignore les motifs réels (\cle{inconscient}) ?
\item \textbf{Définir la notion mobilisée} en une phrase précise (la conscience comme présence à soi et au monde ; l'inconscient comme ensemble de contenus psychiques échappant à la conscience mais agissant sur elle).
\item \textbf{Appliquer la définition aux faits} : montrer point par point comment les éléments de la situation illustrent la notion (ex. : lapsus, rêve, oubli, acte manqué $\rightarrow$ manifestations de l'inconscient selon Freud).
\item \textbf{Nuancer} : rappeler la valeur et les limites de la notion (la conscience peut se tromper ; la psychanalyse est une interprétation, pas une science expérimentale au sens strict).
\item \textbf{Conclure} par une phrase de bilan qui répond explicitement à la question posée.
\end{enumerate}
\textbf{Réflexe} : toujours passer du général (définition) au particulier (la situation), puis du particulier au général (conclusion). Jamais de définition « parachutée » sans application.
\end{methode}

\begin{exoresolu}[Le cas de Manga : appliquer la notion d'inconscient]
\textbf{Énoncé.} Manga, élève en classe de Terminale technique à Douala, rêve fréquemment qu'il arrive en retard à l'examen du Baccalauréat sans avoir révisé. Le jour d'une épreuve blanche, il oublie sa convocation à la maison alors qu'il l'avait préparée la veille. Interrogé, il affirme : « Je ne suis pas du tout stressé par cet examen. » À l'aide de la notion d'inconscient, analysez cette situation.
\tcblower
\textbf{Solution détaillée.}
\begin{enumerate}
\item \textbf{Repérage des faits.} Trois faits sont significatifs : (a) un rêve récurrent d'échec à l'examen ; (b) l'oubli de la convocation ; (c) le discours conscient de Manga qui nie toute angoisse.
\item \textbf{Notion en jeu.} Il y a contradiction entre ce que Manga dit (conscience) et ce qu'il fait et rêve. C'est donc la notion d'\cle{inconscient} qui est pertinente.
\item \textbf{Définition.} Selon Freud, l'inconscient est la partie de la vie psychique constituée de désirs et de représentations refoulés, qui échappent à la conscience mais qui agissent sur le comportement et se manifestent de manière déguisée (rêves, lapsus, actes manqués, oublis).
\item \textbf{Application.}
\begin{itemize}
\item Le \textbf{rêve} est, selon Freud, la « voie royale » vers l'inconscient : le rêve de retard à l'examen exprime de manière déguisée une angoisse d'échec que Manga ne s'avoue pas.
\item L'\textbf{oubli} de la convocation est un \emph{acte manqué} : il n'est pas dû au hasard, mais traduit un désir inconscient (ici, peut-être le désir d'échapper à l'épreuve ou la peur de l'affronter).
\item Le \textbf{déni} (« je ne suis pas stressé ») illustre le \emph{refoulement} : la conscience écarte un contenu pénible, mais celui-ci continue d'agir en coulisses.
\end{itemize}
\item \textbf{Nuance.} La psychanalyse offre ici une interprétation cohérente, mais elle a ses limites : un oubli peut aussi avoir une cause purement matérielle (fatigue, distraction). On ne peut pas prouver expérimentalement l'interprétation freudienne ; elle reste une hypothèse explicative.
\item \textbf{Conclusion.} La situation de Manga montre que l'homme n'est pas totalement maître de lui-même : « le moi n'est pas maître dans sa propre maison » (Freud). Ses rêves et ses actes manqués révèlent une angoisse d'examen refoulée, que sa conscience refuse d'admettre. L'inconscient apparaît ainsi comme une réalité agissante que la conscience ne contrôle pas.
\end{enumerate}
\end{exoresolu}

\courssub{Savoir-faire 2 : Rédiger et justifier une réponse argumentée (question de cours)}

\begin{methode}[Construire une réponse argumentée à une question de cours]
Pour répondre à une question directe (ex. : « Quelle est la valeur de la conscience ? », « Quelles sont les limites de la psychanalyse ? ») :
\begin{enumerate}
\item \textbf{Dégager l'enjeu} de la question en une ou deux phrases (pourquoi la question se pose-t-elle ?).
\item \textbf{Définir les termes-clés} du sujet (conscience, inconscient, valeur, limite).
\item \textbf{Annoncer un plan} simple en deux ou trois parties (thèse / antithèse, ou aspects successifs).
\item \textbf{Développer chaque partie} selon le schéma : idée $\rightarrow$ argument $\rightarrow$ référence à un auteur (Descartes, Freud, Sartre, Nietzsche, Alain...) $\rightarrow$ exemple concret.
\item \textbf{Justifier chaque affirmation} : aucune thèse ne doit être posée sans raison ni illustration.
\item \textbf{Conclure} en répondant nettement à la question, sans introduire d'idée nouvelle.
\end{enumerate}
\textbf{Formules à retenir} : « Nous montrerons d'abord que..., ensuite que... » ; « Cette thèse trouve sa limite dans... » ; « Ainsi, ... ».
\end{methode}

\begin{exoresolu}[La conscience me suffit-elle à me connaître ?]
\textbf{Énoncé.} Question de réflexion : la conscience suffit-elle à la connaissance de soi ? Rédigez une réponse argumentée et justifiée (une vingtaine de lignes).
\tcblower
\textbf{Solution détaillée (corrigé rédigé).}

\emph{Introduction.} L'homme se croit transparent à lui-même : qui mieux que moi pourrait savoir ce que je pense et ce que je veux ? Pourtant, l'expérience montre que nous nous trompons souvent sur nos propres motifs. La conscience, définie comme la présence de l'esprit à lui-même et au monde, suffit-elle alors à la connaissance de soi ?

\emph{Première partie — La valeur de la conscience.} La conscience est d'abord le fondement de toute connaissance de soi. Descartes, avec le \emph{cogito} (« je pense, donc je suis »), montre que la conscience de soi est une évidence indubitable : même en doutant de tout, je ne peux douter que je suis un être qui pense. La conscience est aussi le siège de la liberté et de la responsabilité morale : selon Kant, c'est parce que je suis conscient de mes actes que je peux en répondre devant la loi morale. Au quotidien, un candidat au Baccalauréat technique qui réfléchit à ses forces et à ses faiblesses avant de choisir sa spécialité use de sa conscience réflexive pour se connaître.

\emph{Deuxième partie — Les limites de la conscience.} Cependant, la conscience a des limites. Elle peut être trompeuse : Nietzsche la compare à la surface d'une eau profonde, elle n'est que le sommet visible de la vie psychique. Freud révèle l'existence de l'\cle{inconscient} : des désirs refoulés, inaccessibles à la conscience, déterminent nos conduites (rêves, lapsus, actes manqués). De plus, la conscience immédiate se trompe : Alain remarque que celui qui croit agir par générosité agit parfois par vanité. Enfin, Bergson souligne que la conscience spontanée déforme le moi profond en le figeant dans des mots.

\emph{Conclusion.} La conscience est donc nécessaire mais non suffisante à la connaissance de soi. Elle en est le point de départ indispensable, mais elle doit être complétée par un travail critique sur soi, voire par l'exploration de l'inconscient. Se connaître exige de dépasser l'évidence immédiate du \emph{cogito}.
\end{exoresolu}

\courssub{Savoir-faire 3 : Construire une problématique et une introduction de dissertation}

\begin{methode}[Rédiger l'introduction d'une dissertation philosophique]
L'introduction se construit en quatre temps :
\begin{enumerate}
\item \textbf{L'accroche (amorce)} : partir d'un constat général, d'une expérience commune ou d'une opinion spontanée liée au sujet. Éviter les définitions de dictionnaire et les formules creuses (« depuis la nuit des temps... »).
\item \textbf{La mise en difficulté} : montrer que le sujet contient un problème (une tension, une contradiction, un doute). C'est le cœur de l'introduction.
\item \textbf{La problématique} : formuler le problème sous forme d'une question précise, qui reprend les termes du sujet sans les répéter mécaniquement.
\item \textbf{L'annonce du plan} : présenter brièvement les deux ou trois moments de la réflexion (thèse, antithèse, synthèse ou dépassement).
\end{enumerate}
\textbf{Réflexes} : l'introduction tient en un seul paragraphe de 8 à 12 lignes ; chaque étape doit enchaîner logiquement avec la précédente ; la problématique doit être une vraie question à laquelle le devoir répondra.
\end{methode}

\begin{exoresolu}[Sujet : « L'homme est-il maître de lui-même ? »]
\textbf{Énoncé.} Rédigez l'introduction complète d'une dissertation sur le sujet : « L'homme est-il maître de lui-même ? »
\tcblower
\textbf{Solution détaillée.}

\emph{Analyse préalable.} Termes-clés : « homme », « maître » (qui a le pouvoir de contrôler, de décider), « lui-même » (ses pensées, ses désirs, ses actes). Opinion spontanée : oui, car l'homme est conscient et libre. Difficulté : les passions, les erreurs, les comportements incompréhensibles suggèrent le contraire. Problème : la conscience suffit-elle à garantir la maîtrise de soi ?

\emph{Introduction rédigée.}

Il semble aller de soi que l'homme se gouverne : contrairement à l'animal conduit par l'instinct, il délibère, choisit et répond de ses actes. La conscience, cette lumière intérieure par laquelle chacun se saisit comme un « je », apparaît comme le fondement de cette maîtrise : être conscient, n'est-ce pas être présent à soi et donc se posséder ? Pourtant, l'expérience contredit cette évidence. Qui n'a jamais dit un mot qu'il regrettait aussitôt, cédé à une colère qu'il ne voulait pas, oublié précisément ce qu'il tenait à retenir ? Freud a donné à ces faits une signification radicale : derrière la conscience agirait un inconscient, fait de désirs refoulés, qui déciderait à notre insu. Le moi, écrit-il, « n'est pas maître dans sa propre maison ». Dès lors, la question se pose : \textbf{l'homme est-il véritablement maître de lui-même, ou sa conscience n'est-elle que le jouet de forces qui la dépassent ?} Pour répondre, nous examinerons d'abord la conscience comme pouvoir de maîtrise de soi ; nous montrerons ensuite que l'inconscient en révèle les limites ; enfin, nous demanderons si la maîtrise de soi, sans être totale, ne demeure pas une exigence et une tâche possible.
\end{exoresolu}

\courssub{Savoir-faire 4 : Rédiger un paragraphe argumenté (développement)}

\begin{methode}[Construire un paragraphe de développement]
Chaque paragraphe de dissertation suit une structure rigoureuse :
\begin{enumerate}
\item \textbf{Idée directrice} : une phrase qui énonce clairement la thèse du paragraphe.
\item \textbf{Explication} : développer l'idée, définir les notions, dérouler le raisonnement.
\item \textbf{Argument} : donner la raison qui fonde la thèse (argument logique ou référence à un auteur, cité brièvement et expliqué, jamais plaqué).
\item \textbf{Exemple} : illustrer par un cas concret, de préférence proche de la vie des élèves (contexte camerounais possible).
\item \textbf{Transition} : une phrase qui montre la limite de l'idée et prépare le paragraphe suivant.
\end{enumerate}
\textbf{Formule à retenir} : « Une idée, un argument, un exemple » par paragraphe. Une citation non expliquée ne vaut rien.
\end{methode}

\begin{exoresolu}[Paragraphe sur la valeur de la psychanalyse]
\textbf{Énoncé.} Rédigez un paragraphe de développement montrant la valeur de la psychanalyse freudienne pour la connaissance de l'homme.
\tcblower
\textbf{Solution détaillée.}

\emph{Plan du paragraphe.} Idée : la psychanalyse élargit la connaissance de l'homme au-delà de la conscience. Argument : elle rend intelligibles des conduites jusque-là inexplicables. Référence : Freud, \emph{La Psychopathologie de la vie quotidienne}. Exemple : le lapsus. Transition : mais cette interprétation est-elle vérifiable ?

\emph{Paragraphe rédigé.}

La psychanalyse constitue d'abord un progrès considérable dans la connaissance que l'homme a de lui-même, car elle étend le champ du psychique au-delà des limites de la conscience. Avant Freud, les conduites étranges — rêves absurdes, oublis inexplicables, paroles qui échappent — étaient traitées comme des accidents insignifiants. Freud, dans \emph{La Psychopathologie de la vie quotidienne}, soutient au contraire que rien dans la vie psychique n'est dû au hasard : le lapsus, l'acte manqué, le rêve sont des formations de l'inconscient, c'est-à-dire l'expression déguisée de désirs refoulés que la censure morale empêche d'accéder directement à la conscience. Ainsi, lorsqu'un élève de Ngaoundéré, invité à présenter publiquement son exposé, « oublie » le nom de son professeur au moment de le remercier, la psychanalyse n'y voit pas une simple défaillance de mémoire, mais l'aveu indirect d'une hostilité ou d'une crainte refoulée. La valeur de la psychanalyse est donc de donner un sens là où la conscience ne voyait que chaos, et de révéler que l'homme est plus vaste que ce qu'il sait de lui-même. Toutefois, une interprétation qui prétend tout expliquer ne risque-t-elle pas de n'être jamais vérifiable ? C'est la question des limites de la psychanalyse qu'il faut maintenant examiner.
\end{exoresolu}

\courssub{Savoir-faire 5 : Commenter un texte philosophique}

\begin{methode}[Méthode du commentaire de texte]
Pour commenter un texte d'auteur (Descartes, Freud, Sartre...) :
\begin{enumerate}
\item \textbf{Lecture analytique} : identifier l'auteur, l'œuvre, le thème ; repérer la thèse défendue et les articulations du texte (mots charnières : « mais », « donc », « cependant »).
\item \textbf{Dégager la thèse} : formuler en une phrase ce que l'auteur soutient.
\item \textbf{Relever la structure} : diviser le texte en deux ou trois moments et justifier chaque coupure.
\item \textbf{Expliquer le texte pas à pas} : suivre l'ordre du texte, expliquer les notions, expliciter les arguments, illustrer par des exemples personnels. Ne jamais paraphraser : expliquer, c'est dire \emph{pourquoi} l'auteur avance chaque idée.
\item \textbf{Discussion critique} : évaluer la thèse (sa force, ses limites, une objection éventuelle), en mobilisant d'autres auteurs.
\item \textbf{Conclusion} : bilan de l'intérêt philosophique du texte et ouverture raisonnée.
\end{enumerate}
\textbf{Réflexe} : toujours citer une courte portion du texte avant de l'expliquer, pour ancrer le commentaire.
\end{methode}

\begin{exoresolu}[Commentaire guidé d'un extrait de Freud]
\textbf{Énoncé.} Texte : « Le moi n'est pas maître dans sa propre maison. [...] Le moi se comporte comme s'il était le maître, mais il est en réalité déterminé par des processus psychiques qui lui sont inconnus. » (Freud, adapté). Expliquez et discutez ce texte.
\tcblower
\textbf{Solution détaillée.}

\emph{1. Présentation et thèse.} Cet extrait de Freud résume la thèse centrale de la psychanalyse : la conscience (le « moi ») ne se gouverne pas elle-même ; elle est déterminée par l'inconscient.

\emph{2. Structure.} Le texte comporte deux moments : (a) l'énoncé de la thèse par une métaphore (« maître dans sa propre maison ») ; (b) son explication (« déterminé par des processus psychiques qui lui sont inconnus »).

\emph{3. Explication.} La métaphore de la maison signifie que le moi se croit souverain dans son propre psychisme, alors qu'il n'en est que le locataire partiel : d'autres « habitants » — les pulsions et les désirs refoulés — agissent sans son autorisation. L'expression « processus psychiques qui lui sont inconnus » définit l'inconscient : non pas une absence de psychisme, mais un psychisme réel dont la conscience n'a pas connaissance. La formule « se comporte comme s'il était le maître » souligne le caractère illusoire de la maîtrise consciente : notre sentiment de liberté et de transparence à soi serait une illusion nécessaire. Ainsi, le locuteur qui commet un lapsus croit avoir simplement « mal parlé » ; en réalité, selon Freud, c'est un désir refoulé qui s'est exprimé à sa place.

\emph{4. Discussion.} La force de cette thèse est de rendre compte de faits que la psychologie de la conscience laissait inexpliqués (rêves, symptômes, actes manqués). Mais elle connaît des limites : Karl Popper objecte que la psychanalyse n'est pas réfutable — toute objection peut être interprétée comme une « résistance », ce qui la met à l'abri de toute vérification, donc hors du champ de la science expérimentale. Sartre, de son côté, refuse l'inconscient au nom de la liberté : pour lui, l'homme est « condamné à être libre » et se mentir à soi-même (la « mauvaise foi ») n'exige pas un inconscient, car celui qui se trompe doit, à un certain niveau, savoir ce qu'il cache.

\emph{5. Conclusion.} Ce texte a la valeur de briser l'illusion d'une conscience transparente et souveraine ; mais réduire l'homme à un jouet de l'inconscient menacerait la responsabilité morale et la liberté. La vérité semble intermédiaire : la conscience n'est pas toute-puissante, mais elle peut, par le travail de réflexion, élargir son empire.
\end{exoresolu}

\courssub{Savoir-faire 6 : Rédiger une conclusion et justifier une position personnelle}

\begin{methode}[Conclure un devoir en justifiant sa réponse]
La conclusion n'est pas un résumé mécanique :
\begin{enumerate}
\item \textbf{Rappeler le chemin parcouru} en une phrase (les grandes étapes du devoir).
\item \textbf{Répondre nettement à la problématique} : oui, non, ou réponse nuancée — mais toujours tranchée et assumée.
\item \textbf{Justifier la réponse} : donner la raison principale qui fonde votre position (reprendre l'argument décisif du devoir).
\item \textbf{Ouvrir} éventuellement sur une question nouvelle liée au sujet, sans la traiter.
\end{enumerate}
\textbf{Piège à éviter} : ne jamais introduire dans la conclusion un argument ou un auteur non traité dans le développement.
\end{methode}

\begin{exoresolu}[Conclusion du sujet « L'homme est-il maître de lui-même ? »]
\textbf{Énoncé.} Rédigez la conclusion de la dissertation dont la problématique était : « L'homme est-il véritablement maître de lui-même, ou sa conscience n'est-elle que le jouet de forces qui la dépassent ? » (développement : I. la conscience comme pouvoir de maîtrise ; II. l'inconscient comme limite ; III. la maîtrise de soi comme tâche).
\tcblower
\textbf{Solution détaillée.}

\emph{Conclusion rédigée.}

Notre réflexion nous a conduits de l'évidence du \emph{cogito} cartésien, qui fait de la conscience le fondement de la maîtrise de soi, à la découverte freudienne de l'inconscient, qui révèle que le moi « n'est pas maître dans sa propre maison ». Fallait-il en conclure que l'homme n'est que le jouet de forces obscures ? Nous avons soutenu que non. Si la conscience n'est ni transparente ni toute-puissante, elle conserve le pouvoir de prendre acte de ses propres limites : reconnaître que je puis être déterminé, c'est déjà commencer à ne plus l'être aveuglément. La maîtrise de soi n'est donc pas un donné, mais une \textbf{tâche} : celle d'une lucidité toujours à conquérir, par la réflexion, l'éducation et l'examen critique de soi. L'homme n'est pas maître de lui-même par nature ; il peut le devenir par effort. Reste alors une question : cette conquête de soi a-t-elle des limites que la volonté ne pourra jamais franchir ?
\end{exoresolu}

\begin{attention}[Les 5 erreurs qui coûtent des points au Bac]
\begin{enumerate}
\item \textbf{Confondre inconscient et non-conscient.} L'inconscient freudien n'est pas le simple fonctionnement automatique du corps (respiration, digestion) ni l'absence momentanée d'attention : c'est un psychisme actif, fait de désirs \emph{refoulés}, qui agit sur la conscience. Ne pas écrire « inconscient » pour « inattentif » ou « endormi ».
\item \textbf{Plaquer les citations sans les expliquer.} Écrire « Freud dit que le moi n'est pas maître dans sa propre maison » puis passer à autre chose ne rapporte rien : toute référence doit être expliquée (que signifie la métaphore ? quels faits la justifient ?) et reliée à l'argumentation.
\item \textbf{Présenter la psychanalyse comme une science expérimentale exacte.} Erreur symétrique : la rejeter totalement. La bonne position est nuancée : la psychanalyse a une valeur explicative et thérapeutique, mais ses interprétations sont difficilement vérifiables (objection de Popper : l'irréfutabilité). Oublier cette nuance fait perdre les points de la discussion critique.
\item \textbf{Traiter la conscience comme infaillible dans la première partie sans annoncer ses limites.} Le devoir doit montrer la valeur \emph{et} les limites de la conscience (elle peut se tromper : illusion, mauvaise foi, erreur sur ses propres motifs). Un devoir à sens unique, sans tension, ne dépasse pas la moyenne.
\item \textbf{Négliger la problématique et la conclusion.} Une introduction sans question clairement formulée, ou une conclusion qui ne répond pas à cette question (ou qui introduit une idée neuve), coûte plusieurs points. Vérifier avant de rendre : ma problématique est-elle une question ? Ma conclusion y répond-elle explicitement ?
\end{enumerate}
\end{attention}

\newpage

\coursec{Exercices}

\courssub{Je vérifie mes connaissances}

\begin{exercice}[QCM : Notions fondamentales]
\emph{Choisissez la seule bonne réponse pour chaque question.}
\begin{enumerate}
\item Pour Descartes, la conscience est :
\begin{itemize}
\item[A.] Une propriété du corps qui permet de percevoir le monde extérieur.
\item[B.] La pensée même, c'est-à-dire la connaissance immédiate et certaine que l'esprit a de ses propres opérations (\emph{cogito}).
\item[C.] Un effet de l'inconscient qui nous trompe sur nos véritables motivations.
\item[D.] Une faculté passive qui subit les impressions des sens.
\end{itemize}
\item Le \emph{refoulement} (Verdrängung) chez Freud désigne :
\begin{itemize}
\item[A.] L'oubli naturel des souvenirs anciens par manque d'attention.
\item[B.] Le processus psychique par lequel des représentations incompatibles avec le Moi sont repoussées dans l'Inconscient.
\item[C.] La censure exercée par le Surmoi sur les désirs du Ça pendant le sommeil.
\item[D.] La transformation de l'énergie psychique en symptôme corporel (conversion).
\end{itemize}
\item Spinoza, dans l'\emph{Éthique}, affirme que « les hommes se croient libres parce qu'ils sont conscients de leurs actions et ignorants des causes par lesquelles ils sont déterminés ». Cela signifie que :
\begin{itemize}
\item[A.] La conscience est la preuve de notre libre arbitre absolu.
\item[B.] La conscience est une connaissance adéquate des causes de nos actions.
\item[C.] La conscience est une connaissance \emph{inadéquate} (partielle, confuse) qui masque la nécessité naturelle.
\item[D.] L'inconscient n'existe pas, il n'y a que de l'ignorance temporaire.
\end{itemize}
\item La topique freudienne (premier modèle) distingue trois instances : le Ça, le Moi, le Surmoi. Laquelle est le siège des pulsions, irrationnelle et régie par le principe de plaisir ?
\begin{itemize}
\item[A.] Le Moi (\emph{Ich}).
\item[B.] Le Surmoi (\emph{Über-Ich}).
\item[C.] Le Ça (\emph{Es}).
\item[D.] Le Préconscient.
\end{itemize}
\item Karl Popper critique la psychanalyse au nom du critère de \emph{démarcation} scientifique. Il lui reproche principalement :
\begin{itemize}
\item[A.] D'être trop difficile à comprendre pour le grand public.
\item[B.] De ne pas être \emph{réfutable} (non falsifiable) : elle explique tout a posteriori et ne prédit rien de risqué.
\item[C.] De ne pas soigner les malades mentaux aussi bien que la chimie.
\item[D.] D'accorder trop d'importance à la sexualité infantile.
\end{itemize}
\end{enumerate}
\end{exercice}

\begin{exercice}[Vrai ou Faux ? Justifiez brièvement]
\emph{Indiquez VRAI ou FAUX et justifiez en une phrase en mobilisant une notion du cours.}
\begin{enumerate}
\item La conscience morale (conscience de devoir) est identique à la conscience psychologique (conscience de soi).
\item L'inconscient freudien est un « deuxième moi » rationnel qui agit en secret.
\item Le « lapsus » (acte manqué) est, pour Freud, une preuve de l'existence de l'inconscient et du retour du refoulé.
\item La psychanalyse est une science au même titre que la physique newtonienne car elle utilise l'observation clinique.
\end{enumerate}
\end{exercice}

\begin{exercice}[Définitions et repères d'auteurs]
\emph{Donnez une définition précise (2-3 lignes) pour chacun des concepts suivants et citez l'auteur de référence principal.}
\begin{enumerate}
\item \cle{Conscience} (au sens cartésien / phénoménologique).
\item \cle{Inconscient} (au sens freudien / topique).
\item \cle{Refoulement}.
\item \cle{Ça / Moi / Surmoi} (définir les trois instances et leur principe de fonctionnement respectif).
\item \cle{Connaissance inadéquate} (Spinoza).
\item \cle{Critère de falsifiabilité} (Popper).
\end{enumerate}
\end{exercice}

\begin{exercice}[Mise en relation : Concepts et Auteurs]
\emph{Reliez chaque thèse à son auteur et à la notion clé correspondante.}
\begin{center}
\begin{tabularx}{\linewidth}{|l|X|X|}\hline
\textbf{Thèse / Citation} & \textbf{Auteur} & \textbf{Notion clé} \\ \hline
« Je pense, donc je suis » ; la conscience est transparente à elle-même. & & \\ \hline
L'homme est un « dieu déchu » ; l'inconscient gouverne la vie psychique. & & \\ \hline
La liberté n'est pas le libre arbitre mais la connaissance de la nécessité (connaissance adéquate). & & \\ \hline
Une théorie qui s'arrange avec tous les faits n'en explique aucun ; la psychanalyse est une « pseudo-science ». & & \\ \hline
\end{tabularx}
\end{center}
\emph{Auteurs à placer : Descartes, Freud, Spinoza, Popper. Notions : Cogito / Certitude, Inconscient / Pulsions, Connaissance adéquate / Nécessité, Falsifiabilité / Démarcation.}
\end{exercice}

\courssub{Je m'entraîne}

\begin{exercice}[Analyse de cas : L'acte manqué au marché de Douala]
M. Kamga, commerçant à Douala, doit régler une dette de 500 000 FCFA à M. Njoya, un fournisseur avec qui il a eu un violent différend il y a six mois. Le jour de l'échéance, M. Kamga se rend à la banque, fait la queue, remplit le bordereau de virement, mais écrit par erreur le numéro de compte de son \emph{frère cadet} au lieu de celui de M. Njoya. Il s'en aperçoit deux jours plus tard, après avoir reçu un appel furieux de son frère. M. Kamga affirme : « Je ne comprends pas, je voulais vraiment le payer, c'est une erreur de distraction. »
\begin{enumerate}
\item Identifiez le type de phénomène psychique décrit (terme technique freudien).
\item Expliquez ce comportement en utilisant les notions de \cle{refoulement}, de \cle{Ça} (pulsion) et de \cle{Moi} (réalité).
\item Ce « lapsus » (acte manqué) engage-t-il la responsabilité morale de M. Kamga ? Justifiez en distinguant \cle{responsabilité pénale} et \cle{responsabilité psychique}.
\end{enumerate}
\end{exercice}

\begin{exercice}[Confrontation : Descartes vs Spinoza sur la transparence de la conscience]
\begin{enumerate}
\item Résumez en 3 lignes la thèse de Descartes sur la conscience comme \emph{cogito} (première certitude, transparence, substantialité de l'âme).
\item Résumez en 3 lignes la critique de Spinoza : la conscience comme « idée idée » mais connaissance \emph{inadéquate} des causes réelles (Déterminisme, Conatus).
\item \textbf{Problématisez} : La conscience est-elle un « miroir fidèle » de ce que nous sommes ? Répondez en 10 lignes en opposant les deux conceptions.
\end{enumerate}
\end{exercice}

\begin{exercice}[Le travail du rêve : Mécanismes et interprétation]
Un patient raconte à son analyste le rêve suivant : « Je monte un grand escalier en colimaçon dans une tour. Soudain, l'escalier s'effondre et je tombe dans une pièce remplie de lettres non ouvertes. Je me réveille en sueur. » L'analyste note que le patient a une relation conflictuelle avec son père (autorité, lettres administratives) et une angoisse de performance professionnelle (gravir les échelons).
\begin{enumerate}
\item Identifiez le \cle{contenu manifeste} et le \cle{contenu latent} (hypothétique) de ce rêve.
\item Nommez et définissez deux \cle{mécanismes du travail du rêve} (parmi : condensation, déplacement, figuration, élaboration secondaire) qui pourraient expliquer le passage du latent au manifeste dans cet exemple.
\item Pourquoi Freud dit-il que le rêve est la « voie royale vers l'inconscient » ?
\end{enumerate}
\end{exercice}

\begin{exercice}[Argumentation dirigée : « L'inconscient nous délivre-t-il de la responsabilité ? »]
\emph{Rédigez une réponse structurée (15-20 lignes) en suivant ce plan imposé :}
\begin{enumerate}
\item \textbf{Thèse} : En quoi l'inconscient (déterminisme pulsionnel, refoulement) semble abolir la responsabilité consciente (l'homme n'est pas maître en sa propre maison).
\item \textbf{Antithèse} : En quoi la responsabilité juridique et morale ne peut pas être dissoute par l'inconscient (notion de \cle{Moi} comme médiateur, devoir de \cle{travailler} son inconscient, responsabilité du sujet divisé).
\item \textbf{Synthèse} : Redéfinir la responsabilité non plus comme conscience immédiate de l'acte, mais comme \cle{prise en charge} de son histoire psychique (référence à Ricoeur ou à l'éthique psychanalytique).
\end{enumerate}
\end{exercice}

\courssub{Je me prépare au Bac}

\begin{exercice}[Sujet de Dissertation type Baccalauréat Technique]
\textbf{Sujet :} « La conscience suffit-elle à définir l'homme ? »
\emph{Consignes : Vous traiterez ce sujet sous la forme d'une dissertation philosophique structurée (Introduction, Développement en trois parties dialectiques : Thèse / Antithèse / Synthèse, Conclusion). Vous mobiliserez obligatoirement les références au programme : Descartes, Spinoza, Freud. Vous illustrerez vos arguments par des exemples concrets (dont au moins un exemple camerounais).}
\begin{enumerate}
\item Rédigez l'\textbf{Introduction} complète (Amorce, Définitions des termes, Problématique, Annonce du plan).
\item Présentez le \textbf{Plan détaillé} (Arguments principaux, Auteurs, Exemples pour chaque sous-partie).
\item Rédigez la \textbf{Conclusion} (Réponse synthétique, Ouverture).
\end{enumerate}
\end{exercice}

\begin{exercice}[Explication de texte guidée : Freud, \emph{L'Interprétation des rêves} (1900)]
\emph{Texte :} « Le rêve est la réalisation déguisée d'un désir refoulé. [...] Le travail du rêve ne consiste pas à créer le désir, mais à lui donner une forme qui permette de tromper la censure. Le contenu latent (les pensées inconscientes du rêve) est transformé en contenu manifeste (le rêve tel qu'on le raconte) par des opérations comme le déplacement et la condensation. L'interprétation vise à remonter du manifeste au latent. »
\begin{enumerate}
\item \textbf{Dégagez la thèse principale} du texte en une phrase.
\item \textbf{Expliquez} les concepts suivants à l'aide du texte et de vos connaissances : \cle{réalisation déguisée}, \cle{désir refoulé}, \cle{censure}, \cle{contenu latent / contenu manifeste}, \cle{déplacement}, \cle{condensation}.
\item \textbf{Structurez le raisonnement} : Quelles sont les étapes logiques qui mènent de la définition du rêve à la méthode d'interprétation ?
\item \textbf{Discutez la portée} : En quoi cette conception du rêve bouleverse-t-elle l'idéal cartésien de la maîtrise de soi par la raison ? La psychanalyse est-elle une « science » au sens épistémologique strict (critère de Popper) ?
\end{enumerate}
\end{exercice}

\begin{exercice}[Épreuve type Probatoire : « Peut-on se connaître soi-même ? »]
\emph{Ce sujet est tombé au Probatoire Technique session 2023. Vous devez produire une copie complète d'examen (durée 3h). Pour cet exercice, on vous demande de réaliser les deux premières étapes fondamentales :}
\begin{enumerate}
\item \textbf{Rédigez l'Introduction complète} (minimum 25 lignes) : Accroche contextuelle (ex: inscription delphique « Connais-toi toi-même »), Définitions rigoureuses de \cle{connaissance}, \cle{soi}, \cle{conscience}, \cle{inconscient}, \textbf{Problématique} (tension entre transparence cartésienne et opacité freudienne/spinoziste), \textbf{Annonce de plan dialectique} (3 parties numérotées avec titres explicites).
\item \textbf{Proposez un Plan détaillé} (avec arguments, auteurs, concepts, exemples pour chaque sous-partie) :
\begin{itemize}
\item \textbf{I. Thèse : La conscience comme accès premier et certain à soi} (Descartes, Phénoménologie).
\item \textbf{II. Antithèse : L'opacité radicale du sujet : l'inconscient et la détermination cachée} (Freud, Spinoza, Sciences cognitives).
\item \textbf{III. Synthèse : La connaissance de soi comme tâche infinie, éthique et narrative} (Ricoeur, Sartre, « travail » sur soi).
\end{itemize}
\item \textbf{Question bonus (5 lignes) :} En quoi le contexte camerounais (poids de la tradition orale, rapports aux ancêtres, médecine traditionnelle vs psychanalyse) éclaire-t-il différemment la question de la connaissance de soi ?
\end{enumerate}
\end{exercice}

\newpage

\coursec{Corrigés des exercices}

\coursec{Leçon 18 : La Conscience et l'Inconscient}

\courssub{Entraînement aux notions fondamentales}

\begin{corrige}[Exercice 1 — QCM : Notions fondamentales]
\begin{enumerate}
\item \textbf{Réponse B.} Pour Descartes, la conscience est la pensée elle-même : « Par le terme de penser, j'entends tout ce qui se fait en nous de telle sorte que nous l'apercevons immédiatement par nous-mêmes » (\emph{Principes de la Philosophie}, I, art. 9). Le \emph{cogito} (« je pense, donc je suis ») est la première certitude indubitable. A est faux (la conscience est attribut de l'âme, non du corps), C est la thèse freudienne (postérieure), D contredit l'activité de la pensée cartésienne.
\item \textbf{Réponse B.} Le refoulement est le mécanisme de défense par lequel le Moi repousse hors de la conscience des représentations ou désirs incompatibles avec ses exigences morales. A est faux (le refoulé n'est pas oublié, il reste actif), C est faux (la censure s'exerce pendant le sommeil mais n'est pas le refoulement), D décrit la conversion hystérique, autre mécanisme.
\item \textbf{Réponse C.} Pour Spinoza, la conscience nous donne l'effet (nos actions, nos désirs) mais pas les causes qui nous déterminent : c'est une connaissance inadéquate, source de l'illusion du libre arbitre. A est précisément ce que Spinoza réfute ; B est l'inverse de sa thèse ; D est faux car Spinoza ne nie pas les causes cachées, il affirme la nécessité naturelle.
\item \textbf{Réponse C.} Le Ça (\emph{Es}) est le réservoir des pulsions, totalement inconscient, irrationnel, régi par le principe de plaisir. Le Moi est régi par le principe de réalité, le Surmoi est l'instance morale (interdits, idéal), le Préconscient n'est pas une instance pulsionnelle.
\item \textbf{Réponse B.} Popper reproche à la psychanalyse son irréfutabilité : aucun comportement observable ne peut la contredire, car elle interprète tout a posteriori (le refoulement explique le souvenir comme l'oubli). Une théorie compatible avec tout ne prédit rien et n'est donc pas scientifique au sens du critère de démarcation.
\end{enumerate}
\end{corrige}

\begin{attention}[Points de vigilance — Exercice 1]
Ne pas confondre \emph{refoulement} (mise hors de la conscience) et \emph{censure} (barrière entre inconscient et préconscient, relâchée dans le rêve). Ne pas attribuer à Descartes la thèse de l'inconscient : c'est l'erreur la plus fréquente au Bac.
\end{attention}

\begin{corrige}[Exercice 2 — Vrai ou Faux ? Justifiez brièvement]
\begin{enumerate}
\item \textbf{FAUX.} La conscience psychologique est la présence de l'esprit à lui-même (conscience de soi, perception de ses états), tandis que la conscience morale est le jugement porté sur le bien et le mal de nos actes (conscience du devoir) : ce sont deux sens distincts du mot « conscience ».
\item \textbf{FAUX.} L'inconscient freudien n'est ni un moi ni rationnel : il est un ensemble de représentations et de pulsions refoulées, irrationnelles, ignorantes de la contradiction et du temps, régies par le principe de plaisir (processus primaires).
\item \textbf{VRAI.} Pour Freud (\emph{Psychopathologie de la vie quotidienne}), le lapsus n'est pas un hasard : il est un compromis entre l'intention consciente et un désir refoulé qui s'exprime malgré la censure — c'est un « retour du refoulé » et donc une manifestation de l'inconscient.
\item \textbf{FAUX.} Selon le critère de Popper, une théorie scientifique doit être falsifiable (réfutable par l'expérience) ; or la psychanalyse, expliquant tous les cas possibles, n'est pas réfutable : l'observation clinique seule ne suffit pas à faire une science au sens strict.
\end{enumerate}
\end{corrige}

\begin{corrige}[Exercice 3 — Définitions et repères d'auteurs]
\begin{enumerate}
\item \textbf{Conscience} (Descartes) : connaissance immédiate et certaine que l'esprit a de lui-même et de ses propres opérations (penser, douter, vouloir, sentir). Elle est transparente à elle-même et fonde la première certitude : le \emph{cogito}. Auteur de référence : \textbf{Descartes}.
\item \textbf{Inconscient} (Freud) : partie de l'appareil psychique constituée de contenus (désirs, souvenirs, représentations) maintenus hors de la conscience par le refoulement, mais qui demeurent actifs et se manifestent de façon déguisée (rêves, lapsus, symptômes). Auteur : \textbf{Freud}.
\item \textbf{Refoulement} : mécanisme de défense par lequel le Moi repousse dans l'inconscient les représentations pulsionnelles incompatibles avec les exigences morales et sociales. Le refoulé n'est pas détruit : il cherche à revenir sous forme déguisée. Auteur : \textbf{Freud}.
\item \textbf{Ça / Moi / Surmoi} (Freud, seconde topique) : le \emph{Ça} est le pôle pulsionnel inconscient, régi par le principe de plaisir ; le \emph{Moi} est l'instance de médiation avec le monde extérieur, régie par le principe de réalité ; le \emph{Surmoi} est l'instance morale (interdits parentaux intériorisés, idéal du Moi), juge et censeur du Moi.
\item \textbf{Connaissance inadéquate} (Spinoza) : connaissance partielle et confuse qui saisit les effets sans connaître les causes ; c'est le premier genre de connaissance (imagination, opinion), source des passions et de l'illusion du libre arbitre. Auteur : \textbf{Spinoza} (\emph{Éthique}).
\item \textbf{Critère de falsifiabilité} (Popper) : une théorie n'est scientifique que si elle formule des prédictions risquées susceptibles d'être réfutées par l'expérience ; une théorie compatible avec tous les faits possibles est une pseudo-science. Auteur : \textbf{Popper} (\emph{Conjectures et réfutations}).
\end{enumerate}
\end{corrige}

\begin{corrige}[Exercice 4 — Mise en relation : Concepts et Auteurs]
\begin{center}
\begin{tabularx}{\linewidth}{|X|l|X|}\hline
\rowcolor{popblueL}\textbf{Thèse / Citation} & \textbf{Auteur} & \textbf{Notion clé} \\ \hline
« Je pense, donc je suis » ; la conscience est transparente à elle-même. & \textbf{Descartes} & Cogito / Certitude \\ \hline
L'homme est un « dieu déchu » ; l'inconscient gouverne la vie psychique. & \textbf{Freud} & Inconscient / Pulsions \\ \hline
La liberté n'est pas le libre arbitre mais la connaissance de la nécessité. & \textbf{Spinoza} & Connaissance adéquate / Nécessité \\ \hline
Une théorie qui s'arrange avec tous les faits n'en explique aucun. & \textbf{Popper} & Falsifiabilité / Démarcation \\ \hline
\end{tabularx}
\end{center}
\emph{Justifications :} la première phrase reprend le \emph{cogito} des \emph{Méditations} ; la seconde renvoie à la « troisième blessure narcissique » infligée à l'homme par la psychanalyse (Freud) ; la troisième résume l'\emph{Éthique} (livres IV-V) ; la quatrième formule le critère de démarcation poppérien.
\end{corrige}

\courssub{Application et analyse : cas concrets et mécanismes}

\begin{corrige}[Exercice 5 — Analyse de cas : L'acte manqué au marché de Douala]
\begin{enumerate}
\item Il s'agit d'un \cle{acte manqué} (parapraxie), plus précisément d'une erreur d'écriture comparable au lapsus : l'accomplissement d'une action contraire à l'intention consciente déclarée.
\item \textbf{Analyse freudienne.} M. Kamga éprouve vraisemblablement de l'hostilité ou de la rancune envers M. Njoya depuis leur différend : ce désir agressif (« ne pas le payer », « le punir ») est incompatible avec l'image sociale et morale du commerçant honnête ; il est donc \emph{refoulé} dans l'inconscient. Le \cle{Ça} (pulsion agressive, principe de plaisir) cherche satisfaction ; le \cle{Moi} (principe de réalité) sait que la dette doit être payée et organise le virement. L'acte manqué est une \emph{formation de compromis} : le Moi accomplit formellement le paiement, mais le désir refoulé sabote l'opération en substituant le compte du frère. Le « retour du refoulé » trahit le conflit psychique.
\item \textbf{Responsabilité.} Sur le plan \emph{pénal}, M. Kamga n'a pas commis de faute intentionnelle : il n'y a ni dol ni négligence caractérisée, l'erreur étant involontaire ; sa responsabilité juridique n'est donc pas engagée (il devra simplement régulariser le virement). Sur le plan \emph{psychique} en revanche, l'acte manqué révèle un désir qui lui appartient : pour Freud, le sujet est responsable de ses rêves et de ses actes manqués parce qu'ils expriment \emph{sa} vérité inconsciente. La responsabilité n'est donc pas abolie mais déplacée : M. Kamga est invité à reconnaître et à assumer son hostilité refoulée plutôt qu'à la nier.
\end{enumerate}
\end{corrige}

\begin{attention}[Points de vigilance — Exercice 5]
Bien distinguer les deux plans : le droit juge l'intention consciente et l'acte ; la psychanalyse interroge la vérité du sujet. Ne pas conclure que « l'inconscient excuse tout » : c'est l'écueil classique.
\end{attention}

\begin{corrige}[Exercice 6 — Confrontation : Descartes vs Spinoza]
\begin{enumerate}
\item \textbf{Descartes.} La conscience est la première certitude : dans le doute hyperbolique, je ne peux douter que je pense ; le \emph{cogito} (« je pense, donc je suis ») m'atteint comme substance pensante. L'âme est transparente à elle-même : rien n'est en moi dont je n'aie conscience ; la pensée est l'attribut essentiel de la substance spirituelle, connue plus clairement que le corps.
\item \textbf{Spinoza.} La conscience est bien une « idée de l'idée », mais elle ne saisit que les effets (affections du corps, désirs) sans en connaître les causes : c'est une connaissance \emph{inadéquate}. Tout dans la nature est déterminé par la nécessité (déterminisme) ; chaque être tend à persévérer dans son être (\emph{conatus}). L'homme se croit libre parce qu'il a conscience de ses désirs et ignore les causes qui le déterminent — comme la pierre lancée qui se croirait libre de son mouvement.
\item \textbf{Problématisation (10 lignes).} La conscience est-elle un miroir fidèle de ce que nous sommes ? Pour Descartes, oui : elle est le lieu d'une présence totale à soi, fondement de toute certitude et de la liberté ; je suis exactement ce dont j'ai conscience. Pour Spinoza, non : la conscience n'est que la surface éclairée d'un vaste déterminisme ; elle me donne mes désirs mais me cache leurs causes, et l'illusion du libre arbitre naît précisément de cette ignorance. Le miroir est donc déformant : il reflète les effets et masque les causes. Toutefois, Spinoza ne détruit pas la conscience : il lui assigne une tâche, passer de la connaissance inadéquate à la connaissance adéquate par la raison. La conscience n'est plus alors un miroir passif mais un travail : se connaître, c'est remonter aux causes qui nous font agir. La fidélité de la conscience n'est pas un donné, c'est une conquête.
\end{enumerate}
\end{corrige}

\begin{corrige}[Exercice 7 — Le travail du rêve : Mécanismes et interprétation]
\begin{enumerate}
\item \textbf{Contenu manifeste} : le récit tel qu'il est raconté — l'escalier en colimaçon, l'effondrement, la chute, les lettres non ouvertes, le réveil angoissé. \textbf{Contenu latent (hypothétique)} : les pensées et désirs inconscients que le rêve réalise de façon déguisée — ici, le conflit avec la figure paternelle et autoritaire (les lettres administratives), le désir de réussite sociale (« gravir les échelons ») mêlé à la crainte de l'échec et de la sanction paternelle (l'effondrement de l'escalier).
\item \textbf{Deux mécanismes du travail du rêve :}
\begin{itemize}
\item La \cle{condensation} : plusieurs éléments latents se fondent en une seule image manifeste. L'escalier condense à la fois l'ambition professionnelle (monter les échelons) et la hiérarchie paternelle (s'élever sous le regard du père).
\item Le \cle{déplacement} : l'intensité affective est transférée d'un élément important vers un élément secondaire. L'angoisse liée au père est déplacée sur les lettres non ouvertes, objets apparemment anodins, ce qui permet de tromper la censure.
\end{itemize}
(On pouvait aussi citer la \emph{figuration} — transformation de pensées abstraites en images — et l'\emph{élaboration secondaire} — mise en récit cohérent au réveil.)
\item Le rêve est la « voie royale vers l'inconscient » parce que, pendant le sommeil, la censure qui barre l'accès à la conscience se relâche : les désirs refoulés peuvent alors s'exprimer, certes déguisés par le travail du rêve, mais de manière repérable. L'interprétation, en remontant du manifeste au latent, donne ainsi accès aux contenus inconscients les mieux gardés du sujet.
\end{enumerate}
\end{corrige}

\courssub{Argumentation, dissertation et explication de texte}

\begin{corrige}[Exercice 8 — Argumentation dirigée : « L'inconscient nous délivre-t-il de la responsabilité ? »]
\textbf{Thèse.} La découverte freudienne semble ruiner la responsabilité. Si le psychisme est gouverné par des pulsions inconscientes, des conflits infantiles refoulés et un déterminisme psychique strict, alors le sujet conscient n'est plus l'auteur de ses actes : « le Moi n'est pas maître en sa propre maison » (Freud). Le criminel pouvait-il ne pas agir, si son geste accomplissait un scénario inconscient ? L'acte manqué de M. Kamga montre que nos conduites nous échappent. Comment être responsable de ce dont on n'a pas conscience ? La conscience morale elle-même paraît alors une illusion, comparable à l'illusion du libre arbitre dénoncée par Spinoza.

\textbf{Antithèse.} Pourtant, dissoudre la responsabilité dans l'inconscient serait contradictoire et dangereux. D'abord, le droit ne peut juger que des sujets présumés capables de discernement : sans responsabilité, plus de société possible. Ensuite, Freud lui-même ne déresponsabilise pas : le \cle{Moi} est précisément l'instance de médiation entre le Ça, le Surmoi et la réalité ; il a pour tâche de négocier avec les pulsions, non de les subir. Enfin, le contenu inconscient est \emph{le mien} : mes rêves, mes lapsus, mes symptômes expriment mon histoire. Dire « ce n'est pas moi, c'est mon inconscient » serait une mauvaise foi : le sujet est divisé, mais non dédoublé en deux personnes.

\textbf{Synthèse.} Il faut donc redéfinir la responsabilité. Elle ne réside plus dans la conscience immédiate et transparente de l'acte (idéal cartésien intenable), mais dans la \cle{prise en charge} de son histoire psychique. Être responsable, c'est accepter de reconnaître comme siens les désirs que l'on découvre en soi, et travailler à les élaborer — c'est le sens de la cure psychanalytique et de l'exigence éthique formulée par Ricoeur : assumer la « dialectique de l'archéologie et de la téléologie », comprendre d'où l'on vient pour choisir qui l'on veut être. L'inconscient ne délivre pas de la responsabilité ; il en élargit le champ : je suis responsable non seulement de mes actes conscients, mais de ce que je fais de ma vérité inconsciente une fois qu'elle m'est révélée.
\end{corrige}

\begin{attention}[Barème indicatif — Exercice 8 (sur 20)]
Thèse correctement construite avec Freud (5 pts) ; antithèse juridique et morale avec le rôle du Moi (5 pts) ; synthèse personnelle redéfinissant la responsabilité (6 pts) ; clarté, orthographe, articulations logiques (4 pts).
\end{attention}

\begin{corrige}[Exercice 9 — Sujet de Dissertation : « La conscience suffit-elle à définir l'homme ? »]
\textbf{1. Introduction (corrigé rédigé).}

Depuis l'inscription du temple de Delphes, « Connais-toi toi-même », la philosophie occidentale a fait de la conscience le propre de l'homme : l'animal perçoit, l'homme se sait percevant. Mais que désigne la \emph{conscience} ? Au sens psychologique, elle est la présence immédiate de l'esprit à lui-même et à ses états ; au sens moral, le jugement du bien et du mal. « Définir » signifie donner l'essence, ce sans quoi une chose ne serait pas ce qu'elle est. La question est donc de savoir si la conscience épuise l'essentiel de l'humain, ou si l'homme ne se réduit pas à ce qu'il sait de lui-même. Descartes affirme la première option : « je pense, donc je suis » fait de la conscience le fondement certain de l'être humain. Mais Spinoza montre que la conscience ignore les causes qui nous déterminent, et Freud découvre un inconscient qui gouverne la vie psychique à l'insu du sujet. La conscience suffit-elle alors à définir l'homme, ou l'homme est-il aussi, et peut-être d'abord, ce qui en lui échappe à la conscience ? Nous verrons d'abord que la conscience constitue la dignité et la certitude du sujet (I), puis qu'elle est limitée et trompeuse, l'homme étant traversé par l'inconscient et la nécessité (II), enfin que l'homme se définit moins par la conscience possédée que par le rapport qu'il entretient à sa propre vérité (III).

\textbf{2. Plan détaillé.}
\begin{itemize}
\item \textbf{I. Thèse : la conscience comme essence de l'homme.}
\begin{itemize}
\item Arg. 1 : le \emph{cogito}, première certitude indubitable (Descartes, \emph{Méditations}) ; la pensée est l'attribut de la substance pensante.
\item Arg. 2 : la conscience morale fonde la responsabilité, la liberté et la dignité (on ne peut être tenu responsable que de ce dont on a conscience).
\item Exemple : le tribunal camerounais qui distingue l'acte volontaire de l'accident ; l'élève qui reconnaît sa faute.
\end{itemize}
\item \textbf{II. Antithèse : les limites de la conscience.}
\begin{itemize}
\item Arg. 1 : Spinoza — la conscience est connaissance inadéquate ; illusion du libre arbitre ; déterminisme de la nature.
\item Arg. 2 : Freud — l'inconscient, le refoulement, les formations de l'inconscient (rêves, lapsus, symptômes) ; « le Moi n'est pas maître en sa propre maison ».
\item Exemple camerounais : l'acte manqué du commerçant de Douala ; les rêves interprétés dans certaines traditions comme messages cachés — preuve que la culture elle-même admet que l'homme se dépasse.
\end{itemize}
\item \textbf{III. Synthèse : l'homme comme sujet divisé, défini par sa tâche de connaissance.}
\begin{itemize}
\item Arg. 1 : l'homme n'est ni pure conscience ni pur inconscient : il est le \emph{rapport} entre les deux (le Moi freudien comme médiateur).
\item Arg. 2 : se connaître est un travail infini (passage de la connaissance inadéquate à la connaissance adéquate chez Spinoza ; élaboration dans la cure chez Freud).
\item Conclusion provisoire : la conscience ne suffit pas à définir l'homme, mais l'effort pour s'éclairer soi-même, lui, est proprement humain.
\end{itemize}
\end{itemize}

\textbf{3. Conclusion (corrigé rédigé).}

La conscience ne suffit donc pas à définir l'homme : elle n'en est que la surface éclairée, traversée par des forces — pulsions, nécessité naturelle, histoire refoulée — qui la débordent. Mais ce qui définit l'homme n'est pas davantage l'inconscient seul : c'est la tension même entre ce qu'il sait de lui et ce qui en lui demeure obscur, et le travail incessant par lequel il cherche à s'approprier sa vérité. L'homme est l'être qui peut se méconnaître — et qui peut, précisément pour cela, entreprendre de se connaître. Reste une question ouverte : cette quête de soi a-t-elle un terme, ou l'homme est-il condamné, comme le suggère la durée infinie de toute cure, à n'être jamais qu'en chemin vers lui-même ?
\end{corrige}

\begin{attention}[Barème indicatif — Exercice 9 (sur 20)]
Introduction (accroche, définitions, problématique, plan) : 4 pts. Thèse argumentée avec Descartes : 4 pts. Antithèse avec Spinoza et Freud : 5 pts. Synthèse personnelle : 4 pts. Exemples (dont un camerounais), langue, présentation : 3 pts. \textbf{Points de vigilance :} définir « conscience » et « suffire » dès l'introduction ; citer les trois auteurs obligatoires ; éviter la paraphrase des cours sans articulation dialectique.
\end{attention}

\begin{corrige}[Exercice 10 — Explication de texte guidée : Freud, \emph{L'Interprétation des rêves}]
\begin{enumerate}
\item \textbf{Thèse principale.} Le rêve n'est pas un phénomène absurde mais la réalisation déguisée d'un désir refoulé, et son interprétation permet de remonter du récit apparent (manifeste) aux pensées inconscientes (latent) qui s'y expriment.
\item \textbf{Explication des concepts.}
\begin{itemize}
\item \cle{Réalisation déguisée} : le rêve accomplit imaginairement un désir, mais sous une forme transformée, méconnaissable, afin de ne pas éveiller le rêveur ni heurter sa censure morale.
\item \cle{Désir refoulé} : désir (souvent infantile, agressif ou sexuel) repoussé dans l'inconscient car incompatible avec le Moi et le Surmoi ; il demeure chargé d'énergie et cherche à se réaliser.
\item \cle{Censure} : instance psychique (entre l'inconscient et le préconscient) qui interdit l'accès des contenus refoulés à la conscience ; elle se relâche pendant le sommeil sans disparaître totalement — d'où la nécessité du déguisement.
\item \cle{Contenu latent / contenu manifeste} : le latent est l'ensemble des pensées et désirs inconscients du rêve ; le manifeste est le récit imagé tel que le rêveur s'en souvient et le raconte.
\item \cle{Déplacement} : transfert de l'intensité psychique d'une représentation importante vers une représentation secondaire ou indifférente, qui devient ainsi chargée de sens.
\item \cle{Condensation} : fusion de plusieurs éléments latents en une seule image manifeste, qui devient surdéterminée (elle renvoie à plusieurs chaînes d'associations).
\end{itemize}
\item \textbf{Structure du raisonnement.} (1) Définition : le rêve est réalisation déguisée d'un désir refoulé ; (2) justification du déguisement : la censure, affaiblie mais active durant le sommeil, impose une transformation ; (3) description des opérations de cette transformation : le travail du rêve (déplacement, condensation) convertit le latent en manifeste ; (4) conséquence méthodologique : l'interprétation doit parcourir le chemin inverse, du manifeste au latent, par les associations libres du rêveur. Le raisonnement va donc de la définition à la méthode, en passant par l'analyse des mécanismes.
\item \textbf{Discussion de la portée.} Cette conception porte un coup décisif à l'idéal cartésien : pour Descartes, la pensée est transparente à elle-même et la raison peut se rendre maîtresse de ses passions ; or Freud montre que même notre vie nocturne nous échappe, qu'un désir que nous ignorons écrit nos rêves, et que la conscience n'est qu'une partie de la psyché. Le sujet n'est plus maître en sa propre maison. Quant au statut scientifique de la psychanalyse : au sens strict du critère de Popper, elle n'est pas une science, car ses interprétations ne sont pas falsifiables — tout élément du rêve peut être interprété a posteriori comme confirmation de la théorie, et aucune observation ne pourrait la réfuter. Cela n'ôte pas à la psychanalyse sa valeur herméneutique et thérapeutique, mais interdit de la ranger parmi les sciences expérimentales.
\end{enumerate}
\end{corrige}

\begin{attention}[Points de vigilance — Exercice 10]
Dans une explication de texte, toujours partir du texte (citer les expressions) avant de mobiliser ses connaissances ; distinguer clairement explication (fidélité à l'auteur) et discussion (esprit critique argumenté).
\end{attention}

\begin{corrige}[Exercice 11 — Épreuve type Probatoire : « Peut-on se connaître soi-même ? »]
\textbf{1. Introduction complète (corrigé rédigé).}

Au fronton du temple de Delphes était gravée l'injonction : « Connais-toi toi-même ». Plus de deux millénaires après Socrate, qui en avait fait le programme de toute philosophie, cette exigence demeure aussi pressante qu'énigmatique. Car si rien ne semble plus proche de moi que moi-même — je suis présent à chacune de mes pensées, à chacun de mes actes —, rien peut-être ne m'est plus étranger : qui n'a jamais été surpris par ses propres réactions, déçu par ses propres choix, découvert en lui des sentiments qu'il ne se connaissait pas ?

Précisons d'abord les termes. \emph{Connaître} signifie saisir un objet par des représentations vraies et justifiées, par opposition à la simple opinion ou à la familiarité immédiate. Le \emph{soi} désigne ce qui fait l'identité singulière d'une personne : son caractère, ses désirs, ses motifs profonds, son histoire. La \emph{conscience} est la présence immédiate de l'esprit à lui-même et à ses états ; l'\emph{inconscient}, au sens freudien, désigne l'ensemble des contenus psychiques refoulés qui agissent en nous sans que nous en ayons connaissance. La question « peut-on se connaître soi-même ? » demande donc si le sujet peut devenir pour lui-même un objet de connaissance véritable, ou si une part de lui-même lui reste à jamais dérobée.

La difficulté tient à une tension fondamentale. D'un côté, la conscience paraît offrir un accès privilégié et immédiat à soi : Descartes en fait la première certitude, le \emph{cogito}, car je ne puis douter que je pense sans penser encore. De l'autre, cette transparence est mise en doute : Spinoza soutient que la conscience, connaissance inadéquate, nous donne nos désirs mais nous cache leurs causes ; Freud découvre un inconscient qui fait de nous des étrangers à nous-mêmes. La connaissance de soi est-elle alors un donné immédiat de la conscience, une impossibilité radicale, ou une tâche à accomplir ?

Nous examinerons d'abord la thèse de la conscience comme accès premier et certain à soi (I), puis l'antithèse de l'opacité du sujet, déterminé par des causes cachées et gouverné par l'inconscient (II), enfin la synthèse qui fait de la connaissance de soi une tâche infinie, éthique et narrative (III).

\textbf{2. Plan détaillé.}
\begin{itemize}
\item \textbf{I. Thèse : la conscience comme accès premier et certain à soi.}
\begin{itemize}
\item Arg. 1 : le \emph{cogito} cartésien — dans le doute, la pensée se saisit elle-même avec certitude ; l'âme est plus facile à connaître que le corps.
\item Arg. 2 : l'introspection — chacun a un accès direct et privilégié à ses états intérieurs (joie, douleur, doute) qu'aucun autrui ne peut contester.
\item Arg. 3 : la conscience morale — se connaître, c'est aussi pouvoir s'examiner, se juger, se repentir : condition de la responsabilité.
\item Exemple : l'élève qui, après une faute, reconnaît sincèrement sa colère et ses motifs.
\end{itemize}
\item \textbf{II. Antithèse : l'opacité radicale du sujet.}
\begin{itemize}
\item Arg. 1 : Spinoza — la conscience est connaissance inadéquate : nous sommes conscients de nos actions mais ignorants des causes qui nous déterminent ; illusion du libre arbitre.
\item Arg. 2 : Freud — l'inconscient, le refoulement ; les rêves, lapsus et symptômes révèlent des désirs que le sujet ignore et désavoue ; « le Moi n'est pas maître en sa propre maison ».
\item Arg. 3 : les biais de l'auto-observation — la mémoire reconstruit, l'amour-propre déforme, la mauvaise foi (que Sartre analysera) fait de nous nos propres trompeurs.
\item Exemple : l'acte manqué du commerçant de Douala qui « se trompe » de compte bancaire.
\end{itemize}
\item \textbf{III. Synthèse : la connaissance de soi comme tâche infinie, éthique et narrative.}
\begin{itemize}
\item Arg. 1 : ni transparence ni opacité totales — l'inconscient se manifeste (rêves, actes manqués) et peut donc être partiellement interprété : la cure psychanalytique est un travail de connaissance.
\item Arg. 2 : Spinoza lui-même ouvre la voie : passer de la connaissance inadéquate à la connaissance adéquate des causes, c'est transformer sa vie (éthique de la joie).
\item Arg. 3 : Ricoeur — le soi se connaît à travers le récit qu'il fait de lui-même (identité narrative) et à travers le regard d'autrui : se connaître est un détour, non une intuition.
\item Exemple : le dialogue avec un aîné, un griot ou un thérapeute comme médiation vers soi.
\end{itemize}
\end{itemize}

\textbf{3. Question bonus.} Le contexte camerounais éclaire singulièrement la question : la tradition orale (contes, proverbes, palabres) fait de la connaissance de soi une affaire \emph{communautaire} et transmise par la parole des anciens, non une introspection solitaire ; le rapport aux ancêtres inscrit l'identité dans une généalogie qui dépasse l'individu conscient ; et si la médecine traditionnelle interprète certains troubles comme des déséquilibres relationnels ou spirituels, la psychanalyse les lit comme des conflits intrapsychiques — deux herméneutiques du sujet qui, loin de s'exclure, montrent que se connaître soi-même suppose toujours une médiation culturelle.
\end{corrige}

\begin{attention}[Barème indicatif — Exercice 11 (sur 20)]
Introduction complète et problématisée : 5 pts. Partie I (Descartes, arguments, exemples) : 4 pts. Partie II (Spinoza, Freud, arguments) : 5 pts. Partie III (synthèse personnelle, Ricoeur) : 4 pts. Langue, méthode, présentation : 2 pts. \textbf{Points de vigilance :} ne pas transformer la dissertation en exposé sur Freud ; maintenir la question (« peut-on ? ») comme fil conducteur ; annoncer un plan dialectique numéroté.
\end{attention}

\newpage

\coursec{Fiche bilan}

\coursec{La Conscience et l'Inconscient}

\begin{popfigure}
\centering
\begin{tikzpicture}[
    mindmap,
    grow cyclic,
    every node/.style={concept, execute at begin node=\hskip0pt},
    concept color=popblue,
    level 1/.append style={level distance=4.5cm, sibling angle=360/6},
    level 2/.append style={level distance=3cm, sibling angle=360/4},
    text=white,
    font=\small\bfseries
]
\node[concept, scale=1.2] {Conscience \& Inconscient}
    child[concept color=poporange] { node[concept] {Nature de la conscience}
        child { node[concept, scale=0.7] {Cogito (Descartes)} }
        child { node[concept, scale=0.7] {Perception interne} }
        child { node[concept, scale=0.7] {Unité / Transparence} }
        child { node[concept, scale=0.7] {Intentionalité (Brentano/Husserl)} }
    }
    child[concept color=popgreen] { node[concept] {Valeur de la conscience}
        child { node[concept, scale=0.7] {Fondement moral (Kant)} }
        child { node[concept, scale=0.7] {Liberté \& Responsabilité} }
        child { node[concept, scale=0.7] {Connaissance de soi} }
        child { node[concept, scale=0.7] {Critique (Sartre)} }
    }
    child[concept color=poppurple] { node[concept] {Limites de la conscience}
        child { node[concept, scale=0.7] {Inconscient (Freud)} }
        child { node[concept, scale=0.7] {Conditionnement social} }
        child { node[concept, scale=0.7] {Erreur / Illusion} }
        child { node[concept, scale=0.7] {Corps \& Cérébralité} }
    }
    child[concept color=poppink] { node[concept] {L'inconscient freudien}
        child { node[concept, scale=0.7] {Topique : Cs / Pcs / Ics} }
        child { node[concept, scale=0.7] {Dynamic : Refoulement} }
        child { node[concept, scale=0.7] {Rêve : Voie royale} }
        child { node[concept, scale=0.7] {Pulsions : Vie / Mort} }
    }
    child[concept color=popgold] { node[concept] {Valeur de la psychanalyse}
        child { node[concept, scale=0.7] {Thérapie par la parole} }
        child { node[concept, scale=0.7] {Libération du sujet} }
        child { node[concept, scale=0.7] {Herméneutique (Ricœur)} }
        child { node[concept, scale=0.7] {Anthropologie philosophique} }
    }
    child[concept color=popdark] { node[concept] {Limites de la psychanalyse}
        child { node[concept, scale=0.7] {Non falsifiabilité (Popper)} }
        child { node[concept, scale=0.7] {Determinisme psychique} }
        child { node[concept, scale=0.7] {Durée / Coût / Transfert} }
        child { node[concept, scale=0.7] {Critique féministe / sociale} }
    };
\end{tikzpicture}
\legende{Carte mentale : Leçon 18 -- Conscience et Inconscient (Terminale Technique)}
\end{popfigure}

\begin{definition}[Définitions clés]
\begin{itemize}
    \item \cle{Conscience (psychologique)} : Connaissance immédiate, réflexive et unifiée que le sujet a de ses propres actes, pensées et sentiments (transparence cartésienne).
    \item \cle{Conscience morale} : Jugement pratique du bien et du mal ; voix de la raison pure pratique commandant l'action par le devoir (Kant : \emph{impératif catégorique}).
    \item \cle{Intentionalité} : Caractère de toute conscience d'être conscience \emph{de} quelque chose (Brentano, Husserl) ; il n'y a pas de conscience vide, elle vise toujours un objet.
    \item \cle{Inconscient (Freud)} : « Autre scène » psychique, réservoir des pulsions refoulées, régi par le principe de plaisir, intemporel, non contradictoire, inaccessible directement à la conscience.
    \item \cle{Refoulement} : Mécanisme de défense originaire repoussant dans l'Inconscient (Ics) des représentations ou affects incompatibles avec le Moi (Cs) ; fondement de l'inconscient dynamique.
    \item \cle{Psychanalyse} : Méthode d'investigation (association libre, interprétation des rêves, actes manqués) et thérapeutique visant à « rendre conscient l'inconscient » (\emph{Wo Es war, soll Ich werden} : « Là où c'était le Ça, doit advenir le Moi »).
\end{itemize}
\end{definition}

\begin{aretenir}[Thèses \& Auteurs de référence (à citer)]
\begin{center}
\begin{tabularx}{\linewidth}{|l|X|l|}
\hline
\rowcolor{popblueL}
\textbf{Notion} & \textbf{Thèse principale} & \textbf{Auteur} \\
\hline
Nature Cs & \emph{Cogito, ergo sum} : certitude indubitable, transparence, substantialité (res cogitans) & Descartes \\
\hline
Nature Cs & Conscience comme perception interne ; degrés (apperception) ; existence de « petites perceptions » inconscientes & Leibniz \\
\hline
Valeur Cs & Fondement de la dignité, de l'autonomie, du devoir ; \emph{Tu dois, donc tu peux} & Kant \\
\hline
Limites Cs & « Le Moi n'est pas maître en sa propre maison » ; déterminisme psychique, refoulement, pulsions & Freud \\
\hline
Inconscient & 1\textsuperscript{re} topique (Cs/Pcs/Ics) ; 2\textsuperscript{de} topique (Moi/Ça/Surmoi) ; langage de l'inconscient (rêve, acte manqué, symptôme) & Freud \\
\hline
Valeur Psychanalyse & Archéologie du sujet ; herméneutique du soupçon ; éthique du désir ; « rendre l'inconscient conscient » & Ricœur, Lacan \\
\hline
Limites Psychanalyse & Pseudoscience (non réfutable, Popper) ; réductionnisme biologisant ; rapport de pouvoir analyste/analysant ; critique du phallocentrisme & Popper, Grünbaum, Foucault, Irigaray \\
\hline
\end{tabularx}
\end{center}
\end{aretenir}

\begin{methode}[Méthodes express : Dissertation et Commentaire]
\begin{enumerate}
    \item \textbf{Analyser le sujet} : Distinguer impérativement \emph{conscience psychologique} (connaissance, transparence) et \emph{conscience morale} (valeur, obligation). Repérer si le sujet porte sur l'inconscient comme \emph{limite} de la conscience (négativité) ou comme \emph{instance positive} (créativité, désir, vérité du sujet).
    \item \textbf{Problématiser} : Exemple : « La conscience est-elle transparente à elle-même ? » $\rightarrow$ Tension entre l'évidence cartésienne (thèse) et l'opacité freudienne (antithèse). La synthèse pensera la conscience comme \emph{conquête} et non comme don originaire.
    \item \textbf{Structurer le plan} (type dialectique classique attendu au Bac) :
        \begin{itemize}
            \item \textbf{Thèse} : La conscience comme maîtrise, unité, fondement de la responsabilité et de la vérité (Descartes, Kant, idéalisme).
            \item \textbf{Antithèse} : La conscience comme limitée, trompeuse, déterminée par l'inconscient (Freud), l'idéologie (Marx), la vie (Nietzsche), le conditionnement social (Durkheim, Bourdieu).
            \item \textbf{Synthèse} : La conscience comme \emph{conquête} (Ricœur, Sartre, Merleau-Ponty) : lucidité critique, éthique de l'aveu, responsabilité assumée \emph{malgré} l'inconscient. Le sujet se construit \emph{par} l'épreuve de l'opacité.
        \end{itemize}
    \item \textbf{Illustrer par des exemples concrets camerounais} :
        \begin{itemize}
            \item \emph{Acte manqué} : Oubli répété d'un rendez-vous administratif crucial révélant un conflit d'autorité inconscient (rejet du pouvoir étatique ou paternaliste).
            \item \emph{Rêve codé} : Symboles culturels (totem, ancêtres, eau, serpent) interprétables comme réalisation déguisée d'un désir refoulé (ex: réussite aux examens, mariage).
            \item \emph{Symptôme social / Non-dit familial} : Gestion d'un héritage ou d'une succession bloquée par des non-dits, des secrets de famille (enfants cachés, sorcellerie) agissant comme un « refoulement social » générant des conflits ouverts.
        \end{itemize}
\end{enumerate}
\end{methode}

\begin{aretenir}[Checklist avant le Bac]
\begin{itemize}
    \item $\square$ Savoir définir \textbf{conscience} (psychologique vs morale) et \textbf{inconscient} (freudien : topique, dynamique, économique).
    \item $\square$ Expliquer le \textbf{Cogito} comme fondement de la certitude et de la transparence (Descartes) ; nuancer avec les « petites perceptions » de Leibniz.
    \item $\square$ Exposer la \textbf{conscience morale} comme autonomie et devoir (Kant : impératif catégorique, postulat de la liberté).
    \item $\square$ Démontrer les \textbf{limites de la conscience} : refoulement, pulsions, conditionnements sociaux (Freud, Marx, sociologie), corporéité (Merleau-Ponty).
    \item $\square$ Décrire la \textbf{méthode psychanalytique} (règle fondamentale : association libre ; interprétation ; cure type : transfert, résistance, levée du refoulement) et son but (\emph{Ich} à la place du \emph{Es}).
    \item $\square$ Citer \textbf{3 arguments pour la valeur} de la psychanalyse (efficacité thérapeutique, herméneutique du sens/Ricœur, anthropologie philosophique/Lacan).
    \item $\square$ Citer \textbf{3 arguments contre / limites} (critère de scientificité/Popper, déterminisme psychique niant la responsabilité, critiques féministes/sociales sur le pouvoir et la norme).
    \item $\square$ Maîtriser l'articulation \textbf{Conscience / Inconscient} : rupture épistémologique (Freud vs Descartes) ou continuité (Leibniz, Ricœur : l'inconscient comme « conscience blessée ») ?
    \item $\square$ Produire un \textbf{exemple concret camerounais} (acte manqué, rêve, symptôme social, non-dit familial) et l'analyser avec les concepts du cours.
    \item $\square$ Savoir rédiger une \textbf{introduction} (accroche, définitions des termes, problématique, annonce de plan) et une \textbf{conclusion} (réponse synthétique à la problématique, ouverture pertinente).
    \item $\square$ Respecter la \textbf{structure dialectique} (Thèse / Antithèse / Synthèse) avec transitions logiques (connecteurs : \emph{certes, cependant, en réalité, dès lors}).
\end{itemize}
\end{aretenir}

\findecours

\end{document}
